% Lectura y comentario : « Los medios de comunicación » ; textes argumentatifs / injonctifs ; cultura de la paz — Espagnol Tle A — Cours signé OnBuch+
% Programme officiel MINESEC. Compiler avec : tectonic E22-medios-de-comunicacion-cultura-de-paz.tex
\newcommand{\DOCMATIERE}{Espagnol}
\newcommand{\DOCNIVEAU}{Tle A}
\newcommand{\DOCMODULE}{Módulo 5 : Médias, communication et culture de la paix}
\newcommand{\DOCLECON}{E22}
\newcommand{\DOCTITRE}{Lectura y comentario : « Los medios de comunicación » ; textes argumentatifs / injonctifs ; cultura de la paz}
% OnBuch+ — préambule « cours signés » ENRICHI (compilé avec tectonic / XeLaTeX)
% Système visuel « Pop » d'OnBuch+ : fond crème (jamais blanc), bordures encre
% épaisses, ombres « dures » décalées, titres Archivo Black en capitales.
% Version MATHÉMATIQUES : graphes (pgfplots/tikz), boîtes pédagogiques
% (définition, propriété, méthode, attention, le savais-tu, exercice résolu,
% exercice, TP, bilan), page de garde et signature OnBuch+ ancrée en bas.
%
% Macros attendues dans main.tex AVANT \input{preamble} :
%   \DOCMATIERE  \DOCNIVEAU  \DOCMODULE  \DOCLECON (numéro)  \DOCTITRE
\documentclass[11pt]{article}

\usepackage{fontspec}
\usepackage[spanish,french]{babel} % français = langue principale ; l'espagnol via \esp{..} / espagnol
\usepackage[a4paper,margin=2cm,top=3.4cm,bottom=2.8cm,headheight=1.4cm,footskip=1.3cm]{geometry}
\usepackage{amsmath,amssymb,amsfonts}
\usepackage{mathtools} % \xmapsto, \coloneqq... souvent utilisés par les modèles
\usepackage{cancel} % simplifications barrées (bilans de forces)
% \abs{z} (module, valeur absolue) et \norm{u} (norme), utilisés par les modèles
\providecommand{\abs}{}\let\abs\relax\DeclarePairedDelimiter{\abs}{\lvert}{\rvert}
\providecommand{\norm}{}\let\norm\relax\DeclarePairedDelimiter{\norm}{\lVert}{\rVert}
\usepackage[table]{xcolor} % [table] : fournit aussi \rowcolors (alternance de lignes)
\usepackage{pagecolor}
\usepackage{tikz}
\usetikzlibrary{arrows.meta,positioning,calc,shapes.geometric,shapes.misc,shapes.arrows,decorations.pathmorphing,decorations.pathreplacing,decorations.markings,patterns,shadows,mindmap,trees,fit,backgrounds,matrix,intersections,angles,quotes}
\usepackage{pgfplots}
\usepackage[version=4]{mhchem} % équations nucléaires \ce{^{235}_{92}U}
\usepackage[european,siunitx]{circuitikz} % schémas électriques
\pgfplotsset{compat=1.18}
\usepgfplotslibrary{fillbetween,groupplots} % aires entre courbes (calcul intégral)
\usepackage{enumitem}
\usepackage{fancyhdr}
% [most] charge toutes les bibliothèques tcolorbox (listings, minted, raster,
% documentation...) et gonfle inutilement la mémoire TeX sur un document long
% (~300 boîtes) : on ne charge que ce qui est réellement utilisé.
\usepackage[skins,breakable]{tcolorbox}
\usepackage{graphicx}
\usepackage{colortbl}
\usepackage{booktabs}
\usepackage{tabularx}
\usepackage{diagbox} % case d'en-tête diagonale des tableaux à double entrée
% Colonnes extensibles centrée / gauche / droite (C, L, R), souvent utilisées par les modèles
\newcolumntype{C}{>{\centering\arraybackslash}X}
\newcolumntype{L}{>{\raggedright\arraybackslash}X}
\newcolumntype{R}{>{\raggedleft\arraybackslash}X}
\usepackage{array}
\usepackage{multicol}
\usepackage{siunitx}
% parse-numbers=false : accepte aussi \SI{2,5 \times 10^{-3}}{...}
\sisetup{locale=FR,output-decimal-marker={,},group-minimum-digits=4,parse-numbers=false}
\DeclareSIUnit\ohm{\ensuremath{\symup{\Omega}}} % U+2126 absent de la police
\DeclareSIUnit\division{div} % réglages d'oscilloscope (ms/div, V/div)
\DeclareSIUnit\tr{tr} % tours (vitesse de rotation, ex. \SI{1500}{\tr\per\minute})
\DeclareSIUnit\molar{M} % concentration molaire (ex. \SI{3}{\molar}, \SI{1}{\micro\molar})
\DeclareSIUnit\an{an} % année (le modèle confond parfois avec \year, un compteur TeX) — ex. \SI{5}{\kilo\meter\per\an}
\DeclareSIUnit\FCFA{FCFA} % franc CFA (ex. \SI{500}{\FCFA\per\kilo\gram})
\DeclareSIUnit\degreeBrix{\textdegree Brix} % teneur en sucre d'un sirop/jus (réfractométrie)
\DeclareSIUnit\month{mois} % le modèle utilise parfois \month, un compteur TeX, comme unité de durée
\usepackage{qrcode}
% \degree : commande fréquemment utilisée par les modèles pour un angle ou une
% température en dehors de \SI{}{} (non fournie par les paquets chargés).
\providecommand{\degree}{^{\circ}}
% \qed (fin de démonstration), utilisé par les modèles sans amsthm
\providecommand{\qed}{\hfill$\square$}
% \barre : double barre des valeurs interdites dans un tableau de variations (tkz-tab)
\providecommand{\barre}{\ensuremath{\|}}
% environnement proof (amsthm non chargé) utilisé par les modèles
\ifdefined\proof\else\newenvironment{proof}[1][Démonstration]{\par\noindent\textit{#1.}\ }{\qed\par}\fi
\usepackage{lastpage}
\usepackage{hyperref}
\hypersetup{hidelinks}

% ============================================================
% Palette « Pop » OnBuch+ — mêmes hex que la classe P (plus_ui.dart)
% ============================================================
\definecolor{poporange}{HTML}{F59321}
\definecolor{popdark}{HTML}{E07A0C}
\definecolor{popcream}{HTML}{FFF6E8}
\definecolor{popink}{HTML}{1C1714}
\definecolor{poppurple}{HTML}{7A5AE0}
\definecolor{popgreen}{HTML}{2E9E6A}
\definecolor{poppink}{HTML}{E0457B}
\definecolor{popblue}{HTML}{2F7DE1}
\definecolor{popgold}{HTML}{C98A12}
\definecolor{popmuted}{HTML}{8A7F76}
\definecolor{popline}{HTML}{EADFD2}
\definecolor{popgray}{HTML}{8A7F76}
\colorlet{popgrey}{popgray}
% Alias tolérants (le modèle nomme parfois les couleurs différemment)
\colorlet{popwhite}{white}
\colorlet{popblack}{popink}
\colorlet{popred}{poppink}
\colorlet{popyellow}{popgold}
% Teintes claires (fonds de tableaux/encadrés) — le modèle nomme parfois
% les couleurs avec un suffixe L (ex. popblueL) sans les avoir définies.
\colorlet{poporangeL}{poporange!20}
\colorlet{popdarkL}{popdark!20}
\colorlet{poppurpleL}{poppurple!20}
\colorlet{popgreenL}{popgreen!20}
\colorlet{poppinkL}{poppink!20}
\colorlet{popblueL}{popblue!20}
\colorlet{popgoldL}{popgold!20}
\colorlet{popredL}{popred!20}
\colorlet{popgrayL}{popgray!20}
% Teintes claires pour fonds de tableaux / zones
\colorlet{popblueL}{popblue!10!white}
\colorlet{poporangeL}{poporange!14!white}
\colorlet{popgreenL}{popgreen!12!white}
\colorlet{poppurpleL}{poppurple!12!white}
\colorlet{poppinkL}{poppink!12!white}
\colorlet{popgoldL}{popgold!14!white}

\pagecolor{popcream}

% ============================================================
% Typographie
% ============================================================
\defaultfontfeatures{Ligatures=TeX}
\setmainfont{PlusJakartaSans-Medium.ttf}[Path=../../../fonts/,BoldFont=PlusJakartaSans-Bold.ttf]
% Maths en sans-serif (Fira Math) avec chiffres et lettres droites de Plus
% Jakarta, pour que les formules se fondent dans le texte.
\usepackage[math-style=ISO,bold-style=ISO]{unicode-math}
\setmathfont{FiraMath-Regular.otf}[Path=../../../fonts/]
\setmathfont{PlusJakartaSans-Medium.ttf}[Path=../../../fonts/,range={up/num,up/{Latin,latin},\mathup}]
\usepackage{icomma} % virgule décimale sans espace en mode math
% Caractères mathématiques Unicode absents des polices de texte (Plus Jakarta,
% Archivo Black) : rendus par leur équivalent mathématique, en texte comme en
% formule — sinon ils s'affichent en carré vide (ex. « Arithmétique dans ℤ »).
\usepackage{newunicodechar}
\newunicodechar{ℝ}{\ensuremath{\mathbb{R}}}
\newunicodechar{ℤ}{\ensuremath{\mathbb{Z}}}
\newunicodechar{ℂ}{\ensuremath{\mathbb{C}}}
\newunicodechar{ℕ}{\ensuremath{\mathbb{N}}}
\newunicodechar{ℚ}{\ensuremath{\mathbb{Q}}}
\newunicodechar{𝔻}{\ensuremath{\mathbb{D}}}
\newunicodechar{α}{\ensuremath{\alpha}}
\newunicodechar{β}{\ensuremath{\beta}}
\newunicodechar{γ}{\ensuremath{\gamma}}
\newunicodechar{δ}{\ensuremath{\delta}}
\newunicodechar{ε}{\ensuremath{\varepsilon}}
\newunicodechar{θ}{\ensuremath{\theta}}
\newunicodechar{λ}{\ensuremath{\lambda}}
\newunicodechar{μ}{\ensuremath{\mu}}
\newunicodechar{σ}{\ensuremath{\sigma}}
\newunicodechar{τ}{\ensuremath{\tau}}
\newunicodechar{φ}{\ensuremath{\varphi}}
\newunicodechar{ψ}{\ensuremath{\psi}}
\newunicodechar{ω}{\ensuremath{\omega}}
\newunicodechar{Σ}{\ensuremath{\Sigma}}
% majuscules produites par \MakeUppercase dans les titres (α-aminés -> Α)
\newunicodechar{Α}{\ensuremath{\alpha}}
\newunicodechar{Β}{\ensuremath{\beta}}
\newunicodechar{Γ}{\ensuremath{\Gamma}}
\newunicodechar{Δ}{\ensuremath{\Delta}}
\newunicodechar{Ε}{\ensuremath{\varepsilon}}
\newunicodechar{Θ}{\ensuremath{\Theta}}
\newunicodechar{Λ}{\ensuremath{\Lambda}}
\newunicodechar{Μ}{\ensuremath{\mu}}
\newunicodechar{Τ}{\ensuremath{\tau}}
\newunicodechar{Φ}{\ensuremath{\Phi}}
\newunicodechar{Ψ}{\ensuremath{\Psi}}
\newunicodechar{Ω}{\ensuremath{\Omega}}
\newunicodechar{↦}{\ensuremath{\mapsto}}
\newunicodechar{∈}{\ensuremath{\in}}
\newunicodechar{∉}{\ensuremath{\notin}}
\newunicodechar{∧}{\ensuremath{\wedge}}
\newunicodechar{⇔}{\ensuremath{\Leftrightarrow}}
\newunicodechar{⟺}{\ensuremath{\Longleftrightarrow}}
\newunicodechar{⇒}{\ensuremath{\Rightarrow}}
\newunicodechar{⟹}{\ensuremath{\Longrightarrow}}
\newunicodechar{∘}{\ensuremath{\circ}}
\newunicodechar{∪}{\ensuremath{\cup}}
\newunicodechar{∩}{\ensuremath{\cap}}
\newunicodechar{≡}{\ensuremath{\equiv}}
\newunicodechar{⊂}{\ensuremath{\subset}}
\newunicodechar{⊥}{\ensuremath{\perp}}
\newunicodechar{∃}{\ensuremath{\exists}}
\newunicodechar{∀}{\ensuremath{\forall}}
\newunicodechar{∛}{\ensuremath{\sqrt[3]{\,}}}
\newunicodechar{∜}{\ensuremath{\sqrt[4]{\,}}}
\newunicodechar{⃗}{\ensuremath{{}^{\to}}}
\newunicodechar{ⁿ}{\ifmmode^{n}\else\textsuperscript{\ensuremath{n}}\fi}
\newunicodechar{ˣ}{\ifmmode^{x}\else\textsuperscript{\ensuremath{x}}\fi}
\newunicodechar{ᵇ}{\ifmmode^{b}\else\textsuperscript{\ensuremath{b}}\fi}
\newunicodechar{ʳ}{\ifmmode^{r}\else\textsuperscript{\ensuremath{r}}\fi}
\newunicodechar{ᵘ}{\ifmmode^{u}\else\textsuperscript{\ensuremath{u}}\fi}
\newunicodechar{ʸ}{\ifmmode^{y}\else\textsuperscript{\ensuremath{y}}\fi}
\newunicodechar{ᵃ}{\ifmmode^{a}\else\textsuperscript{\ensuremath{a}}\fi}
\newunicodechar{ᵉ}{\ifmmode^{e}\else\textsuperscript{\ensuremath{e}}\fi}
\newunicodechar{ᵏ}{\ifmmode^{k}\else\textsuperscript{\ensuremath{k}}\fi}
\newunicodechar{ᵖ}{\ifmmode^{p}\else\textsuperscript{\ensuremath{p}}\fi}
\newunicodechar{ᴺ}{\ifmmode^{N}\else\textsuperscript{\ensuremath{N}}\fi}
\newunicodechar{ᶜ}{\ifmmode^{c}\else\textsuperscript{\ensuremath{c}}\fi}
\newunicodechar{ʰ}{\ifmmode^{h}\else\textsuperscript{\ensuremath{h}}\fi}
\newunicodechar{ᵠ}{\ifmmode^{q}\else\textsuperscript{\ensuremath{q}}\fi}
\newunicodechar{ₙ}{\ifmmode_{n}\else\textsubscript{\ensuremath{n}}\fi}
\newunicodechar{ₐ}{\ifmmode_{a}\else\textsubscript{\ensuremath{a}}\fi}
\newunicodechar{ᵢ}{\ifmmode_{i}\else\textsubscript{\ensuremath{i}}\fi}
\newunicodechar{ᵣ}{\ifmmode_{r}\else\textsubscript{\ensuremath{r}}\fi}
\newunicodechar{ₖ}{\ifmmode_{k}\else\textsubscript{\ensuremath{k}}\fi}
\newunicodechar{ᵦ}{\ifmmode_{\beta}\else\textsubscript{\ensuremath{\beta}}\fi}
\newunicodechar{ₓ}{\ifmmode_{x}\else\textsubscript{\ensuremath{x}}\fi}
\newfontfamily\popdisplay{ArchivoBlack-Regular.ttf}[Path=../../../fonts/]
\newfontfamily\poplogo{SpaceGrotesk-Bold.ttf}[Path=../../../fonts/]
\newfontfamily\popsemi{PlusJakartaSans-SemiBold.ttf}[Path=../../../fonts/]
\newfontfamily\popxbold{PlusJakartaSans-ExtraBold.ttf}[Path=../../../fonts/]
\newfontfamily\popxboldls{PlusJakartaSans-ExtraBold.ttf}[Path=../../../fonts/,LetterSpace=3.0]

\setlist[enumerate]{leftmargin=1.7em,itemsep=5pt,topsep=4pt}
\setlist[itemize]{leftmargin=1.7em,itemsep=4pt,topsep=4pt,label={\color{poporange}\textbullet}}
\setlength{\parindent}{0pt}
\setlength{\parskip}{5pt}
\renewcommand{\arraystretch}{1.25}

% pgfplots : style Pop par défaut
\pgfplotsset{
  every axis/.append style={axis line style={popink,line width=1pt},tick style={popink},
    grid style={popline,line width=0.5pt},label style={font=\small},tick label style={font=\footnotesize}},
  popcurve/.style={poporange,line width=1.8pt,no markers,smooth},
}
% Même style disponible dans un \draw TikZ simple (hors pgfplots)
\tikzset{popcurve/.style={poporange,line width=1.8pt}}
\tikzset{popgrid/.style={popline,very thin}}
\colorlet{popcurve}{poporange} % le nom du style est parfois utilisé comme couleur

% ============================================================
% Pill / squiggle
% ============================================================
\newcommand{\poppill}[3]{%
  \tikz[baseline=-0.55ex]\node[draw=popink,line width=1.2pt,rounded corners=2.6pt,%
    fill=#2,inner xsep=6.5pt,inner ysep=3pt]{%
    \popxboldls\fontsize{7}{7}\selectfont\color{#3}\MakeUppercase{#1}};%
}
\newcommand{\popsquiggle}[1]{%
  \tikz[baseline=-0.4ex]{
    \draw[#1,line width=2.6pt,line cap=round]
      plot [smooth] coordinates {(0,0.09) (0.34,0.01) (0.68,0.21) (1.02,0.01) (1.36,0.10)};
  }%
}
% Mot-clé surligné (marqueur orange sous le texte)
\newcommand{\cle}[1]{{\popxbold\color{popdark}#1}}

% ============================================================
% En-tête / pied
% ============================================================
\pagestyle{fancy}
\fancyhf{}
\renewcommand{\headrulewidth}{0pt}
\renewcommand{\footrulewidth}{0pt}
\lhead{\raisebox{-0.15ex}{\poplogo\fontsize{13}{13}\selectfont\color{popink}OnBuch\color{poporange}+}}
\rhead{\poppill{\DOCMATIERE\ \textbullet\ \DOCNIVEAU\ \textbullet\ Leçon \DOCLECON}{white}{popink}}
\lfoot{\popxbold\fontsize{7.6}{7.6}\selectfont\color{popmuted}\MakeUppercase{Cours signé OnBuch+}\ \textbullet\ {\popsemi\DOCTITRE}}
\rfoot{\tikz[baseline=-0.6ex]\node[rounded corners=8.5pt,fill=popink,minimum height=17pt,minimum width=17pt,inner xsep=5pt,inner sep=0]{\popxbold\fontsize{8}{8}\selectfont\color{popcream}\thepage\,/\,\pageref*{LastPage}};}

% ============================================================
% Titres
% ============================================================
\newcommand{\chaptitre}[2]{%
  \begin{tcolorbox}[enhanced,colback=poporange,colframe=popink,boxrule=2.6pt,arc=11pt,
      drop shadow=popink,left=15pt,right=15pt,top=12pt,bottom=11pt,before skip=3pt,after skip=18pt]
    \poppill{#2}{popink}{popcream}\par\vspace{9pt}
    {\raggedright\hyphenpenalty=10000\exhyphenpenalty=10000\popdisplay\fontsize{22}{24}\selectfont\color{popcream}\MakeUppercase{#1}\par}
  \end{tcolorbox}
}
% Section numérotée : pastille orange avec le numéro + titre Archivo
\newcounter{popsec}
\newcommand{\coursec}[1]{%
  \stepcounter{popsec}\setcounter{popsub}{0}%
  \par\vspace{16pt}\noindent
  \begin{minipage}{\linewidth}
  \tikz[baseline=-0.7ex]\node[draw=popink,line width=1.6pt,fill=poporange,rounded corners=4pt,minimum size=22pt,inner sep=0]{\popdisplay\fontsize{12}{12}\selectfont\color{popcream}\thepopsec};%
  \hspace{8pt}{\popdisplay\fontsize{15}{17}\selectfont\color{popink}\MakeUppercase{#1}}\par
  \vspace{4pt}\noindent{\color{poporange}\rule{36pt}{3.6pt}}
  \end{minipage}\par\nopagebreak\vspace{8pt}
}
\newcounter{popsub}
\newcommand{\courssub}[1]{%
  \stepcounter{popsub}%
  \par\vspace{9pt}\noindent{\popxbold\fontsize{11.5}{13}\selectfont\color{popink}{\color{poporange}\thepopsec.\thepopsub}\hspace{6pt}#1}\par\nopagebreak\vspace{4pt}
}

% ============================================================
% Boîtes pédagogiques — même recette : fond blanc, bordure épaisse colorée,
% ombre dure encre, badge plein. Titre optionnel [#1].
% ============================================================
\tcbset{popbase/.style={breakable,enhanced,colback=white,boxrule=2.3pt,arc=9pt,
  shadow={3pt}{-3pt}{0pt}{popink},left=14pt,right=14pt,top=11pt,bottom=11pt,before skip=11pt,after skip=15pt,
  fontupper=\popsemi\fontsize{10.6}{14.5}\selectfont\color{popink},
  fontlower=\popsemi\fontsize{10.6}{14.5}\selectfont\color{popink}}}
% Coche dessinée (le glyphe ✓ n'existe pas dans les polices du document)
\newcommand{\popcheck}[1]{\tikz[baseline=-0.5ex]\draw[#1,line width=1.6pt,line cap=round,line join=round](-0.13,0.0)--(-0.04,-0.09)--(0.13,0.1);}
\providecommand{\coche}{\popcheck{popgreen}}
\providecommand{\resultat}[1]{\par\smallskip\noindent\fcolorbox{popgreen}{popgreenL}{\parbox{\dimexpr\linewidth-2\fboxsep-2\fboxrule\relax}{\textbf{Résultat.} #1}}\par\smallskip}
\newcommand{\popboxhead}[3]{\poppill{#1}{#2}{white}\ifx\relax#3\relax\else\hspace{6pt}{\popxbold\fontsize{10.6}{12}\selectfont\color{popink}#3}\fi\par\vspace{7pt}}

\newtcolorbox{aretenir}[1][]{popbase,colframe=popgreen,before upper={\popboxhead{À retenir}{popgreen}{#1}}}
% Texte en espagnol (document de lecture, dialogue, poème, extrait) : cadre
% d'étude. Un titre optionnel (source, auteur) s'affiche dans l'en-tête.
\newtcolorbox{textebox}[1][]{popbase,colframe=popblue,before upper={\popboxhead{Texto}{popblue}{#1}\begin{otherlanguage}{spanish}},after upper={\end{otherlanguage}}}
% Marques ✓ / ✗ (forme correcte / fautive) souvent utilisées par les modèles
\providecommand{\cross}{\textcolor{poppink}{\ensuremath{\times}}}
\providecommand{\xmark}{\cross}\providecommand{\crossmark}{\cross}\providecommand{\cmark}{\ensuremath{\checkmark}}
% Ligne de réponse / justification à compléter (inventée par les modèles)
\providecommand{\justification}[1]{\par\noindent\textit{Justification~:} \dotfill\par}
% \textipa (package tipa non chargé) : la police ne couvre pas l'API -> simple passage
\providecommand{\textipa}[1]{#1}
% Mot ou phrase en espagnol à l'intérieur d'un paragraphe français
\newcommand{\esp}[1]{\ifmmode\text{\textit{#1}}\else\begingroup\selectlanguage{spanish}\itshape #1\endgroup\fi}
\newtcolorbox{exemplebox}[1][]{popbase,colframe=poporange,before upper={\popboxhead{Exemple}{poporange}{#1}}}
\newtcolorbox{definition}[1][]{popbase,colframe=popblue,before upper={\popboxhead{Définition}{popblue}{#1}}}
\newtcolorbox{propriete}[1][]{popbase,colframe=poppurple,before upper={\popboxhead{Propriété}{poppurple}{#1}}}
\newtcolorbox{methode}[1][]{popbase,colframe=popgold,before upper={\popboxhead{Méthode}{popgold}{#1}}}
\newtcolorbox{attention}[1][]{popbase,colframe=poppink,before upper={\popboxhead{Attention}{poppink}{#1}}}
\newtcolorbox{experience}[1][]{popbase,colframe=popblue,colback=popblueL,before upper={\popboxhead{Expérience}{popblue}{#1}}}
% « Le savais-tu ? » : carte violette pleine (contexte, culture, Cameroun)
\newtcolorbox{savaistu}[1][]{popbase,colback=poppurple,colframe=popink,
  fontupper=\popsemi\fontsize{10.4}{14.2}\selectfont\color{white},
  before upper={\poppill{Le savais-tu\,?}{white}{poppurple}\ifx\relax#1\relax\else\hspace{6pt}{\popxbold\color{white}#1}\fi\par\vspace{7pt}}}
% Exercice résolu : énoncé en haut, \tcblower puis la solution sur fond crème
\newcounter{popexo}\newcounter{popexr}
\newtcolorbox{exoresolu}[1][]{popbase,colframe=poporange,colbacklower=popcream,
  segmentation style={popink,line width=1.2pt,dashed},
  before upper={\stepcounter{popexr}\popboxhead{Exercice résolu \thepopexr}{poporange}{#1}},
  before lower={\poppill{Solution}{popink}{popcream}\par\vspace{6pt}}}
% Exercice (non résolu en place) — la correction est dans la partie Corrigés
\newtcolorbox{exercice}[1][]{popbase,colframe=popink,
  before upper={\stepcounter{popexo}\popboxhead{Exercice \thepopexo}{popink}{#1}}}
% Corrigé
\newtcolorbox{corrige}[1][]{popbase,colframe=popgreen,colback=popgreenL,
  before upper={\popboxhead{Corrigé}{popgreen}{#1}}}
% Objectifs / prérequis / bilan
\newtcolorbox{objectifs}{popbase,colframe=popink,colback=poporangeL,
  before upper={\popboxhead{Objectifs de la leçon}{popink}{}}}
\newtcolorbox{prerequis}{popbase,colframe=popink,colback=popblueL,
  before upper={\popboxhead{Prérequis}{popblue}{}}}
\newtcolorbox{bilan}{popbase,colframe=popink,colback=white,boxrule=2.8pt,
  before upper={\popboxhead{Fiche bilan}{poporange}{L'essentiel à connaître pour l'examen}}}
% Figure encadrée avec légende
\newtcolorbox{popfigure}[1][]{enhanced,breakable=false,colback=white,colframe=popline,boxrule=1.4pt,arc=8pt,
  left=10pt,right=10pt,top=10pt,bottom=8pt,before skip=10pt,after skip=12pt,halign=center,
  lower separated=false,halign lower=center,
  fontlower=\popsemi\fontsize{9}{11.5}\selectfont\color{popmuted},
  #1}
\newcommand{\legende}[1]{\par\vspace{6pt}{\popsemi\fontsize{9}{11.5}\selectfont\color{popmuted}#1}}

% ============================================================
% Signature OnBuch+
% ============================================================
\newcommand{\poplogoinline}[1]{{\poplogo\fontsize{#1}{#1}\selectfont\color{popink}OnBuch\color{poporange}+}}
\newcommand{\signatureonbuch}{%
  \begin{tcolorbox}[enhanced,colback=popink,colframe=popink,boxrule=2.2pt,arc=10pt,
      drop shadow=poporange,left=14pt,right=14pt,top=10pt,bottom=10pt,before skip=0pt,after skip=0pt]
    \tikz[baseline=-0.7ex]\node[circle,fill=popgreen,draw=popcream,line width=1.3pt,minimum size=22pt,inner sep=0]{\popcheck{white}};%
    \hspace{9pt}\begin{minipage}[c]{0.62\linewidth}
      {\popxbold\fontsize{12}{13}\selectfont\color{popcream}Signé\ \textbullet\ {\poplogo OnBuch\color{poporange}+}}\par\vspace{2pt}
      {\raggedright\popsemi\fontsize{8}{10.5}\selectfont\color{popline}\MakeUppercase{Équipe pédagogique OnBuch+}\\\MakeUppercase{Conforme au programme officiel MINESEC}\par}
    \end{minipage}\hfill
    {\poplogo\fontsize{22}{22}\selectfont\color{popcream}OnBuch\color{poporange}+}
  \end{tcolorbox}%
}

% ============================================================
% Page de garde
% #1 titre · #2 matière · #3 niveau · #4 module · #5 n° leçon · #6 durée ·
% #7 description / objectifs (texte ou itemize)
% ============================================================
\newcommand{\pagegarde}[7]{%
  \thispagestyle{empty}
  \newgeometry{a4paper,margin=1.7cm,top=1.6cm,bottom=1.4cm}%
  \begin{tikzpicture}[remember picture,overlay]
    \fill[poporange!18] ([xshift=-3.2cm,yshift=-3.6cm]current page.north east) circle (4.6cm);
    \fill[poppurple!14] ([xshift=2.4cm,yshift=7.6cm]current page.south west) circle (3.8cm);
  \end{tikzpicture}%
  {\poplogo\fontsize{30}{30}\selectfont\color{popink}OnBuch\color{poporange}+}\hfill
  \raisebox{0.6ex}{\poppill{Cours signé \textbullet\ Premium}{popink}{popcream}}\par
  \vspace{1pt}\popsquiggle{poppurple}\par
  \vspace{8pt}
  \begin{tcolorbox}[enhanced,colback=poporange,colframe=popink,boxrule=2.8pt,arc=13pt,
      drop shadow=popink,left=18pt,right=18pt,top=12pt,bottom=13pt,after skip=10pt]
    \poppill{#2 \textbullet\ #3}{popink}{popcream}\hspace{5pt}\poppill{#4}{white}{popink}\par\vspace{8pt}
    {\popdisplay\fontsize{14}{14}\selectfont\color{popink}LEÇON #5}\par\vspace{4pt}
    {\raggedright\hyphenpenalty=10000\exhyphenpenalty=10000\popdisplay\fontsize{25}{27}\selectfont\color{popcream}\MakeUppercase{#1}\par}
    \vspace{8pt}
    {\popxbold\fontsize{9.5}{9.5}\selectfont\color{popink}\MakeUppercase{Durée conseillée : #6}}
  \end{tcolorbox}
  \begin{tcolorbox}[enhanced,colback=white,colframe=popink,boxrule=2.2pt,arc=10pt,
      drop shadow=popink,left=16pt,right=16pt,top=10pt,bottom=10pt,after skip=8pt]
    \poppill{Dans cette leçon}{poppurple}{white}\par\vspace{5pt}
    \popsemi\fontsize{9.3}{12.6}\selectfont\color{popink}#7
  \end{tcolorbox}
  \noindent\begin{minipage}[t]{0.487\linewidth}
    \begin{tcolorbox}[enhanced,colback=popgreen,colframe=popink,boxrule=2.2pt,arc=10pt,
        drop shadow=popink,left=11pt,right=11pt,top=10pt,bottom=10pt,height=3.1cm,valign=center]
      \tcbox[colback=white,colframe=popink,boxrule=1.3pt,arc=5pt,boxsep=0pt,nobeforeafter,
        left=3pt,right=3pt,top=3pt,bottom=3pt,valign=center]{\qrcode[height=1.9cm]{https://play.google.com/store/apps/details?id=com.luvvix.onbuch&hl=fr}}%
      \hspace{8pt}\begin{minipage}[c]{0.5\linewidth}
        {\popdisplay\fontsize{11}{12.5}\selectfont\color{white}\MakeUppercase{Télécharge OnBuch}}\par\vspace{4pt}
        {\raggedright\popsemi\fontsize{8.4}{11}\selectfont\color{white}Gratuit \textbullet\ Google Play\\Tuteur IA, annales, concours, résultats\par}
      \end{minipage}
    \end{tcolorbox}
  \end{minipage}\hfill
  \begin{minipage}[t]{0.487\linewidth}
    \begin{tcolorbox}[enhanced,colback=poppurple,colframe=popink,boxrule=2.2pt,arc=10pt,
        drop shadow=popink,left=11pt,right=11pt,top=10pt,bottom=10pt,height=3.1cm,valign=center]
      \tikz[baseline=-0.7ex]\node[circle,fill=white,draw=popink,line width=1.3pt,minimum size=1.6cm,inner sep=0]{\popdisplay\fontsize{24}{24}\selectfont\color{poppurple}+};%
      \hspace{8pt}\begin{minipage}[c]{0.6\linewidth}
        {\popdisplay\fontsize{11}{12.5}\selectfont\color{white}\MakeUppercase{Passe à OnBuch+}}\par\vspace{4pt}
        {\raggedright\popsemi\fontsize{8.4}{11}\selectfont\color{white}Tous les cours signés, Tuteur IA illimité, fiches PDF — onglet OnBuch+ de l'app\par}
      \end{minipage}
    \end{tcolorbox}
  \end{minipage}
  \vspace{8pt}
  \popcontenu
  \vfill
  \signatureonbuch
  \restoregeometry
  \newpage
}

% Rangée « Ce cours contient » (page de garde)
\newcommand{\poptile}[3]{% #1 couleur #2 gros texte #3 légende
  \begin{tcolorbox}[enhanced,colback=#1,colframe=popink,boxrule=2pt,arc=9pt,shadow={2.5pt}{-2.5pt}{0pt}{popink},
      left=6pt,right=6pt,top=9pt,bottom=9pt,halign=center,valign=center,height=1.75cm,width=\linewidth,nobeforeafter]
    {\popdisplay\fontsize{12.5}{13}\selectfont\color{white}\MakeUppercase{#2}}\par\vspace{3pt}
    {\centering\popsemi\fontsize{7.8}{9.5}\selectfont\color{white}#3\par}
  \end{tcolorbox}}
\newcommand{\popcontenu}{%
  \par\noindent{\popxboldls\fontsize{7.5}{7.5}\selectfont\color{popmuted}CE COURS SIGNÉ CONTIENT}\par\vspace{7pt}
  \noindent\begin{minipage}[t]{0.235\linewidth}\poptile{popblue}{Cours}{complet, illustré, pas à pas}\end{minipage}\hfill
  \begin{minipage}[t]{0.235\linewidth}\poptile{poporange}{Méthodes}{et exercices résolus}\end{minipage}\hfill
  \begin{minipage}[t]{0.235\linewidth}\poptile{poppink}{Exercices}{type examen + corrigés}\end{minipage}\hfill
  \begin{minipage}[t]{0.235\linewidth}\poptile{popgreen}{Fiche bilan}{l'essentiel à retenir}\end{minipage}\par}

% Sommaire
\newtcolorbox{sommaire}{enhanced,colback=white,colframe=popink,boxrule=2.2pt,arc=10pt,
  drop shadow=popink,left=18pt,right=18pt,top=13pt,bottom=13pt,before skip=6pt,after skip=20pt,
  before upper={\poppill{Sommaire}{poppurple}{white}\par\vspace{9pt}\popsemi\fontsize{10.6}{14.5}\selectfont\color{popink}}}

% Signature de fin de cours — poussée tout en bas de la dernière page
\newcommand{\findecours}{%
  \par\vfill\nopagebreak
  \begin{center}\popsquiggle{poporange}\end{center}\vspace{4pt}
  \signatureonbuch
}
\AtBeginDocument{\def\llbracket{\mathopen{[\mkern-3.5mu[}}\def\rrbracket{\mathclose{]\mkern-3.5mu]}}} % intervalles d'entiers (glyphe ⟦ absent de la police)
% Glyphes absents de Fira Math : redéfinis à partir de glyphes disponibles
\AtBeginDocument{%
  \def\setminus{\mathbin{\backslash}}\let\smallsetminus\setminus
  \def\vdots{\mathord{\vbox{\baselineskip3.2pt\lineskiplimit0pt\kern4pt\hbox{.}\hbox{.}\hbox{.}}}}%
  \def\ddots{\mathinner{\mkern1mu\raise6.5pt\hbox{.}\mkern2mu\raise3.6pt\hbox{.}\mkern2mu\raise0.7pt\hbox{.}\mkern1mu}}%
  \def\triangle{\mathord{\scalebox{0.9}{$\symup{\Delta}$}}}%
  \def\nabla{\mathord{\raisebox{\depth}{\scalebox{1}[-1]{$\symup{\Delta}$}}}}%
  \def\checkmark{\popcheck{popdark}}%
  \def\star{\mathbin{\ast}}\def\bigstar{\mathord{\ast}}%
  \def\diamond{\mathbin{\scalebox{0.6}{$\Diamond$}}}%
  \def\bigcup{\mathop{\mathchoice{\raisebox{-0.3em}{\scalebox{1.7}{$\cup$}}}{\scalebox{1.25}{$\cup$}}{\cup}{\cup}}\displaylimits}%
  \def\bigcap{\mathop{\mathchoice{\raisebox{-0.3em}{\scalebox{1.7}{$\cap$}}}{\scalebox{1.25}{$\cap$}}{\cap}{\cap}}\displaylimits}%
  \def\lesseqqgtr{\mathrel{\lessgtr}}\def\bigoplus{\mathop{\oplus}\displaylimits}%
}
\providecommand{\ssi}{si et seulement si} % abréviation fréquente
\AtBeginDocument{\providecommand{\wideparen}{\overparen}} % arc de cercle

\begin{document}

\pagegarde{Lectura y comentario : « Los medios de comunicación » ; textes argumentatifs / injonctifs ; cultura de la paz}{Espagnol}{Tle A}{Módulo 5 : Médias, communication et culture de la paix}{E22}{2 h}{Cette leçon de lecture-commentaire propose l'étude d'un texte argumentatif original sur les médias et la culture de la paix. Elle permet d'acquérir le lexique des médias et de la paix, la méthode du commentaire de texte et la formulation d'injonctions et d'appels à la paix, compétences clés pour l'épreuve du Bac.
\vspace{4pt}
\begin{multicols}{2}\small
\begin{itemize}
\item À la fin de cette leçon, je sais lire et comprendre globalement et finement un texte argumentatif sur les médias.
\item Je sais identifier la thèse, les arguments et les exemples dans un texte argumentatif.
\item Je sais reconnaître les marques du texte injonctif (impératif, subjonctif d'exhortation, futur d'ordre).
\item Je sais utiliser le lexique des médias : prensa, televisión, radio, periodista, noticia, informativo, redes sociales.
\item Je sais utiliser le lexique de la paix : paz, conflicto, diálogo, violencia, tolerancia, reconciliación.
\item Je sais rédiger un commentaire structuré (introduction, développement, conclusion) sur un texte argumentatif.
\end{itemize}\end{multicols}}

\begin{sommaire}
\begin{multicols}{2}
\begin{enumerate}
\item Texte support : « Los medios de comunicación, ¿armas de guerra o instrumentos de paz? »
\item Compréhension du texte
\item Lexique thématique : médias et culture de la paix
\item Textes argumentatifs et injonctifs : repérage et grammaire de l'appel
\item Méthode du commentaire de texte et commentaire modèle
\item Production guidée : écrire un appel à la paix
\item Activité d'intégration
\item Méthodes et exercices résolus
\item Exercices
\item Corrigés des exercices
\item Fiche bilan
\end{enumerate}
\end{multicols}
\end{sommaire}

\begin{objectifs}
\begin{itemize}
\item À la fin de cette leçon, je sais lire et comprendre globalement et finement un texte argumentatif sur les médias.
\item Je sais identifier la thèse, les arguments et les exemples dans un texte argumentatif.
\item Je sais reconnaître les marques du texte injonctif (impératif, subjonctif d'exhortation, futur d'ordre).
\item Je sais utiliser le lexique des médias : prensa, televisión, radio, periodista, noticia, informativo, redes sociales.
\item Je sais utiliser le lexique de la paix : paz, conflicto, diálogo, violencia, tolerancia, reconciliación.
\item Je sais rédiger un commentaire structuré (introduction, développement, conclusion) sur un texte argumentatif.
\item Je sais formuler des injonctions et des appels à la paix à l'impératif et avec « ojalá » / « que + subjonctif ».
\item Je sais produire un court texte argumentatif ou injonctif (affiche, tribune) en faveur de la culture de la paix.
\end{itemize}
\end{objectifs}

\begin{prerequis}
\begin{itemize}
\item L'impératif affirmatif et négatif (leçon de Première).
\item Le présent de l'indicatif et du subjonctif des verbes réguliers.
\item La lecture méthodique d'un texte journalistique (Première, module sur la presse).
\item Les connecteurs logiques de base : porque, sin embargo, además, por eso.
\item Le vocabulaire général de la société et de l'actualité.
\end{itemize}
\end{prerequis}

\newpage

\coursec{Texte support : « Los medios de comunicación, ¿armas de guerra o instrumentos de paz? »}

Chaque matin à Yaoundé, des millions de Camerounais allument la radio, consultent leur téléphone ou regardent le journal télévisé : les médias façonnent notre vision du monde. Mais que se passe-t-il quand une fausse information circule sur WhatsApp et attise un conflit entre communautés ? Dans le monde hispanophone, des campagnes comme \esp{No más violencia} montrent que les médias peuvent aussi construire la paix. Comment lire un texte qui défend une thèse sur le rôle des médias ? Comment repérer ses arguments et formuler, à notre tour, un appel au dialogue ? C'est tout l'enjeu de cette leçon, conforme au Module 5 du programme officiel de Terminale A (Médias, communication et culture de la paix).

\courssub{Présentation du texte}

Le texte que nous allons étudier est un \cle{article d'opinion} (\esp{artículo de opinión}) publié dans une revue scolaire hispanophone fictive, \esp{La Voz del Estudiante}. Ce type de texte est caractéristique de l'unité 5 du manuel Didáctica V, consacrée aux médias, à la communication et à la culture de la paix. Il s'inscrit pleinement dans le module 5 du programme de Terminale A : « Médias, communication et culture de la paix ».

\begin{definition}[Qu'est-ce qu'un texte argumentatif ?]
Un \cle{texte argumentatif} (\esp{texto argumentativo}) est un texte qui défend une \cle{thèse} (une opinion, \esp{tesis}) à l'aide d'\cle{arguments} (des raisons, \esp{argumentos}) et d'\cle{exemples} (des illustrations, \esp{ejemplos}) pour convaincre le lecteur. Sa structure typique est : introduction (présentation du sujet et de la thèse), développement (arguments et exemples), conclusion (rappel de la thèse, ouverture).
\end{definition}

\begin{definition}[Qu'est-ce qu'un texte injonctif ?]
Un \cle{texte injonctif} (\esp{texto injuntivo}) est un texte qui cherche à \emph{faire agir} le lecteur : ordres, conseils, appels, slogans. Ses marques typiques en espagnol sont : l'\cle{impératif} (\esp{¡Escucha!}, \esp{¡Dialoguemos!}), la tournure \esp{hay que + infinitivo} (« il faut »), et \esp{que + subjuntivo} (\esp{¡Que viva la paz!}). Les appels à la paix relèvent de ce genre.
\end{definition}

\courssub{Lecture du texte}

\textbf{Première étape : lecture silencieuse.}\par
Lisez le texte en silence, sans dictionnaire, en vous laissant guider par les mots transparents. Votre objectif : saisir le \emph{sens général} et identifier la thèse défendue par l'auteur.

\textbf{Deuxième étape : lecture à voix haute par l'enseignant.}\par
Écoutez attentivement la prononciation : notez l'intonation montante de la question du titre (\esp{¿armas de guerra o instrumentos de paz?}) et le ton exclamatif de l'appel final (\esp{¡usemos los medios para unir, no para dividir!}). Quelle est votre première impression ? Le texte est-il pessimiste, optimiste, engagé ?

\begin{textebox}[Los medios de comunicación, ¿armas de guerra o instrumentos de paz? — La Voz del Estudiante]
En nuestro mundo conectado, una noticia falsa puede encender un conflicto en pocas horas. Basta un mensaje anónimo en una red social para que dos barrios, dos comunidades o dos países se miren con desconfianza. Sin embargo, la radio, la prensa y la televisión también pueden construir la paz: informan con verdad, dan voz a las víctimas y promueven el diálogo entre las comunidades.

El periodista responsable verifica sus fuentes antes de publicar. Sabe que cada titular puede calmar los ánimos o, al contrario, aumentar la violencia. Por eso, su profesión es también una misión de paz. Cuando un informativo explica los hechos con objetividad, los ciudadanos comprenden mejor a sus vecinos y el miedo disminuye.

Además, las redes sociales, usadas con inteligencia, difunden mensajes de tolerancia y reconciliación. En muchos países hispanos, campañas como «No más violencia» han unido a miles de jóvenes alrededor de valores comunes: el respeto, la justicia y la convivencia. En Colombia, por ejemplo, las radios comunitarias han ayudado a la reconciliación después de años de conflicto armado.

Los medios, por lo tanto, no son buenos ni malos en sí mismos: todo depende del uso que hacemos de ellos. Un micrófono puede gritar odio, pero también puede cantar la fraternidad. Una pantalla puede mostrar la guerra, pero también puede enseñar la cultura del otro.

Por eso decimos: ¡usemos los medios para unir, no para dividir! ¡Verifiquemos antes de compartir! ¡Escuchemos todas las voces! La paz empieza con una información verdadera, y cada uno de nosotros es responsable de lo que publica, de lo que comenta y de lo que comparte. El futuro de nuestras sociedades se escribe también en nuestras pantallas.
\end{textebox}

\textbf{Traduction du texte (Version).}\par
« Dans notre monde connecté, une fausse nouvelle peut déclencher un conflit en quelques heures. Il suffit d'un message anonyme sur un réseau social pour que deux quartiers, deux communautés ou deux pays se regardent avec méfiance. Cependant, la radio, la presse et la télévision peuvent aussi construire la paix : elles informent avec vérité, donnent la parole aux victimes et promeuvent le dialogue entre les communautés.

Le journaliste responsable vérifie ses sources avant de publier. Il sait que chaque gros titre peut calmer les esprits ou, au contraire, augmenter la violence. C'est pourquoi sa profession est aussi une mission de paix. Quand un journal télévisé explique les faits avec objectivité, les citoyens comprennent mieux leurs voisins et la peur diminue.

De plus, les réseaux sociaux, utilisés avec intelligence, diffusent des messages de tolérance et de réconciliation. Dans de nombreux pays hispaniques, des campagnes comme "Plus jamais de violence" ont uni des milliers de jeunes autour de valeurs communes : le respect, la justice et la vie commune. En Colombie, par exemple, les radios communautaires ont aidé à la réconciliation après des années de conflit armé.

Les médias, par conséquent, ne sont ni bons ni mauvais en eux-mêmes : tout dépend de l'usage que nous en faisons. Un micro peut crier la haine, mais il peut aussi chanter la fraternité. Un écran peut montrer la guerre, mais il peut aussi enseigner la culture de l'autre.

C'est pourquoi nous disons : utilisons les médias pour unir, pas pour diviser ! Vérifions avant de partager ! Écoutons toutes les voix ! La paix commence avec une information véridique, et chacun de nous est responsable de ce qu'il publie, de ce qu'il commente et de ce qu'il partage. L'avenir de nos sociétés s'écrit aussi sur nos écrans. »

\courssub{Glossaire et mots transparents}

\begin{definition}[Glossaire du texte : médias et paix]
\begin{itemize}
\item \esp{la prensa} : la presse (écrite)
\item \esp{el informativo} : le journal télévisé ou radiodiffusé
\item \esp{la noticia} : la nouvelle, l'information (une seule)
\item \esp{el titular} : le titre, le gros titre (d'un article)
\item \esp{la desinformación} : la désinformation
\item \esp{el diálogo} : le dialogue
\item \esp{la reconciliación} : la réconciliation
\item \esp{la convivencia} : la vie commune, le vivre-ensemble
\item \esp{la fuente} : la source (d'information)
\item \esp{la red social} : le réseau social
\item \esp{la paz} : la paix
\item \esp{el conflicto} : le conflit
\item \esp{la tolerancia} : la tolérance
\item \esp{la objetividad} : l'objectivité
\item \esp{la fraternidad} : la fraternité
\item \esp{verificar} : vérifier
\item \esp{difundir} : diffuser, répandre
\item \esp{encender} : allumer, déclencher (un conflit)
\end{itemize}
\end{definition}

\begin{aretenir}[Les mots transparents : un trésor pour le lecteur francophone]
Les \cle{mots transparents} sont des mots espagnols très proches du français, faciles à reconnaître. Dans notre texte : \esp{comunicación} (communication), \esp{televisión} (télévision), \esp{conflicto} (conflit), \esp{violencia} (violence), \esp{tolerancia} (tolérance), \esp{objetividad} (objectivité), \esp{fraternidad} (fraternité), \esp{responsable} (responsable), \esp{información} (information), \esp{sociedades} (sociétés). Stratégie de lecture : repérez d'abord ces mots pour deviner le thème général avant même de lire en détail.
\end{aretenir}

\begin{attention}[Faux amis et pièges de traduction]
\begin{itemize}
\item \esp{noticia} signifie « nouvelle, information » (une seule), pas « notice » (mode d'emploi).
\item \esp{actualidad} signifie « actualité », pas « actuality » (réalité).
\item \esp{compartir} signifie « partager » (sur les réseaux sociaux) : \esp{compartir un mensaje} « partager un message ». Ne pas confondre avec \esp{repartir} (répartir).
\item Attention au genre : \esp{la red} (le réseau) est féminin, mais \esp{el internet} et \esp{el titular} sont masculins.
\item \esp{para que + subjuntivo} : \esp{para que dos barrios se miren} $\rightarrow$ « pour que deux quartiers se regardent » (but). Ne pas traduire l'indicatif \emph{se miran} (ils se regardent) mais le subjonctif \emph{se miren}.
\item \esp{en sí mismos} : « en eux-mêmes » (réfléchi), pas « en si mêmes ».
\end{itemize}
\end{attention}

\courssub{Carte mentale et schéma du raisonnement}

\begin{popfigure}
\begin{tikzpicture}[mindmap, grow cyclic, every node/.style={concept, align=center, font=\small},
root concept/.append style={concept color=popblue, font=\small\bfseries},
level 1/.append style={level distance=3.2cm, sibling angle=90},
level 2/.append style={level distance=2.4cm, sibling angle=60}]
\node[root concept]{Los medios de comunicación}
  child[concept color=poporange] { node {prensa}
    child[concept color=popgreen] { node[align=center]{paz:\\ titular\\ veraz} }
    child[concept color=poppink] { node[align=center]{guerra:\\ desinformación} } }
  child[concept color=poporange] { node {televisión}
    child[concept color=popgreen] { node[align=center]{paz:\\ informativo\\ objetivo} }
    child[concept color=poppink] { node[align=center]{guerra:\\ imágenes\\ manipuladas} } }
  child[concept color=poporange] { node {radio}
    child[concept color=popgreen] { node[align=center]{paz:\\ radios\\ comunitarias} }
    child[concept color=poppink] { node[align=center]{guerra:\\ discursos\\ de odio} } }
  child[concept color=poporange] { node {redes sociales}
    child[concept color=popgreen] { node[align=center]{paz:\\ campañas\\ de tolerancia} }
    child[concept color=poppink] { node[align=center]{guerra:\\ noticias\\ falsas} } };
\end{tikzpicture}
\legende{Carte mentale : chaque média peut servir la guerre ou la paix selon l'usage qu'on en fait.}
\end{popfigure}

\begin{popfigure}
\begin{tikzpicture}[node distance=0.6cm, every node/.style={align=center, font=\small}]
\node[draw=popgreen, fill=popgreenL, rounded corners, thick, minimum width=3cm, minimum height=0.9cm, align=center] (a1){información\\ veraz};
\node[draw=popgreen, fill=popgreenL, rounded corners, thick, minimum width=3cm, minimum height=0.9cm, right=1.2cm of a1] (a2) {diálogo};
\node[draw=popgreen, fill=popgreenL, rounded corners, thick, minimum width=3cm, minimum height=0.9cm, right=1.2cm of a2] (a3) {paz};
\draw[-{Stealth}, very thick, popgreen] (a1) -- (a2);
\draw[-{Stealth}, very thick, popgreen] (a2) -- (a3);
\node[draw=poppink, fill=poppinkL, rounded corners, thick, minimum width=3cm, minimum height=0.9cm, below=1.4cm of a1] (b1) {desinformación};
\node[draw=poppink, fill=poppinkL, rounded corners, thick, minimum width=3cm, minimum height=0.9cm, right=1.2cm of b1] (b2) {conflicto};
\node[draw=poppink, fill=poppinkL, rounded corners, thick, minimum width=3cm, minimum height=0.9cm, right=1.2cm of b2] (b3) {violencia};
\draw[-{Stealth}, very thick, poppink] (b1) -- (b2);
\draw[-{Stealth}, very thick, poppink] (b2) -- (b3);
\end{tikzpicture}
\legende{Les deux chemins des médias : la chaîne de la paix contre la chaîne de la violence.}
\end{popfigure}

\courssub{Premier repérage : thèse, arguments, exemples}

Observons comment l'auteur organise sa démonstration. La \cle{thèse} est annoncée dès le titre sous forme de question, puis affirmée au quatrième paragraphe : \esp{Los medios no son buenos ni malos en sí mismos: todo depende del uso que hacemos de ellos.} (« Les médias ne sont ni bons ni mauvais en eux-mêmes : tout dépend de l'usage que nous en faisons. »)

\begin{center}
\begin{tabularx}{\linewidth}{|X|X|X|}\hline
\rowcolor{popblueL} Élément & Passage du texte & Traduction \\ \hline
Thèse & \esp{todo depende del uso que hacemos de ellos} & tout dépend de l'usage que nous en faisons \\ \hline
Argument 1 & \esp{una noticia falsa puede encender un conflicto} & une fausse nouvelle peut déclencher un conflit \\ \hline
Argument 2 & \esp{el periodista responsable verifica sus fuentes} & le journaliste responsable vérifie ses sources \\ \hline
Argument 3 & \esp{las redes sociales difunden mensajes de tolerancia} & les réseaux sociaux diffusent des messages de tolérance \\ \hline
Exemple & \esp{las radios comunitarias en Colombia} & les radios communautaires en Colombie \\ \hline
Appel (injonctif) & \esp{¡usemos los medios para unir, no para dividir!} & utilisons les médias pour unir, pas pour diviser ! \\ \hline
\end{tabularx}
\end{center}

\begin{propriete}[Les connecteurs logiques du texte argumentatif]
Pour enchaîner thèse et arguments, l'espagnol utilise des \cle{connecteurs} qu'il faut savoir repérer :
\begin{itemize}
\item opposition : \esp{sin embargo} (cependant), \esp{al contrario} (au contraire), \esp{pero} (mais) ;
\item addition : \esp{además} (de plus), \esp{también} (aussi) ;
\item cause et conséquence : \esp{por eso} (c'est pourquoi), \esp{por lo tanto} (par conséquent), \esp{porque} (parce que) ;
\item exemple : \esp{por ejemplo} (par exemple).
\end{itemize}
Ces mots sont les « panneaux indicateurs » du raisonnement : les souligner, c'est déjà comprendre la structure du texte.
\end{propriete}

\begin{exoresolu}[Repérer la thèse et les arguments]
Énoncé. Lisez cet extrait : \esp{La radio es el medio más cercano a la gente. Sin embargo, a veces difunde rumores. Por eso, los periodistas deben verificar la información antes de emitirla.} 1) Quelle est la thèse implicite ? 2) Quel connecteur exprime l'opposition ? 3) Quel connecteur exprime la conséquence ?
\tcblower
Solution détaillée. 1) La thèse implicite est que la radio est un média utile mais qui exige de la responsabilité : l'auteur valorise la radio (\esp{el medio más cercano a la gente}, « le média le plus proche des gens ») tout en signalant son danger. 2) Le connecteur d'opposition est \esp{sin embargo} (« cependant ») : il introduit la contrepartie négative (\esp{a veces difunde rumores}, « elle diffuse parfois des rumeurs »). 3) Le connecteur de conséquence est \esp{por eso} (« c'est pourquoi ») : il introduit la solution préconisée, formulée avec \esp{deben + infinitivo} (\esp{deben verificar}, « ils doivent vérifier »), marque typique du texte injonctif.
\end{exoresolu}

\begin{savaistu}[La radio, média de la paix en Afrique centrale et en Amérique latine]
La radio reste le média le plus écouté en Afrique centrale : au Cameroun, elle informe les villages les plus éloignés, souvent dans les langues locales (bassa, bamoun, fulfulde...). En Amérique latine, des \esp{radios comunitarias} (« radios communautaires ») jouent un rôle majeur dans les processus de paix : en Colombie, après les accords de paix de 2016 entre le gouvernement et les FARC, ces radios ont donné la parole aux anciens combattants et aux victimes pour favoriser la \esp{reconciliación}. La Guinée équatoriale, seul pays hispanophone d'Afrique, partage avec le Cameroun cette tradition de la radio comme lien social et outil d'éducation populaire.
\end{savaistu}

\begin{exercice}[Questions de compréhension]
Répondez en espagnol par des phrases complètes.
\begin{enumerate}
\item ¿Qué puede encender una noticia falsa en pocas horas?
\item ¿Qué hace el periodista responsable antes de publicar?
\item ¿Qué difunden las redes sociales usadas con inteligencia?
\item ¿Qué ejemplo de país hispano menciona el texto?
\item ¿Cuál es la conclusión del autor sobre los medios?
\end{enumerate}
\end{exercice}

\begin{corrige}[Exercice : Questions de compréhension]
\begin{enumerate}
\item \esp{Una noticia falsa puede encender un conflicto en pocas horas.}
\item \esp{El periodista responsable verifica sus fuentes antes de publicar.}
\item \esp{Difunden mensajes de tolerancia y reconciliación.}
\item \esp{Menciona el ejemplo de Colombia y sus radios comunitarias.}
\item \esp{El autor concluye que los medios no son buenos ni malos: todo depende del uso que hacemos de ellos, y que la paz empieza con una información verdadera.}
\end{enumerate}
\end{corrige}

\courssub{Lexique thématique approfondi : médias et culture de la paix}

Pour réussir l'expression écrite et orale sur ce thème, il faut maîtriser les champs lexicaux suivants. Apprenez-les par catégories.

\begin{center}
\begin{tabularx}{\linewidth}{|>{\bfseries}X|X|X|}\hline
\rowcolor{popblueL} \textbf{Médias \& Information} & \textbf{Verbes d'action} & \textbf{Paix \& Citoyenneté} \\ \hline
\esp{el medio de comunicación} (média) & \esp{informar} (informer) & \esp{la paz} (paix) \\
\esp{la prensa} (presse) & \esp{comunicar} (communiquer) & \esp{el diálogo} (dialogue) \\
\esp{la radio} (radio) & \esp{verificar} (vérifier) & \esp{la reconciliación} (réconciliation) \\
\esp{la televisión} (télévision) & \esp{contrastar} (recouper/croiser les sources) & \esp{la convivencia} (vivre-ensemble) \\
\esp{las redes sociales} (réseaux sociaux) & \esp{difundir} (diffuser) & \esp{la tolerancia} (tolérance) \\
\esp{el periodista} (journaliste) & \esp{manipular} (manipuler) & \esp{el respeto} (respect) \\
\esp{la noticia} (nouvelle) & \esp{censurar} (censurer) & \esp{la justicia} (justice) \\
\esp{el titular} (titre/gros titre) & \esp{compartir} (partager) & \esp{la no violencia} (non-violence) \\
\esp{la fuente} (source) & \esp{publicar} (publier) & \esp{la cultura de paz} (culture de paix) \\
\esp{la desinformación} (désinformation) & \esp{denunciar} (dénoncer) & \esp{los derechos humanos} (droits de l'homme) \\
\esp{la fake news / bulo} (infox / rumeur) & \esp{concienciar} (sensibiliser) & \esp{la ciudadanía} (citoyenneté) \\ \hline
\end{tabularx}
\end{center}

\courssub{Grammaire de l'injonction : l'impératif et les valeurs du subjonctif}

Le texte final est un modèle de \cle{texte injonctif}. Il utilise trois procédés grammaticaux majeurs pour appeler à l'action. Maîtrisez-les pour votre production écrite.

\begin{propriete}[L'impératif (Imperativo) : formes et emplois]
L'impératif exprime un ordre, un conseil, une prière, une interdiction. Il n'existe qu'aux personnes \textbf{tú, vosotros, nosotros, ustedes}.
\begin{center}
\begin{tabular}{|l|c|c|c|c|}\hline
\rowcolor{popblueL} \textbf{Verbe} & \textbf{tú} & \textbf{vosotros} & \textbf{nosotros} & \textbf{ustedes} \\ \hline
\esp{usar} (régulier -ar) & \esp{usa} / \esp{no uses} & \esp{usad} / \esp{no uséis} & \esp{usemos} / \esp{no usemos} & \esp{usen} / \esp{no usen} \\ \hline
\esp{verificar} (régulier -ar) & \esp{verifica} / \esp{no verifiques} & \esp{verificad} / \esp{no verifiquéis} & \esp{verifiquemos} / \esp{no verifiquemos} & \esp{verifiquen} / \esp{no verifiquen} \\ \hline
\esp{escuchar} (régulier -ar) & \esp{escucha} / \esp{no escuches} & \esp{escuchad} / \esp{no escuchéis} & \esp{escuchemos} / \esp{no escuchemos} & \esp{escuchen} / \esp{no escuchen} \\ \hline
\esp{dialogar} (régulier -ar) & \esp{dialoga} / \esp{no dialogues} & \esp{dialogad} / \esp{no dialoguéis} & \esp{dialoguemos} / \esp{no dialoguemos} & \esp{dialoguen} / \esp{no dialoguen} \\ \hline
\esp{compartir} (régulier -ir) & \esp{comparte} / \esp{no compartas} & \esp{compartid} / \esp{no compartáis} & \esp{compartamos} / \esp{no compartamos} & \esp{compartan} / \esp{no compartan} \\ \hline
\esp{publicar} (régulier -ar) & \esp{publica} / \esp{no publiques} & \esp{publicad} / \esp{no publiquéis} & \esp{publicemos} / \esp{no publicemos} & \esp{publicen} / \esp{no publiquen} \\ \hline
\end{tabular}
\end{center}
\textbf{Rappel :} À la forme négative, \textbf{toutes les personnes prennent les terminaisons du subjonctif présent}. Pour \emph{nosotros}, l'impératif affirmatif et négatif sont identiques au subjonctif présent (\esp{usemos / no usemos}).
\end{propriete}

\begin{propriete}[Autres structures de l'injonction]
\begin{itemize}
\item \textbf{Hay que + infinitif} : impersonnel, valeur d'obligation générale, de conseil. \esp{Hay que verificar las fuentes.} (« Il faut vérifier les sources. »)
\item \textbf{Deber / Tener que + infinitif} : obligation morale (\esp{deber}) ou contrainte (\esp{tener que}). \esp{Los periodistas deben ser objetivos.}
\item \textbf{Que + Subjonctif présent} : souhait, exhortation, slogan (souvent sans sujet). \esp{¡Que viva la paz!} (« Vive la paix ! »), \esp{¡Que no haya más violencia!} (« Qu'il n'y ait plus de violence ! »).
\item \textbf{Infinitif comme impératif} (affiches, manuels, recettes) : \esp{Verificar antes de compartir.} (« Vérifier avant de partager. »)
\end{itemize}
\end{propriete}

\begin{attention}[Pièges grammaticaux fréquents : Por / Para, Ser / Estar]
\begin{itemize}
\item \textbf{Por vs Para} dans le texte :
  \begin{itemize}
  \item \esp{para unir / para dividir} : \cle{But, finalité} (destination) $\rightarrow$ \textbf{PARA}.
  \item \esp{por eso / por lo tanto} : \cle{Cause, conséquence} $\rightarrow$ \textbf{POR}.
  \item \esp{por ejemplo} : \cle{Exemplification} $\rightarrow$ \textbf{POR}.
  \item \esp{gracias a / por} (cause) : \esp{gracias a la radio} / \esp{por la radio}.
  \end{itemize}
\item \textbf{Ser vs Estar} dans le texte :
  \begin{itemize}
  \item \esp{no son buenos ni malos} (\esp{ser}) : \cle{Caractéristique essentielle, nature} des médias.
  \item \esp{están conectados} (implicite, \esp{estar}) : \cle{État, situation} (le monde \emph{est} connecté).
  \item \esp{es también una misión} (\esp{ser}) : \cle{Identité, définition} de la profession.
  \item \esp{usados con inteligencia} (participe passé \emph{usado} + \esp{estar} si passif d'état) : ici participe adjectival \esp{usadas} (accord avec \esp{redes sociales}), pas de \esp{estar} explicite.
  \end{itemize}
\end{itemize}
\end{attention}

\courssub{Méthode : Commenter un texte argumentatif (Bac LV2)}

\begin{methode}[Étapes du commentaire composé]
\begin{enumerate}
\item \textbf{Introduction (1 paragraphe)} :
  \begin{itemize}
  \item \textbf{Accroche} : contextualiser le thème (médias, paix, désinformation) + référence au texte (titre, auteur fictif, source, date supposée, genre : article d'opinion).
  \item \textbf{Sujet} : de quoi parle le texte ? (Le double pouvoir des médias : diviser ou unir).
  \item \textbf{Thèse} : quelle est la position de l'auteur ? (Les médias sont des outils neutres ; l'usage responsable construit la paix).
  \item \textbf{Annonce du plan} : 2 ou 3 parties logiques (ex: I. Les dangers de la désinformation ; II. La responsabilité journalistique comme levier de paix ; III. L'appel citoyen à une utilisation éthique).
  \end{itemize}
\item \textbf{Développement (2 ou 3 parties, 1 paragraphe chacune)} :
  \begin{itemize}
  \item Chaque partie = 1 idée-force (argument principal du texte).
  \item Structure interne : \textbf{Idée} $\rightarrow$ \textbf{Citation courte} (intégrée dans la phrase) $\rightarrow$ \textbf{Analyse} (vocabulaire, connecteurs, figures de style, grammaire : impératif, subjonctif, connecteurs).
  \item Utiliser les connecteurs logiques : \esp{En primer lugar}, \esp{Además}, \esp{Por el contrario}, \esp{Por consiguiente}.
  \end{itemize}
\item \textbf{Conclusion (1 paragraphe)} :
  \begin{itemize}
  \item \textbf{Bilan} : résumer la démonstration (thèse confirmée).
  \item \textbf{Ouverture} : lien avec l'actualité (Cameroun, Guinée équatoriale, réseaux sociaux), avis personnel nuancé, citation courte (ex: \esp{« La pluma es la lengua del alma »} — Cervantes).
  \end{itemize}
\end{enumerate}
\end{methode}

\courssub{Commentaire modèle du texte « Los medios de comunicación... »}

\textbf{Introduction}\par
Dans un monde hyperconnecté où l'information circule à la vitesse de la lumière, les médias occupent une place centrale dans la fabrique de l'opinion publique. L'article d'opinion « Los medios de comunicación, ¿armas de guerra o instrumentos de paz? », publié dans la revue scolaire \emph{La Voz del Estudiante}, interroge cette responsabilité. L'auteur y défend la thèse selon laquelle les médias ne sont ni bons ni mauvais par nature : leur impact dépend de l'usage éthique que les professionnels et les citoyens en font. Nous analyserons d'abord les dangers de la désinformation (I), puis le rôle pacificateur d'un journalisme rigoureux (II), pour finir sur l'appel final à une citoyenneté numérique responsable (III).

\textbf{Développement}\par
\textbf{I. La désinformation : une arme de guerre moderne.}\par
Le texte s'ouvre sur une tonalité alarmante : \esp{« una noticia falsa puede encender un conflicto en pocas horas »}. L'hyperbole \esp{« en pocas horas »} et le verbe \esp{encender} (métaphore de l'incendie) soulignent la viralité destructrice des \esp{fake news}. L'exemple du \esp{« mensaje anónimo en una red social »} ancre le propos dans la réalité du smartphone, objet familier de l'élève camerounais. L'opposition marquée par \esp{« sin embargo »} (l. 4) pivote vers la thèse positive, mais le premier paragraphe a posé le diagnostic : sans vérification, le média devient \esp{arma de guerra}.

\textbf{II. Le journalisme responsable : une mission de paix.}\par
Le deuxième paragraphe valorise la figure du \esp{periodista responsable} dont l'éthique repose sur la \esp{verificación de fuentes}. L'antithèse \esp{« calmar los ánimos o, al contrario, aumentar la violencia »} montre le pouvoir performatif du \esp{titular}. La métaphore professionnelle \esp{« su profesión es también una misión de paz »} élève le journaliste au rang d'artisan de la \esp{convivencia}. L'objectivité (\esp{objetividad}) est présentée comme condition \emph{sine qua non} de la compréhension mutuelle (\esp{comprenden mejor a sus vecinos}) et de la baisse de la peur (\esp{el miedo disminuye}).

\textbf{III. De la contrainte à l'appel citoyen : la grammaire de l'engagement.}\par
Le connecteur \esp{« además »} ajoute l'argument des réseaux sociaux comme levier de \esp{tolerancia} et \esp{reconciliación}. L'exemple concret des \esp{radios comunitarias en Colombia} (post-accords de 2016) ancre l'universel dans l'historique. La thèse est alors explicitée (\esp{« todo depende del uso »}), suivie d'une série de parallélismes antithétiques : \esp{« gritar odio / cantar la fraternidad »}, \esp{« mostrar la guerra / enseñar la cultura del otro »}. Le texte bascule enfin dans l'injonction directe : l'impératif \emph{nosotros} (\esp{usemos, verifiquemos, escuchemos}) inclut l'auteur et le lecteur dans une communauté de destin. La phrase finale \esp{« cada uno de nosotros es responsable »} individualise la responsabilité collective.

\textbf{Conclusion}\par
Ce texte démontre avec rigueur que la technique médiatique est neutre, mais que son usage est éthique. Il rejoint l'urgence camerounaise où la rumeur sur WhatsApp peut embraser un quartier. L'ouverture vers une « culture de la paix » (\esp{cultura de paz}), notion portée par l'UNESCO, invite chaque lycéen à devenir un \esp{ciudadano digital} critique. Comme l'écrivait Machado : \esp{« Caminante, no hay camino, se hace camino al andar »} — la paix se construit, clic après clic, par la vérification et le dialogue.

\courssub{Production guidée : Écrire un appel à la paix (Texte injonctif)}

\begin{methode}[Rédiger un appel / un slogan (Texte injonctif)]
\begin{enumerate}
\item \textbf{Cible et ton} : Qui parle ? (Nous, les jeunes / Je, citoyen). Ton : solennel, mobilisateur, positif.
\item \textbf{Verbes d'action} : Choisir 3 à 5 verbes forts à l'impératif \emph{nosotros} (inclusif) ou \emph{ustedes} (appel large) : \esp{usemos, verifiquemos, escuchemos, construyamos, difundamos, rechacemos, promovamos}.
\item \textbf{Structure binaire} : Opposer le négatif à rejeter et le positif à construire (\esp{no para dividir, sino para unir}).
\item \textbf{Connecteurs} : \esp{Por eso, Por tanto, ¡Basta ya de...!}
\item \textbf{Rhétorique} : Parallélisme, anaphore (répétition en tête de phrase), métaphore forte.
\end{enumerate}
\end{methode}

\begin{exemplebox}[Modèle d'appel à la paix (production d'élève)]
\textbf{Titre :} \esp{¡Construyamos la paz en nuestras pantallas!}\par
\esp{Ante el ruido de la mentira, elijamos la verdad.}\par
\esp{Ante el grito del odio, cantemos la fraternidad.}\par
\esp{¡Verifiquemos antes de reenviar!}\par
\esp{¡Escuchemos la voz del otro, aunque nos moleste!}\par
\esp{¡Transformemos cada 'like' en un gesto de respeto!}\par
\esp{La paz no cae del cielo: se teclea, se comparte, se vive.}\par
\esp{¡Camerún, Guinea Ecuatorial, el mundo hispano: unidos por la información veraz!}
\end{exemplebox}

\begin{exercice}[À vous de jouer : Écrivez votre appel]
Rédigez un court texte injonctif (80--100 mots) intitulé \esp{« Jóvenes por la paz digital »}. Vous êtes délégué(e) de votre lycée à Yaoundé ou Douala. Adressez-vous aux élèves camerounais et hispanophones.
\textbf{Contraintes :}
\begin{itemize}
\item Utilisez au moins \textbf{4 impératifs} (dont 2 à \emph{nosotros}).
\item Intégrez \textbf{5 mots du lexique} du tableau (p. ex. \esp{verificar, difundir, convivencia, tolerancia, desinformación}).
\item Employez \textbf{une structure \esp{que + subjuntivo}} et \textbf{une structure \esp{hay que + infinitivo}}.
\item Marquez l'opposition \esp{no ... sino ...} ou \esp{no para ... sino para ...}.
\end{itemize}
\end{exercice}

\begin{corrige}[Exercice : Jóvenes por la paz digital (Proposition de correction)]
\esp{¡Jóvenes de Camerún y del mundo hispano, construyamos la paz digital!}\par
\esp{No dejemos que la desinformación divida nuestros barrios. ¡Verifiquemos cada noticia antes de compartirla!}\par
\esp{Hay que contrastar las fuentes: un clic responsable evita un conflicto real.}\par
\esp{¡Difundamos mensajes de tolerancia y convivencia en WhatsApp y Facebook!}\par
\esp{¡Que la verdad sea nuestra única viralidad!}\par
\esp{No usemos las redes para destruir, sino para unir corazones. ¡La paz empieza en tu pantalla!}
\end{corrige}

\courssub{Entraînement : Thème (Français $\rightarrow$ Espagnol)}

Traduisez les phrases suivantes en espagnol en mobilisant le lexique, la grammaire (impératif, subjonctif, por/para, ser/estar) et les connecteurs vus dans la leçon.

\begin{exercice}[Thème : Médias et paix]
\begin{enumerate}
\item Les médias ne sont pas des armes de guerre, ce sont des instruments de paix. (ser)
\item Il faut vérifier les sources avant de publier une nouvelle. (hay que / antes de + infinitif)
\item Les journalistes responsables informent avec objectivité pour calmer les esprits. (para + infinitif)
\item Que la radio communautaire favorise la réconciliation en Colombie ! (que + subjonctif)
\item Utilisons les réseaux sociaux pour diffuser la tolérance, pas pour propager la haine. (nosotros imperatif / no para... sino para)
\item Une fausse information peut déclencher un conflit entre deux communautés. (encender / entre)
\item Chaque citoyen est responsable de ce qu'il partage sur son écran. (ser responsable de / lo que + subjonctif)
\item Écoutons toutes les voix pour construire une culture de paix. (imperatif nosotros / para + infinitif)
\end{enumerate}
\end{exercice}

\begin{corrige}[Thème : Médias et paix]
\begin{enumerate}
\item \esp{Los medios no son armas de guerra, son instrumentos de paz.}
\item \esp{Hay que verificar las fuentes antes de publicar una noticia.}
\item \esp{Los periodistas responsables informan con objetividad para calmar los ánimos.}
\item \esp{¡Que la radio comunitaria favorezca la reconciliación en Colombia!}
\item \esp{Usemos las redes sociales para difundir la tolerancia, no para propagar el odio.}
\item \esp{Una noticia falsa puede encender un conflicto entre dos comunidades.}
\item \esp{Cada ciudadano es responsable de lo que comparte en su pantalla.}
\item \esp{Escuchemos todas las voces para construir una cultura de paz.}
\end{enumerate}
\end{corrige}

\begin{aretenir}[Synthèse de la leçon E22 — Ce qu'il faut retenir pour le Bac]
\begin{itemize}
\item \textbf{Genre textuel} : Article d'opinion (\esp{artículo de opinión}) = argumentatif + injonctif final.
\item \textbf{Structure} : Titre interrogatif $\rightarrow$ Constat (dangers) $\rightarrow$ Arguments (responsabilité, exemples) $\rightarrow$ Thèse (neutralité technique, responsabilité humaine) $\rightarrow$ Appel à l'action (impératifs \emph{nosotros}).
\item \textbf{Lexique clé} : \esp{prensa, informativo, noticia, titular, fuente, red social, desinformación, veracidad, diálogo, reconciliación, convivencia, tolerancia, cultura de paz}.
\item \textbf{Grammaire de l'injonction} : Impératif (\emph{usemos, verifiquemos}), \esp{hay que + inf.}, \esp{que + subj.}, \esp{deber + inf.}, Infinitif d'affiche.
\item \textbf{Connecteurs} : \esp{sin embargo, al contrario, además, por eso, por lo tanto, por ejemplo, para que + subj.}.
\item \textbf{Pièges francophones} : \esp{noticia} (info) $\neq$ notice ; \esp{por/para} (cause/but) ; \esp{ser/estar} (nature/état) ; \esp{compartir} (partager) ; accord \esp{usadas} (fém. pl. avec \esp{redes sociales}).
\item \textbf{Compétences Bac} : Commentaire structuré (Intro/Thèse/Plan, Développement citation+analyse, Conclusion/Ouverture) ; Thème/Version précis ; Production injonctive (appel, slogan, lettre ouverte).
\end{itemize}
\end{aretenir}

\coursec{Texte support : « Los medios de comunicación, ¿armas de guerra o instrumentos de paz? »}

\begin{textebox}[Los medios de comunicación, ¿armas de guerra o instrumentos de paz?]
Los medios de comunicación, ¿armas de guerra o instrumentos de paz? La pregunta no es inocente. En muchos países, la prensa, la radio y la televisión informan cada día a millones de personas, y las redes sociales multiplican esa información en pocos segundos. Sin embargo, esa rapidez tiene un precio: una noticia falsa puede encender un conflicto en una comunidad entera. Hemos visto rumores que han provocado violencia en mercados, en escuelas e incluso entre vecinos que vivían en paz desde hacía años.

Frente a este peligro, el periodista responsable verifica sus fuentes antes de publicar. Además, los medios serios dan voz a las víctimas, es decir, permiten que las personas que sufren cuenten su historia y sean escuchadas. Así, la información verdadera calma los ánimos y abre el diálogo entre los grupos enfrentados.

Las redes sociales también pueden difundir la tolerancia: jóvenes de regiones diferentes comparten sus canciones, sus recetas y sus sueños, y descubren que se parecen mucho más de lo que creían.

Por eso, los medios pueden construir la paz si informan con verdad. ¡Usemos los medios para dialogar y no para dividir!
\end{textebox}

\courssub{Traduction et glossaire commenté}

\begin{center}
\begin{tabularx}{\linewidth}{|X|X|X|}
\hline
\rowcolor{popblueL} \textbf{Español} & \textbf{Français} & \textbf{Note linguistique / Contexte} \\
\hline
\esp{la prensa} & la presse (écrite) & Ne pas confondre avec \esp{la presión} (la pression). \esp{Prensa} vient de \esp{prensa} (presse à imprimer). \\
\hline
\esp{las redes sociales} & les réseaux sociaux & \esp{Red} = filet, réseau. \esp{Social} = social. Accord au féminin pluriel. \\
\hline
\esp{multiplicar} & multiplier, diffuser massivement & Verbe régulier en \esp{-ar}. Ici : \esp{multiplican} (indic. présent, 3e pers. pl.). \\
\hline
\esp{una noticia falsa} & une fausse nouvelle, une \emph{fake news} & \esp{Noticia} = nouvelle (information). Faux ami : \esp{noticia} $\neq$ \emph{notice} (avis). \\
\hline
\esp{encender un conflicto} & déclencher / allumer un conflit & Métaphore : \esp{encender} (allumer un feu, une lumière). \esp{Encender} $\rightarrow$ \esp{enciendo} (diphthongue). \\
\hline
\esp{el rumor} & la rumeur & \esp{Provocar violencia} : provoquer la violence (collocation forte). \\
\hline
\esp{los vecinos} & les voisins & \esp{Vecino} (voisin) vient de \esp{vecindad} (voisinage). \\
\hline
\esp{verificar sus fuentes} & vérifier ses sources & \esp{Verificar} (régulier). \esp{Fuente} = source (d'information, d'eau). \\
\hline
\esp{dar voz a las víctimas} & donner la parole aux victimes & Locution figée. \esp{Es decir} = c'est-à-dire (explicatif). \\
\hline
\esp{calmar los ánimos} & apaiser les esprits / les cœurs & \esp{Ánimo} = moral, courage, esprit. \esp{Calmar} (régulier). \\
\hline
\esp{los grupos enfrentados} & les groupes opposés / en conflit & Participe passé \esp{enfrentado} (de \esp{enfrentar}, se faire face). \\
\hline
\esp{difundir la tolerancia} & diffuser la tolérance & \esp{Difundir} (irrégulier : \esp{difundo}, \esp{difunda} subj.) = diffuser, répandre. \\
\hline
\esp{las recetas / los sueños} & les recettes / les rêves & Vocabulaire concret pour illustrer l'échange culturel. \\
\hline
\esp{dialogar / dividir} & dialoguer / diviser & \esp{Dialogar} (régulier). \esp{Dividir} (régulier en \esp{-ir}). \\
\hline
\end{tabularx}
\end{center}

\begin{attention}[Faux amis fréquents : Médias \& Information]
\begin{itemize}
\item \esp{Actualidad} = l'actualité (ensemble des événements) ; \textbf{pas} \emph{actualité} au sens de « modernité » (\esp{modernidad}).
\item \esp{Reportaje} = reportage, enquête journalistique ; \textbf{pas} \emph{rapport} (\esp{informe}).
\item \esp{Sensible} = sensible (qui ressent), raisonnable ; \textbf{pas} \emph{sensible} au sens de « délicat » (\esp{delicado}, \esp{espinoso}).
\item \esp{Realizar} = réaliser (accomplir), concrétiser ; pour « réaliser » au sens de « comprendre », on dit \esp{darse cuenta de}.
\item \esp{Noticia} = nouvelle (info) ; \esp{Novedad} = nouveauté (chose nouvelle).
\end{itemize}
\end{attention}

\coursec{Compréhension du texte : thèse, arguments, structure}

\courssub{Sujet vs Thèse : distinction cruciale pour le Bac}

\begin{definition}[Sujet et Thèse]
Le \cle{sujet} (el \esp{tema}) est le thème général du texte : \textbf{de quoi parle-t-on ?} Ici : \esp{los medios de comunicación} (les médias).
La \cle{thèse} (la \esp{tesis}) est l'opinion défendue par l'auteur : \textbf{que dit-il sur le sujet ?} Elle se trouve souvent en conclusion, introduite par un connecteur conclusif (\esp{por eso}, \esp{por tanto}, \esp{en conclusión}).
\end{definition}

\begin{attention}[Piège d'examen]
Ne jamais répondre « Le texte parle des médias » si la question est \esp{¿Cuál es la tesis?} (Quelle est la thèse ?). Cela vaut 0 point. Il faut citer la thèse : \esp{« Los medios pueden construir la paz si informan con verdad. »}
\end{attention}

\courssub{Repérage des arguments et connecteurs logiques}

Le texte suit une structure \cle{argumentative} classique : Problème $\rightarrow$ Arguments / Contre-arguments $\rightarrow$ Thèse/Conclusion.

\begin{propriete}[Les connecteurs logiques : la carte d'identité du texte argumentatif]
\begin{center}
\begin{tabularx}{\linewidth}{|X|X|X|}
\hline
\rowcolor{popblueL} \textbf{Connecteur (Español)} & \textbf{Français} & \textbf{Fonction rhétorique dans le texte} \\
\hline
\esp{La pregunta no es inocente} & La question n'est pas anodine & \textbf{Accroche (Hook)} : annonce le débat. \\
\hline
\esp{Sin embargo} & Cependant, toutefois & \textbf{Concession / Tournant} : introduit le risque (vitesse = danger). \\
\hline
\esp{Frente a este peligro} & Face à ce danger & \textbf{Opposition} : introduit la solution (journalisme responsable). \\
\hline
\esp{Además} & De plus, en outre & \textbf{Addition} : renforce l'argument 2 (donner la parole). \\
\hline
\esp{Es decir} & C'est-à-dire & \textbf{Explication / Reformulation} : définit « dar voz ». \\
\hline
\esp{Así} & Ainsi, de cette façon & \textbf{Conséquence} : résultat de l'info vraie (calme, dialogue). \\
\hline
\esp{También} & Aussi & \textbf{Extension} : argument 3 (réseaux sociaux positifs). \\
\hline
\esp{Por eso} & C'est pourquoi, donc & \textbf{Conclusion / Thèse} : synthèse + ouverture injonctive. \\
\hline
\end{tabularx}
\end{center}

\begin{popfigure}
\begin{tikzpicture}[
node distance=0.5cm,
main/.style={draw=popblue, fill=popblueL, rounded corners, align=center, text width=5.8cm, inner sep=4pt, font=\small\bfseries},
arg/.style={draw=popgreen, fill=popgreenL, rounded corners, align=center, text width=5.8cm, inner sep=4pt, font=\small},
conc/.style={draw=poporange, fill=poporangeL, rounded corners, align=center, text width=5.8cm, inner sep=4pt, font=\small\bfseries},
arrow/.style={-{Stealth[length=2.5mm]}, thick, popdark}]
\node[main, align=center] (t){\textbf{Problématique (Titre)}\\ ¿Armas de guerra o instrumentos de paz?};
\node[arg, below=of t, align=center] (a1){\textbf{Arg. 1 : Danger}\\ \esp{Una noticia falsa puede encender un conflicto}\\ \emph{Ex: rumores $\rightarrow$ violencia en mercados, escuelas}};
\node[arg, below=of a1, align=center] (a2){\textbf{Arg. 2 : Solution pro} \\ \esp{Periodista verifica fuentes}\\ \esp{Medios serios dan voz a víctimas}\\ \emph{Conséquence: calmar ánimos, abrir diálogo}};
\node[arg, below=of a2, align=center] (a3){\textbf{Arg. 3 : Potentiel RS} \\ \esp{Redes sociales difunden tolerancia}\\ \emph{Ex: jóvenes comparten cultura/sueños}};
\node[conc, below=of a3, align=center] (c){\textbf{Thèse + Appel à l'action}\\ \esp{Por eso, los medios pueden construir la paz si informan con verdad}\\ \esp{¡Usemos los medios para dialogar y no para dividir!}};
\draw[arrow] (t) -- (a1) node[midway, right, font=\tiny, popdark] {\esp{Sin embargo}};
\draw[arrow] (a1) -- (a2) node[midway, right, font=\tiny, popdark] {\esp{Frente a / Además}};
\draw[arrow] (a2) -- (a3) node[midway, right, font=\tiny, popdark] {\esp{También}};
\draw[arrow] (a3) -- (c) node[midway, right, font=\tiny, popdark] {\esp{Por eso}};
\end{tikzpicture}
\legende{Structure argumentative : du problème à l'injonction finale.}
\end{popfigure}

\end{propriete}

\courssub{Compréhension fine : méthodologie et questions types Bac}

\begin{methode}[Répondre à une question de compréhension (Écrit / Oral)]
\begin{enumerate}
\item \textbf{Analyser la question} : Identifier le mot-clé (\esp{¿Por qué?} $\rightarrow$ cause ; \esp{¿Qué significa?} $\rightarrow$ définition ; \esp{¿Cuál?} $\rightarrow$ choix/identification).
\item \textbf{Localiser l'information} : Relire le paragraphe concerné. Souligner la phrase exacte.
\item \textbf{Citer} : Copier le segment textuel entre guillemets espagnols \esp{« ... »}.
\item \textbf{Expliquer / Reformuler} : Réécrire avec \emph{vos propres mots} (en espagnol si la question est en espagnol, en français si elle est en français). Ne pas recopier bêtement.
\end{enumerate}
\textbf{Règle d'or} : \cle{Citation + Explication = Points complets}. Citation seule = incomplet. Explication seule = non justifié.
\end{methode}

\begin{exoresolu}[Question type : Cause / Conséquence]
Énoncé : \esp{Según el texto, ¿por qué es peligrosa la rapidez de las redes sociales?}
\tcblower
\textbf{Citation} : \esp{« una noticia falsa puede encender un conflicto en una comunidad entera »} (línea 4-5).
\textbf{Explication} : \esp{La rapidez es peligrosa porque permite que una información falsa (fake news) se difunda masivamente antes de ser verificada, lo que puede provocar violencia real entre personas, como pasó en mercados y escuelas.} \\
\textit{Traduction : La rapidité est dangereuse car elle permet à une fausse info de se diffuser massivement avant vérification, ce qui peut provoquer une violence réelle (ex: marchés, écoles).}
\end{exoresolu}

\begin{exoresolu}[Question type : Sens d'expression / Métaphore]
Énoncé : \esp{Explique la expresión « dar voz a las víctimas » basándote en el texto.}
\tcblower
\textbf{Citation} : \esp{« es decir, permiten que las personas que sufren cuenten su historia y sean escuchadas »} (línea 8-9).
\textbf{Explication} : \esp{Significa que los medios deben ofrecer un espacio público a quienes sufren para que narren su experiencia directamente, sin intermediarios, y así sean oídas por la sociedad. Es una metáfora: la "voz" representa el derecho a la palabra pública.}
\end{exoresolu}

\begin{exoresolu}[Question type : Grammaire dans le texte / Valeur du mode]
Énoncé : \esp{Analiza la forma verbal « usemos » en la última frase. ¿Qué modo y tiempo? ¿Qué valor expresa?}
\tcblower
\textbf{Forme} : \esp{Usemos} = \textbf{Subjonctif présent}, 1ère personne du pluriel, du verbe \esp{usar} (régulier -ar).
\textbf{Valeur} : \cle{Valeur hortative / injonctive collective}. Comme l'impératif n'existe pas à la 1ère personne du pluriel en espagnol, on utilise le subjonctif présent pour exprimer un appel à l'action incluant l'énonciateur : « \textbf{Utilisons} les médias... » (équivalent français : impératif 1ère pers. pl. ou « Il faut que nous utilisions »).
\textbf{Contraste} : Impératif 2e pers. \esp{¡Usad!} (vous) / \esp{¡Usa!} (tu) vs Subjonctif hortatif \esp{¡Usemos!} (nous).
\end{exoresolu}

\courssub{Entraînement : Questions à traiter}

\begin{exercice}[Compréhension écrite - Niveau Bac]
\begin{enumerate}
\item \esp{¿Cuál es la tesis central del texto? Cítala textualmente entre comillas.}
\item \esp{Identifica tres conectores lógicos del texto y explica la función de cada uno.}
\item \esp{¿Qué papel juega el periodista responsable según el autor? Cita y explica.}
\item \esp{El título es una pregunta retórica. Explica qué es y qué efecto produce en el lector.}
\item Traduction (Version) : \esp{« Frente a este peligro, el periodista responsable verifica sus fuentes antes de publicar. Además, los medios serios dan voz a las víctimas... »} (Jusqu'à \esp{...abre el diálogo entre los grupos enfrentados.}).
\end{enumerate}
\end{exercice}

\begin{corrige}[Exercice Compréhension]
\begin{enumerate}
\item \esp{« Los medios pueden construir la paz si informan con verdad. »} (Dernier paragraphe, introduit par \esp{Por eso}).
\item Exemples : \esp{Sin embargo} (concession/opposition) ; \esp{Además} (addition d'argument) ; \esp{Por eso} (conséquence/conclusion) ; \esp{Es decir} (explication) ; \esp{Frente a} (opposition spatiale/argumentative).
\item \textbf{Citation} : \esp{« el periodista responsable verifica sus fuentes antes de publicar »} / \esp{« los medios serios dan voz a las víctimas »}. \textbf{Explication} : Il a un rôle déontologique : vérifier (rigueur) et donner la parole aux sans-voix (éthique), ce qui apaise les tensions.
\item Une question rhétorique ne demande pas de réponse mais provoque la réflexion. Ici, l'antithèse \esp{armas de guerra / instrumentos de paz} cadre le débat : les médias sont ambivalents.
\item \textbf{Version} : « Face à ce danger, le journaliste responsable vérifie ses sources avant de publier. De plus, les médias sérieux donnent la parole aux victimes, c'est-à-dire qu'ils permettent aux personnes qui souffrent de raconter leur histoire et d'être entendues. Ainsi, l'information vraie apaise les esprits et ouvre le dialogue entre les groupes opposés. »
\end{enumerate}
\end{corrige}

\coursec{Lexique thématique : Médias, Information, Culture de la paix}

\courssub{Familles de mots et collocations (Mots-outils pour l'expression)}

\begin{definition}[Champs lexicaux à maîtriser pour le Bac]
\begin{itemize}
\item \textbf{Acteurs} : \esp{el periodista} (journaliste), \esp{el reportero / la reportera} (reporter), \esp{el corresponsal} (correspondant), \esp{el editor / la editora} (rédacteur en chef / éditeur), \esp{el locutor / la locutora} (speaker, animateur radio), \esp{el influencer / la influencer} (influenceur - anglicisme admis), \esp{el ciudadano} (citoyen).
\item \textbf{Supports} : \esp{el periódico / el diario} (quotidien), \esp{la revista} (magazine), \esp{la radio}, \esp{la televisión (la tele)}, \esp{las redes sociales}, \esp{el medio digital}, \esp{el blog}, \esp{el podcast}.
\item \textbf{Contenus} : \esp{la noticia} (info), \esp{el titular} (titre/manchette), \esp{el artículo} (article), \esp{el reportaje} (reportage), \esp{la entrevista} (interview), \esp{el editorial} (éditorial), \esp{la columna de opinión} (chronique), \esp{la publicidad} (pub), \esp{el bulo / la fake news} (intox, rumeur volontaire), \esp{la desinformación} (désinformation).
\item \textbf{Verbes d'action médiatique} : \esp{informar}, \esp{comunicar}, \esp{difundir} (diffuser), \esp{publicar}, \esp{emitir} (diffuser à l'antenne), \esp{cubrir un evento} (couvrir), \esp{investigar}, \esp{verificar} / \esp{contrastar} (recouper), \esp{censurar}, \esp{manipular}, \esp{viralizar}.
\item \textbf{Paix \& Conflit} : \esp{la paz}, \esp{el conflicto}, \esp{la violencia}, \esp{la guerra}, \esp{el diálogo}, \esp{la reconciliación}, \esp{la tolerancia}, \esp{el respeto}, \esp{la convivencia} (vivre-ensemble), \esp{la mediación} (médiation), \esp{el acuerdo de paz}, \esp{cesar el fuego}, \esp{construir la paz}, \esp{fomentar la cultura de paz}.
\end{itemize}
\end{definition}

\begin{exemplebox}[Collocations utiles pour l'oral et l'écrit (Thème / Version)]
\begin{itemize}
\item \esp{Verificar la veracidad de una noticia} (Vérifier la véracité d'une info).
\item \esp{Luchar contra la desinformación} (Lutter contre la désinformation).
\item \esp{Dar visibilidad a un problema} (Donner de la visibilité / Médiatiser un problème).
\item \esp{Generar debate público} (Générer un débat public).
\item \esp{Promover la cultura de paz} (Promouvoir la culture de paix).
\item \esp{Resolver conflictos de forma pacífica} (Résoudre les conflits pacifiquement).
\item \esp{Construir puentes, no muros} (Construire des ponts, pas des murs - métaphore politique fréquente).
\end{itemize}
\end{exemplebox}

\courssub{Exercices lexicaux}

\begin{exercice}[Lexique : Complétion et Traduction]
\begin{enumerate}
\item Complète en espagnol : \esp{El \_\_\_\_\_\_ (journaliste) debe \_\_\_\_\_\_ (vérifier) las \_\_\_\_\_\_ (sources) para evitar \_\_\_\_\_\_ (la désinformation).}
\item Traduis en espagnol : « Les réseaux sociaux peuvent viraliser de fausses nouvelles et provoquer des conflits. »
\item Associe : 1. \esp{Editorial} \quad 2. \esp{Reportaje} \quad 3. \esp{Entrevista} \quad 4. \esp{Titular} \\
A. Manchette / Titre principal \quad B. Article d'opinion du journal \quad C. Enquête de terrain \quad D. Tête-à-tête questions-réponses.
\end{enumerate}
\end{exercice}

\begin{corrige}[Exercice Lexique]
\begin{enumerate}
\item \esp{El \textbf{periodista} debe \textbf{verificar / contrastar} las \textbf{fuentes} para evitar \textbf{la desinformación}.}
\item \esp{Las redes sociales pueden \textbf{viralizar noticias falsas} y \textbf{provocar conflictos}.} (Ou \esp{difundir bulos / generar conflictos}).
\item 1-B, 2-C, 3-D, 4-A.
\end{enumerate}
\end{corrige}

\coursec{Textes argumentatifs et injonctifs : Grammaire de l'appel à la paix}

\courssub{Repérage : Argumentatif vs Injonctif}

\begin{definition}[Typologie textuelle]
\begin{itemize}
\item \textbf{Texte argumentatif} : Vise à \textbf{convaincre} (logique, raisons, preuves). Structure : Thèse - Arguments - Exemples - Conclusion. Temps dominants : Indicatif (présent, passé, futur). Connecteurs : \esp{porque, por tanto, sin embargo, además}.
\item \textbf{Texte injonctif (ou exhortatif)} : Vise à \textbf{faire agir} (ordre, conseil, appel, supplique). Structure : Appel direct $\rightarrow$ Motivation (souvent implicite) $\rightarrow$ Action visée. Formes verbales spécifiques (voir ci-dessous). Ponctuation : \esp{¡ !}.
\end{itemize}
\end{definition}

\begin{attention}[Hybridation fréquente]
Notre texte support est \textbf{hybride} : il commence par l'argumentation (paragraphes 1 à 3) pour finir par une \textbf{conclusion injonctive} (dernière phrase : \esp{¡Usemos los medios para dialogar...!}). C'est la structure type de l'\emph{editorial} ou de la \emph{tribune libre}.
\end{attention}

\courssub{Grammaire de l'injonction : Les 4 formes de l'appel à l'action}

\begin{propriete}[Les 4 piliers de l'injonction en espagnol (Niveau Terminale)]
Pour \cle{formuler des injonctions et appels à la paix} (compétence visée), tu dois maîtriser ces 4 structures, de la plus directe à la plus indirecte :
\end{propriete}

\begin{center}
\begin{tabularx}{\linewidth}{|X|X|X|X|}
\hline
\rowcolor{popblueL} \textbf{Structure} & \textbf{Formation} & \textbf{Valeur / Registre} & \textbf{Exemple « Paix / Médias »} \\
\hline
\textbf{1. Impératif (Imperativo)} & \esp{Habla! / Hable! / Hablad! / Hablen!} & Ordre direct, conseil, instruction. \textbf{2e personne} (tú, usted, vosotros, ustedes). & \esp{¡Verifica las fuentes antes de compartir!} \\
\hline
\textbf{2. Subjonctif Hortatif (Exhortativo)} & \esp{Que + Subj. Pres.} / \esp{Subj. Pres. 1ère pl.} & Appel collectif (\esp{nosotros}), souhait, bénédiction, ordre indirect (3e pers.). & \esp{¡Usemos los medios para la paz!} (Nous) \\
& & & \esp{¡Que la verdad prevalezca!} (3e pers. - Que la vérité l'emporte !) \\
\hline
\textbf{3. Infinitif d'injonction (Infinitivo de mandato)} & \esp{Infinitif + !} & Règles, pancartes, modes d'emploi, recettes. Impersonnel, neutre. & \esp{¡No difundir rumores sin verificar!} \\
\hline
\textbf{4. Périphrases modales (Deber / Hay que / Tener que + Inf.)} & \esp{Deber / Hay que / Tener que + Infinitif} & Obligation morale, conseil fort, nécessité. Tous sujets. & \esp{Debemos construir puentes.} / \esp{Hay que informar con verdad.} \\
\hline
\end{tabularx}
\end{center}

\begin{aretenir}[Impératif vs Subjonctif : Points de vigilance pour francophones]
\begin{enumerate}
\item \textbf{Formes \esp{nosotros}} : \textbf{Jamais d'impératif} à la 1ère pers. pl. $\rightarrow$ Toujours \textbf{Subjonctif Présent} : \esp{¡Hablemos! ¡Vayamos! ¡Durmamos! ¡Usemos!} (Verbes irréguliers gardent leur irrégularité au subjonctif).
\item \textbf{Négation} : L'impératif \textbf{négatif} utilise \textbf{TOUJOURS le subjonctif} pour toutes les personnes : \esp{¡No hables! ¡No hable! ¡No hablemos! ¡No habléis! ¡No hablen!}.
\item \textbf{Pronoms} : Affirmatif $\rightarrow$ enclitiques (accrochés) : \esp{¡Verifícalo! ¡Dímelo!} (accent écrit si >2 syllabes). Négatif $\rightarrow$ proclitiques (devant) : \esp{¡No lo verifiques! ¡No me lo digas!}.
\item \textbf{Verbes irréguliers fréquents (Impératif + / -)} :
\begin{center}
\begin{tabular}{|c|c|c|c|}
\hline
\rowcolor{poporangeL} \textbf{Infinitif} & \textbf{Imperativo + (tú)} & \textbf{Imperativo - (tú)} & \textbf{Subj. Pres. (nosotros)} \\
\hline
\esp{Ser} & \esp{Sé} & \esp{No seas} & \esp{Seamos} \\
\hline
\esp{Ir} & \esp{Ve} & \esp{No vayas} & \esp{Vayamos} \\
\hline
\esp{Tener} & \esp{Ten} & \esp{No tengas} & \esp{Tengamos} \\
\hline
\esp{Poner} & \esp{Pon} & \esp{No pongas} & \esp{Pongamos} \\
\hline
\esp{Salir} & \esp{Sal} & \esp{No salgas} & \esp{Salgamos} \\
\hline
\esp{Venir} & \esp{Ven} & \esp{No vengas} & \esp{Vengamos} \\
\hline
\esp{Hacer} & \esp{Haz} & \esp{No hagas} & \esp{Hagamos} \\
\hline
\esp{Decir} & \esp{Di} & \esp{No digas} & \esp{Digamos} \\
\hline
\end{tabular}
\end{center}
\end{enumerate}
\end{aretenir}

\begin{exemplebox}[Application : Transformer en appels à la paix]
\begin{itemize}
\item \esp{(Tú) Informa con verdad.} $\rightarrow$ \esp{¡Informa con verdad!} (Impératif +)
\item \esp{(Nosotros) Usar las redes para unir.} $\rightarrow$ \esp{¡Usemos las redes para unir!} (Subj. Hortatif)
\item \esp{(Ustedes) No compartir bulos.} $\rightarrow$ \esp{¡No compartan bulos!} (Impératif - / Subj. Ustedes)
\item \esp{(Regla general) Respetar la dignidad.} $\rightarrow$ \esp{¡Respetar la dignidad humana!} (Infinitif d'injonction)
\item \esp{Deber (nosotros) promover el diálogo.} $\rightarrow$ \esp{Debemos promover el diálogo.} / \esp{¡Promovamos el diálogo!}
\end{itemize}
\end{exemplebox}

\begin{exercice}[Grammaire de l'injonction : Mise en forme]
Mets le verbe entre parenthèses à la forme correcte pour exprimer un appel à l'action (choisis la personne indiquée).
\begin{enumerate}
\item \esp{\_\_\_\_\_\_ (vosotros) \_\_\_\_\_\_ (investigar) antes de publicar!} (Ordre direct, Espagne).
\item \esp{\_\_\_\_\_\_ (nosotros) \_\_\_\_\_\_ (construir) una cultura de paz!} (Appel collectif).
\item \esp{\_\_\_\_\_\_ (ustedes) \_\_\_\_\_\_ (no / difundir) noticias sin verificar!} (Interdiction formelle, AmLat).
\item \esp{\_\_\_\_\_\_ (regla general) \_\_\_\_\_\_ (contrastar) la información!} (Pancarte / Règle).
\item \esp{\_\_\_\_\_\_ (tú) \_\_\_\_\_\_ (ser) responsable con tu palabra!} (Conseil ami).
\end{enumerate}
\end{exercice}

\begin{corrige}[Exercice Grammaire Injonction]
\begin{enumerate}
\item \esp{¡Investigad antes de publicar!} (Impératif \esp{vosotros} -ar $\rightarrow$ -ad).
\item \esp{¡Construyamos una cultura de paz!} (Subj. \esp{nosotros} \esp{construir} $\rightarrow$ \esp{construyamos}).
\item \esp{¡No difundan noticias sin verificar!} (Impératif négatif \esp{ustedes} = Subj. \esp{ustedes}).
\item \esp{¡Contrastar la información!} (Infinitif d'injonction).
\item \esp{¡Sé responsable con tu palabra!} (Impératif irrégulier \esp{tú} : \esp{Sé}). Négatif : \esp{¡No seas...!}.
\end{enumerate}
\end{corrige}

\coursec{Méthode du commentaire de texte (Épreuve écrite du Baccalauréat)}

\begin{methode}[Commentaire de texte argumentatif : 4 étapes pour 20/20]
\begin{enumerate}
\item \textbf{INTRODUCTION (1 paragraphe)} \\
\underline{Accroche} : Contexte général (médias, paix, actualité camerounaise/hispanique). \\
\underline{Présentation du texte} : Titre (entre guillemets \esp{« ... »}), auteur (si connu, sinon « autor anónimo »), nature (article d'opinion, editorial, tribune), source, date. \\
\underline{Sujet} : \esp{El texto trata sobre...} (Le texte traite de...). \\
\underline{Thèse} : \esp{El autor defiende la idea de que... / La tesis central es: « ... »}. \\
\underline{Annonce du plan} : \esp{Analizaremos primero los peligros..., luego el papel del periodista..., finalmente el potencial de las redes...}
\item \textbf{DÉVELOPPEMENT (2 ou 3 paragraphes = 2 ou 3 arguments)} \\
Chaque paragraphe = \textbf{Un argument principal}.
\begin{itemize}
\item Phrase sujet (thème de l'argument).
\item Citation courte intégrée : \esp{El autor afirma que « una noticia falsa puede encender un conflicto »}.
\item Analyse / Commentaire : Expliquer la portée, le vocabulaire, les figures de style (métaphore \esp{encender}, antithèse \esp{armas/instrumentos}), le connecteur utilisé.
\item Lien avec le contexte (Cameroun, Guinée Équatoriale, monde).
\end{itemize}
\item \textbf{CONCLUSION (1 paragraphe)} \\
\underline{Bilan} : Résumer la démonstration en 1 phrase. \esp{En definitiva, el texto demuestra que...} \\
\underline{Ouverture} : Question élargie, lien avec l'actualité, avis personnel nuancé. \esp{Esto nos invita a reflexionar sobre el papel de cada ciudadano en la verificación de la información en la era digital.}
\item \textbf{RELECTURE} : Orthographe (accents, \esp{ñ}, \esp{¿¡}), accords, cohérence connecteurs, respect du nombre de mots (souvent 150-200 mots).
\end{enumerate}
\end{methode}

\begin{exemplebox}[Modèle de commentaire rédigé (Extrait - Introduction + Argument 1)]
\textbf{Introducción} \\
\esp{En la era de la información instantánea, los medios de comunicación se han convertido en el principal motor de la opinión pública. El texto « Los medios de comunicación, ¿armas de guerra o instrumentos de paz? », de autor anónimo, es un artículo de opinión que aborda la doble cara de los medios: su capacidad para destruir o construir la convivencia social. La tesis central del autor es clara: « Los medios pueden construir la paz si informan con verdad ». Para demostrarlo, el texto estructura su argumentación en tres tiempos: el peligro de la falsedad, la responsabilidad profesional y el potencial positivo de las redes sociales.}

\textbf{Desarrollo: Argumento 1 – La velocidad, arma de doble filo} \\
\esp{En primer lugar, el autor denuncia el riesgo inherente a la inmediatez digital. Mediante el conector adversativo « sin embargo », introduce la idea de que « esa rapidez tiene un precio ». La metáfora « encender un conflicto » compara la fake news a una chispa que incendia la convivencia. Los ejemplos concretos —« mercados », « escuelas », « vecinos »— anclan el peligro en lo cotidiano, recordándonos que en Camerún, como en España o Guatemala, un rumor en WhatsApp puede desencadenar violencia real. Así, el texto alerta: sin verificación, el medio se vuelve « arma de guerra ».}
\end{exemplebox}

\coursec{Production guidée : Écrire un appel à la paix (Expression écrite)}

\courssub{Consigne type Bac}
\esp{« Eres portavoz de un club de prensa juvenil en tu instituto. Escribe un llamamiento (150-180 palabras) para publicar en el blog del colegio con motivo del « Día Internacional de la Paz » (21 de septiembre). Debes: denunciar el peligro de los bulos, recordar el deber de verificar, y convocar a los estudiantes a usar las redes para el diálogo. Usa formas injuntivas variadas. »}

\begin{methode}[Rédaction d'un appel / Llamamiento (Structure type)]
\begin{enumerate}
\item \textbf{Titre accrocheur} : \esp{¡Por una red sin bulos! / ¡Jóvenes por la paz digital!}
\item \textbf{Accroche (Hook)} : Situation + Émotion / Constat. \esp{Cada día, miles de rumores...}
\item \textbf{Corps : 3 arguments-injonctions} (Utiliser \textbf{au moins 3 formes différentes} vues en grammaire).
\begin{itemize}
\item Arg 1 (Constat + Interdiction) : \esp{¡No compartáis noticias sin verificar!} (Impératif - / Subj.).
\item Arg 2 (Devoir moral) : \esp{Debemos / Hay que contrastar las fuentes.} (Périphrase).
\item Arg 3 (Appel collectif positif) : \esp{¡Usemos las redes para tender puentes!} (Subj. Hortatif \esp{nosotros}).
\end{itemize}
\item \textbf{Conclusion / Signature} : \esp{¡La paz empieza en tu muro! \textbf{Club de Prensa Juvenil.}}
\end{enumerate}
\end{methode}

\begin{exemplebox}[Modèle de production écrite (Llamamiento)]
\textbf{¡JÓVENES POR LA PAZ DIGITAL: VERIFICA ANTES DE COMPARTIR!}

\esp{¿Sabías que un bulo en WhatsApp puede provocar una pelea en el recreo o un conflicto entre barrios? En este Día de la Paz, desde el Club de Prensa Juvenil del Lycée de Yaoundé, lanzamos un reto: \textbf{¡Convertid vuestros muros en puentes!}}

\esp{\textbf{¡No reenviéis cadenas sin fuente fiable!} (Imperativo negativo). La desinformación es un virus invisible que destruye la confianza. \textbf{Hay que contrastar cada noticia en medios oficiales} (Perífrasis de obligación) antes de darle al botón de « compartir ». Un clic irresponsable duele; mil clics conscientes curan.}

\esp{\textbf{¡Usemos las redes para contar nuestras historias, no para inventar las de otros!} (Subjuntivo hortativo). Compartid vuestra cultura, vuestros sueños, vuestras soluciones. \textbf{Construyamos juntos una ciudadanía digital crítica y solidaria.} (Subjuntivo \esp{nosotros}).}

\esp{La paz no es un sueño, es una práctica diaria. \textbf{¡Verifiquemos, dialoguemos, unámonos!}}

\esp{\textbf{Club de Prensa Juvenil — Lycée Bilingue, Yaoundé.}}
\end{exemplebox}

\begin{exercice}[À ton tour : Production écrite]
Rédige ton propre \esp{llamamiento} (150 mots env.) sur l'un des deux sujets :
\begin{enumerate}
\item \esp{« Día Mundial de la Libertad de Prensa » (3 mayo) : Appel à protéger les journalistes et exiger une info libre/vérifiée.}
\item \esp{« Campaña contra el Ciberacoso » : Appel aux élèves pour signaler, ne pas partager, soutenir les victimes.}
\end{enumerate}
\textbf{Contraintes} : Titre, au moins 3 formes injonctives différentes (Impératif +/-, Subj. Hortatif, Infinitif, \esp{Hay que/Debemos}), vocabulaire du module (prensa, bulo, verificar, paz, diálogo, tolerancia), connecteurs (\esp{por tanto, además, en definitiva}).
\end{exercice}

\begin{corrige}[Pistes de correction Production Écrite]
\textbf{Grille d'auto-évaluation (Coche si présent) :}
\begin{itemize}
\item[ ] Titre entre \esp{¡ !} accrocheur.
\item[ ] Introduction : Contexte + Problème (Bulos / Ciberacoso).
\item[ ] Argument 1 : Forme injonctive 1 (ex: \esp{¡No compartáis...!} / \esp{¡Denunciad...!}).
\item[ ] Argument 2 : Forme injonctive 2 (ex: \esp{Hay que verificar...} / \esp{Debemos apoyar...}).
\item[ ] Argument 3 : Forme injonctive 3 (ex: \esp{¡Construyamos...!} / \esp{¡Protejamos...!}).
\item[ ] Conclusion : Slogan final + Signature.
\item[ ] Vocabulaire précis (min. 5 mots de la liste lexicale).
\item[ ] Orthographe : Accents, \esp{ñ}, \esp{¿¡}, concordance.
\end{itemize}
\end{corrige}

\courssub{Atelier oral : Simulation « Radio Paix »}

\begin{experience}[Jeu de rôle : « En directo : Construyendo la paz »]
\textbf{Contexte} : Vous animez une émission de radio locale (Yaoundé, Douala, Bertoua, Malabo...) spéciale « Culture de la paix ».
\textbf{Rôles (groupes de 4)} :
\begin{enumerate}
\item \esp{El/La Presentador/a} : Lance le sujet, gère le temps, conclut.
\item \esp{El/La Periodista de investigación} : Explique comment vérifier une info (outils : \esp{Google Lens, Maldita.es, Africa Check, sites officiels}).
\item \esp{Un/a Estudiante testigo} : Raconte une anecdote vraie (ou inventée plausible) : \esp{« Una vez vi un rumor en Facebook que decía... y casi causa una pelea... »}.
\item \esp{El/La Experto/a en mediación} : Donne 3 conseils concrets pour désamorcer un conflit né d'une rumeur.
\end{enumerate}
\textbf{Contraintes linguistiques} : Utiliser \textbf{au moins 2 subjonctifs hortatifs} (\esp{¡Hablemos de...! ¡No permitamos...!}), \textbf{1 infinitif d'injonction} (\esp{¡Verificar siempre!}), \textbf{connecteurs} (\esp{por tanto, frente a esto, en conclusión}).
\textbf{Durée} : 3 minutes par groupe. Enregistrement possible sur téléphone pour auto-évaluation.
\end{experience}

\begin{savaistu}[Cameroun \& Guinée Équatoriale : Médias et Paix]
\begin{itemize}
\item \textbf{Radio rurale} : Au Cameroun, les \esp{radios communautaires} (ex: Radio Medumba, Radio Yetem) émettent en langues nationales (Bamiléké, Fulfuldé, Bassa...). Elles sont cruciales pour la \esp{convivencia} (vivre-ensemble) et la prévention des conflits fonciers ou interethniques.
\item \textbf{Guinée Équatoriale} : Seule nation africaine hispanophone officielle. \esp{Radio Nacional de Guinea Ecuatorial} et \esp{TVGE} sont des vecteurs de l'espagnol en zone CEMAC. La \esp{Academia Ecuatoguineana de la Lengua Española} veille à la norme.
\item \textbf{Fact-checking local} : Des initiatives comme \esp{\#StopBulo237} (Cameroun) ou \esp{Africa Check} (bureau régional) forment les jeunes à la \esp{alfabetización mediática} (éducation aux médias).
\item \textbf{Journée clé} : \esp{21 de septiembre} (Journée internationale de la paix) \& \esp{3 de mayo} (Journée mondiale de la liberté de presse). Dates à connaître pour les sujets de rédaction !
\end{itemize}
\end{savaistu}

\courssub{Bilan de la leçon E22 : Check-list des compétences}

\begin{aretenir}[Objectifs validés ? Coche la case \fbox{} si maîtrisé]
\begin{itemize}
\item[ \fbox{} ] \textbf{Lecture} : Je sais identifier le sujet, la thèse, les arguments et les connecteurs d'un texte argumentatif sur les médias.
\item[ \fbox{} ] \textbf{Vocabulaire} : Je maîtrise 20+ mots du champ lexical (médias, paix, conflit) et j'évite les faux amis (\esp{noticia/actualidad/realizar}).
\item[ \fbox{} ] \textbf{Grammaire Injonction} : Je sais former et utiliser les 4 formes de l'appel à l'action (Impératif +/-, Subj. Hortatif, Infinitif, Périphrases \esp{Hay que/Debemos}).
\item[ \fbox{} ] \textbf{Commentaire} : Je connais la structure en 4 étapes (Intro, Développement, Conclusion, Relecture) et je sais citer + analyser.
\item[ \fbox{} ] \textbf{Expression Écrite} : Je peux rédiger un \esp{llamamiento} structuré, varié grammaticalement, et ancré dans mon contexte (Cameroun/Monde hispanophone).
\item[ \fbox{} ] \textbf{Expression Orale} : Je peux participer à un débat / jeu de rôle radio en utilisant le lexique et la grammaire de l'injonction.
\end{itemize}

\end{aretenir}

\begin{attention}[Dernier conseil pour le Bac]
Dans l'épreuve de \textbf{Version} (Espagnol $\rightarrow$ Français), attention au \esp{Por eso} : \textbf{Jamais} « Par ça ». Toujours \emph{C'est pourquoi / Donc / Par conséquent}.
Dans l'épreuve de \textbf{Thème} (Français $\rightarrow$ Espagnol), pour « Utilisons les médias pour la paix » : \textbf{¡Usemos los medios para la paz!} (Subjonctif \esp{nosotros}), \textbf{PAS} \esp{*Usamos} (Indicatif) ni \esp{*Usad} (Impératif \esp{vosotros} si contexte inclusif).
\end{attention}

\esp{¡Excelente trabajo! La lección E22 está completa. ¡A repasar para el Bac!}

\coursec{Lexique thématique : médias et culture de la paix}

Dans la section précédente, nous avons lu et compris le texte \esp{« Los medios de comunicación, ¿armas de guerra o instrumentos de paz? »}. Pour aller plus loin — commenter, traduire, débattre —, il nous faut maintenant des \cle{outils lexicaux} précis. Un bon commentaire de texte, à l'examen comme dans la vie, repose sur le choix exact des mots. Cette section construit donc méthodiquement deux grands champs lexicaux : celui des \cle{médias} et celui de la \cle{culture de la paix}, puis les verbes et expressions qui permettent de les relier.

\courssub{Observation : un texte pour faire émerger le vocabulaire}

Lisons d'abord ce court article de presse, rédigé dans un style journalistique authentique. Soulignez mentalement tous les mots qui appartiennent au monde des médias, puis tous ceux qui appartiennent au monde de la paix.

\begin{textebox}[Un debate en la redacción]
En la redacción del periódico \emph{La Voz del Centro}, en Yaoundé, los periodistas discuten cada mañana sobre la información del día. Hoy, el director ha propuesto un tema delicado: ¿cómo informar sobre un conflicto entre dos comunidades sin sembrar el odio?

— Tenemos que verificar todas las fuentes antes de publicar — dice María, una periodista con diez años de experiencia—. Una noticia falsa puede provocar violencia en la calle.

— Estoy de acuerdo — responde el redactor jefe—. El titular debe ser claro y honesto. Nada de propaganda: nuestro periódico debe promover el diálogo y la tolerancia.

En la televisión y en la radio, los informativos repiten las mismas imágenes del conflicto. En las redes sociales, circulan rumores sin verificar. Algunos canales prefieren el escándalo porque vende más que la verdad. Sin embargo, otros medios deciden dar voz a las víctimas y a los pacifistas: publican reportajes sobre la reconciliación y entrevistas con líderes que llaman a la convivencia y al respeto.

Al final del día, el periódico publica un editorial titulado \emph{« Construir la paz empieza por informar con verdad »}. El texto recuerda que la paz no es solo la ausencia de guerra: es también justicia, escucha y perdón. « Un periodista responsable — concluye el editorial — no describe solamente el conflicto: ayuda a superarlo. »

Los lectores, en los barrios de la ciudad, comentan el artículo. Muchos piensan que los medios de comunicación pueden ser armas de guerra, pero también instrumentos de paz. Todo depende de quienes escriben... y de quienes leen.
\end{textebox}

\textbf{Glossaire :} \esp{la redacción} = la rédaction (d'un journal) ; \esp{delicado} = délicat, sensible ; \esp{sembrar el odio} = semer la haine ; \esp{el redactor jefe} = le rédacteur en chef ; \esp{el rumor} = la rumeur ; \esp{el escándalo} = le scandale ; \esp{las víctimas} = les victimes ; \esp{el editorial} = l'éditorial ; \esp{la ausencia} = l'absence ; \esp{superar} = surmonter ; \esp{los lectores} = les lecteurs.

\textbf{Questions de repérage :}
\begin{enumerate}
\item Citez cinq mots du texte qui désignent des médias ou des supports d'information.
\item Citez trois mots du texte qui appartiennent au champ lexical de la paix.
\item Quel verbe du texte signifie « donner la parole à » ?
\end{enumerate}

\courssub{Le champ lexical des médias}

\begin{definition}[Les médias et leurs acteurs]
\begin{itemize}
\item \esp{la prensa} : la presse (ensemble des journaux) ; \esp{la prensa escrita} : la presse écrite.
\item \esp{el periódico} : le journal ; \esp{el diario} : le quotidien.
\item \esp{el/la periodista} : le/la journaliste ; \esp{el redactor / la redactora} : le/la rédacteur(-trice).
\item \esp{la noticia} : la nouvelle, l'information ; \esp{las noticias} : les informations (journal télévisé).
\item \esp{el informativo} : le journal télévisé ou radiophonique.
\item \esp{el titular} : le gros titre, le titre principal.
\item \esp{la entrevista} : l'interview ; \esp{entrevistar} : interviewer.
\item \esp{el reportaje} : le reportage ; \esp{el artículo} : l'article ; \esp{el editorial} : l'éditorial.
\end{itemize}
\end{definition}

\begin{definition}[Les supports et les canaux de diffusion]
\begin{itemize}
\item \esp{la televisión} : la télévision ; \esp{la radio} : la radio ; \esp{Internet} : Internet.
\item \esp{las redes sociales} : les réseaux sociaux.
\item \esp{el canal} : la chaîne (de télévision) ; \esp{el programa} : l'émission.
\item \esp{el medio de comunicación} : le moyen de communication ; au pluriel, \esp{los medios} désigne « les médias ».
\end{itemize}
\end{definition}

\begin{definition}[Informer, vérifier, tromper]
\begin{itemize}
\item \esp{informar} : informer ; \esp{verificar} : vérifier ; \esp{publicar} : publier ; \esp{difundir} : diffuser.
\item \esp{la fuente} : la source ; \esp{la verdad} : la vérité ; \esp{la mentira} : le mensonge.
\item \esp{la desinformación} : la désinformation ; \esp{las noticias falsas} : les fausses nouvelles ; \esp{la propaganda} : la propagande ; \esp{el rumor} : la rumeur.
\end{itemize}
\end{definition}

\begin{center}
\begin{tabularx}{\linewidth}{|X|X|X|}
\hline
\rowcolor{popblueL} Español & Français & Remarque \\
\hline
la prensa & la presse & nom féminin collectif \\
\hline
el periódico / el diario & le journal / le quotidien & \esp{diario} = aussi « quotidien » (adjectif) \\
\hline
el/la periodista & le/la journaliste & même forme aux deux genres \\
\hline
la noticia & la nouvelle, l'information & \textbf{pas « une notice » !} \\
\hline
el informativo & le journal (TV/radio) & nom masculin, pas l'adjectif « informatif » \\
\hline
el titular & le gros titre & aussi « le titulaire » (d'un poste) \\
\hline
la fuente & la source & aussi « la fontaine » \\
\hline
el reportaje & le reportage & \esp{hacer un reportaje} = faire un reportage \\
\hline
la entrevista & l'interview & \esp{entrevistar a alguien} \\
\hline
el canal / el programa & la chaîne / l'émission & attention au genre : \esp{el programa} \\
\hline
\end{tabularx}
\end{center}

\begin{exemplebox}[Exemples traduits]
\begin{itemize}
\item \esp{Los periodistas verifican las fuentes antes de publicar.} « Les journalistes vérifient les sources avant de publier. »
\item \esp{El informativo de la noche ha dedicado un reportaje a la paz.} « Le journal du soir a consacré un reportage à la paix. »
\item \esp{En las redes sociales circulan muchas noticias falsas.} « Sur les réseaux sociaux circulent beaucoup de fausses nouvelles. »
\item \esp{El titular del periódico anuncia el fin del conflicto.} « Le gros titre du journal annonce la fin du conflit. »
\end{itemize}
\end{exemplebox}

\courssub{Le champ lexical de la paix}

\begin{definition}[Paix, conflit et réconciliation]
\begin{itemize}
\item \esp{la paz} : la paix ; \esp{el conflicto} : le conflit ; \esp{la guerra} : la guerre ; \esp{la violencia} : la violence.
\item \esp{el diálogo} : le dialogue ; \esp{la tolerancia} : la tolérance ; \esp{el respeto} : le respect.
\item \esp{la reconciliación} : la réconciliation ; \esp{la convivencia} : la vie commune, le vivre-ensemble.
\item \esp{la justicia} : la justice ; \esp{el perdón} : le pardon ; \esp{la escucha} : l'écoute.
\end{itemize}
\end{definition}

\begin{propriete}[Verbes d'action pour la paix]
Ces verbes sont tous \cle{réguliers} (sauf \esp{promover}) et très utiles pour l'expression écrite :
\begin{itemize}
\item \esp{dialogar} : dialoguer ; \esp{escuchar} : écouter ; \esp{perdonar} : pardonner.
\item \esp{respetar} : respecter ; \esp{unir} : unir ; \esp{promover} (radical \esp{promuevo} à la 1ère pers. sg. présent indicatif et subjonctif) : promouvoir.
\item \esp{rechazar la violencia} : rejeter la violence ; \esp{construir la paz} : construire la paix ; \esp{superar el conflicto} : surmonter le conflit.
\end{itemize}
\end{propriete}

\begin{center}
\begin{tabularx}{\linewidth}{|X|X|X|}
\hline
\rowcolor{popgreenL} Español & Français & Remarque \\
\hline
la paz & la paix & féminin malgré le -z : \esp{la paz}, jamais « el paz » \\
\hline
el conflicto & le conflit & masculin régulier \\
\hline
la violencia & la violence & féminin en -cia \\
\hline
el diálogo & le dialogue & accent sur le a : \esp{diálogo} \\
\hline
la tolerancia & la tolérance & famille : \esp{tolerante} = tolérant \\
\hline
la reconciliación & la réconciliation & verbe pronominal : \esp{reconciliarse} = se réconcilier \\
\hline
la convivencia & la vie en commun & verbe : \esp{convivir} = vivre ensemble \\
\hline
el respeto & le respect & verbe : \esp{respetar} \\
\hline
la justicia & la justice & adjectif : \esp{justo/a} = juste \\
\hline
\end{tabularx}
\end{center}

\begin{exemplebox}[Expressions utiles traduites]
\begin{itemize}
\item \esp{Los medios deben promover el diálogo.} « Les médias doivent promouvoir le dialogue. »
\item \esp{Rechazamos la violencia.} « Nous rejetons la violence. »
\item \esp{Hay que dar voz a las víctimas.} « Il faut donner la parole aux victimes. »
\item \esp{La propaganda siembra el odio entre las comunidades.} « La propagande sème la haine entre les communautés. »
\item \esp{Construir la paz empieza por escuchar al otro.} « Construire la paix commence par écouter l'autre. »
\end{itemize}
\end{exemplebox}

\courssub{Carte mentale : deux pôles reliés par la vérité}

\begin{popfigure}
\begin{center}
\begin{tikzpicture}[
  polo/.style={rectangle, rounded corners, draw=popdark, thick, fill=popblueL, font=\bfseries, align=center, minimum width=2.6cm, minimum height=0.9cm},
  polo2/.style={rectangle, rounded corners, draw=popdark, thick, fill=popgreenL, font=\bfseries, align=center, minimum width=2.6cm, minimum height=0.9cm},
  hoja/.style={rectangle, rounded corners, draw=popline, fill=popcream, align=center, font=\small, minimum width=2.2cm, minimum height=0.7cm},
  puente/.style={rectangle, rounded corners, draw=poporange, thick, fill=poporangeL, align=center, font=\small\bfseries, minimum width=3.2cm, minimum height=0.8cm}
]
\node[polo] (medios) at (-4.2,0) {MEDIOS};
\node[polo2] (paz) at (4.2,0) {PAZ};
\node[puente, align=center] (verdad) at (0,0){informar\\con verdad};
% Branches MEDIOS
\node[hoja] (prensa) at (-6.6,1.6) {la prensa};
\node[hoja] (radio) at (-6.6,0.55) {la radio};
\node[hoja] (tv) at (-6.6,-0.55) {la televisión};
\node[hoja] (redes) at (-6.6,-1.6) {redes sociales};
\node[hoja] (periodista) at (-2.4,1.9) {el/la periodista};
\node[hoja] (noticia) at (-2.4,-1.9) {la noticia};
% Branches PAZ
\node[hoja] (dialogo) at (6.6,1.6) {el diálogo};
\node[hoja] (tolerancia) at (6.6,0.55) {la tolerancia};
\node[hoja] (reconciliacion) at (6.6,-0.55) {reconciliación};
\node[hoja] (conflicto) at (6.6,-1.6) {el conflicto};
\node[hoja] (convivencia) at (2.4,1.9) {la convivencia};
\node[hoja] (violencia) at (2.4,-1.9) {la violencia};
% Connections
\draw[thick, popblue] (medios) -- (verdad);
\draw[thick, popgreen] (verdad) -- (paz);
\draw[popline] (medios) -- (prensa);
\draw[popline] (medios) -- (radio);
\draw[popline] (medios) -- (tv);
\draw[popline] (medios) -- (redes);
\draw[popline] (medios) -- (periodista);
\draw[popline] (medios) -- (noticia);
\draw[popline] (paz) -- (dialogo);
\draw[popline] (paz) -- (tolerancia);
\draw[popline] (paz) -- (reconciliacion);
\draw[popline] (paz) -- (conflicto);
\draw[popline] (paz) -- (convivencia);
\draw[popline] (paz) -- (violencia);
\end{tikzpicture}
\end{center}
\legende{Carte mentale lexicale : les médias et la paix, reliés par l'information véridique.}
\end{popfigure}

\courssub{Genre et articles : les pièges grammaticaux du lexique}

\begin{propriete}[Le genre des noms de ce chapitre]
\begin{enumerate}
\item \esp{el periodista / la periodista} : les noms en \cle{-ista} ont une seule forme ; seul l'article change selon le sexe de la personne. Même chose pour \esp{el/la joven}, \esp{el/la pacifista}, \esp{el/la testigo} (le/la témoin).
\item \esp{la paz} est \cle{féminin} malgré sa terminaison en -z : \esp{La paz es necesaria.} « La paix est nécessaire. »
\item \esp{el programa}, \esp{el problema}, \esp{el tema}, \esp{el sistema}, \esp{el clima} : masculins malgré le -a final (mots d'origine grecque en -ma).
\item \esp{la radio} : féminin (abréviation de \esp{la radiodifusión}), mais on entend aussi \esp{el radio} en Amérique pour désigner l'appareil récepteur.
\end{enumerate}
\end{propriete}

\begin{attention}[Faux amis et erreurs fréquentes des francophones]
\begin{itemize}
\item \esp{una noticia} = une \textbf{nouvelle} (une information), jamais « une notice ». « Une notice » (mode d'emploi) se dit \esp{las instrucciones} ou \esp{el manual}.
\item \esp{el informativo} = le \textbf{journal télévisé/radiophonique} (nom masculin) ; ce n'est pas l'adjectif français « informatif » (qui se dit \esp{informativo} en espagnol, mais s'emploie différemment : \esp{un programa informativo}).
\item On dit \esp{la paz} et non « el paz » : la terminaison -z ne rend pas le mot masculin (cf. \esp{la luz}, \esp{la vez}, \esp{la nariz}).
\item \esp{la entrevista} = l'interview ; le verbe est \esp{entrevistar}, pas « interviewar » (calque de l'anglais).
\item \esp{el conflicto} s'écrit avec un seul « f » et un seul « l » (comme en français « conflit » mais avec -o final). Attention à la prononciation : [kon.flik.to].
\item \esp{informar} se construit avec \esp{sobre} ou \esp{de} : \esp{informar sobre el conflicto} / \esp{informar de lo ocurrido} = informer sur le conflit / informer de ce qui s'est passé.
\item \esp{empezar por} + infinitif = « commencer par » (ex. : \esp{empezar por escuchar}). Ne pas confondre avec \esp{empezar a} + infinitif = « commencer à ».
\end{itemize}
\end{attention}

\begin{savaistu}[Fake news et journalistes en danger]
L'espagnol possède l'expression \esp{noticias falsas}, mais le mot anglais \emph{fake news} est très répandu dans la presse hispanophone, y compris en Guinée équatoriale et dans les médias espagnols. L'UNESCO, qui défend la liberté de la presse dans le monde entier, a fait du 2 novembre la Journée internationale de la fin de l'impunité pour les crimes contre les journalistes : une manière de rappeler qu'informer avec vérité est un métier parfois dangereux, mais essentiel pour la paix. Au Cameroun comme ailleurs, les journalistes jouent un rôle clé dans la prévention des conflits (bilingue, zones anglophones, élections).
\end{savaistu}

\courssub{Exercice résolu et entraînement}

\begin{exoresolu}[Compléter avec le mot juste]
Énoncé — Complétez les phrases avec le mot espagnol qui convient : \esp{la fuente}, \esp{el titular}, \esp{la paz}, \esp{el informativo}, \esp{la reconciliación}.
\begin{enumerate}
\item Antes de publicar, el periodista debe verificar \_\_\_\_\_\_\_\_.
\item \_\_\_\_\_\_\_\_ del periódico anuncia el fin de la guerra.
\item Después del conflicto, las dos comunidades celebraron \_\_\_\_\_\_\_\_.
\item Vi las noticias en \_\_\_\_\_\_\_\_ de las ocho.
\item Construir \_\_\_\_\_\_\_\_ es tarea de todos.
\end{enumerate}
\tcblower
Solution détaillée :
\begin{enumerate}
\item \esp{Antes de publicar, el periodista debe verificar la fuente.} « Avant de publier, le journaliste doit vérifier la source. » (On vérifie une \emph{source}, pas un titre.)
\item \esp{El titular del periódico anuncia el fin de la guerra.} « Le gros titre du journal annonce la fin de la guerre. » (C'est le titre qui « annonce ».)
\item \esp{Después del conflicto, las dos comunidades celebraron la reconciliación.} « Après le conflit, les deux communautés ont célébré la réconciliation. »
\item \esp{Vi las noticias en el informativo de las ocho.} « J'ai vu les informations au journal de huit heures. » (Masculin : \esp{el informativo}.)
\item \esp{Construir la paz es tarea de todos.} « Construire la paix est l'affaire de tous. » (Féminin : \esp{la paz}.)
\end{enumerate}
\end{exoresolu}

\begin{exercice}[À vous de jouer]
\begin{enumerate}
\item Traduisez en espagnol : « Les médias doivent donner la parole aux jeunes et rejeter la violence. »
\item Traduisez en français : \esp{La propaganda y las noticias falsas siembran el odio y destruyen la convivencia.}
\item Trouvez l'intrus et justifiez : \esp{el reportaje} / \esp{la entrevista} / \esp{la tolerancia} / \esp{el titular}.
\item Complétez avec le bon article (el/la) : \_\_\_ periodista, \_\_\_ paz, \_\_\_ programa, \_\_\_ fuente, \_\_\_ conflicto.
\end{enumerate}
\end{exercice}

\begin{corrige}[Exercice « À vous de jouer »]
\begin{enumerate}
\item \esp{Los medios deben dar voz a los jóvenes y rechazar la violencia.}
\item « La propagande et les fausses nouvelles sèment la haine et détruisent le vivre-ensemble. »
\item L'intrus est \esp{la tolerancia} : c'est le seul mot qui appartient au champ lexical de la paix ; les trois autres appartiennent au champ lexical des médias.
\item \esp{la periodista} (ou \esp{el periodista}), \esp{la paz}, \esp{el programa}, \esp{la fuente}, \esp{el conflicto}.
\end{enumerate}
\end{corrige}

\begin{aretenir}[L'essentiel du lexique]
Deux champs lexicaux à maîtriser pour le commentaire et la traduction : les \cle{médias} (\esp{la prensa, el periodista, la noticia, el informativo, el titular, la fuente, el reportaje}) et la \cle{paix} (\esp{la paz, el conflicto, el diálogo, la tolerancia, la reconciliación, la convivencia}). Le pont entre les deux : \esp{informar con verdad} et les verbes d'action (\esp{promover el diálogo, rechazar la violencia, dar voz a, construir la paz}). Trois pièges à retenir : \esp{la noticia} n'est pas « la notice », \esp{el informativo} n'est pas « l'informatif », et on dit \esp{la paz}, jamais « el paz ».
\end{aretenir}

\coursec{Textes argumentatifs et injonctifs : repérage et grammaire de l'appel}

Dans la section précédente, nous avons constitué le lexique des médias et de la culture de la paix : \esp{prensa}, \esp{periodista}, \esp{noticia}, \esp{conflicto}, \esp{diálogo}... Ces mots ne vivent pas isolément : ils circulent dans des textes qui veulent soit \cle{convaincre} (textes argumentatifs : éditoriaux, chroniques, essais), soit \cle{faire agir} (textes injonctifs : appels, slogans, campagnes de sensibilisation, affiches). Au baccalauréat, le commentaire de texte exige de reconnaître ces deux logiques et de nommer précisément les procédés linguistiques qui les construisent. C'est l'objet de cette section.

\courssub{Observer d'abord : deux textes, deux intentions}

\begin{textebox}[Deux voix sur les médias]
\textbf{Texto A (editorial).} Los medios de comunicación no son neutros: creo que pueden ser armas de guerra o instrumentos de paz, según el uso que hagamos de ellos. En efecto, cuando un periodista difunde rumores sin verificar las fuentes, alimenta el odio entre las comunidades. Sin embargo, la radio comunitaria también puede reunir a vecinos enfrentados y dar voz a las víctimas. Además, las redes sociales permiten que los jóvenes de Yaoundé, Madrid o Malabo dialoguen entre sí. Por eso, considero que la responsabilidad de los medios es inmensa: informar es ya construir la paz o preparar el conflicto.

\textbf{Texto B (campaña).} ¡Jóvenes de Camerún, despertad! No difundáis noticias falsas: verificad siempre las fuentes antes de compartir. Hay que usar las redes sociales para unir, no para dividir. Usemos la palabra contra la violencia. Escuchemos a las víctimas, respetemos al otro, dialoguemos. ¡Que se escuche la voz de la paz en cada barrio! Ojalá reine el diálogo entre los pueblos. No matarás la verdad. El futuro de nuestro país está en vuestras manos: ¡actuad hoy!
\end{textebox}

\textbf{Glossaire :} \esp{según} = selon ; \esp{las fuentes} = les sources ; \esp{alimentar el odio} = nourrir la haine ; \esp{reunir} = rassembler ; \esp{dar voz a} = donner la parole à ; \esp{compartir} = partager ; \esp{despertad} = réveillez-vous ; \esp{reinar} = régner.

\textbf{Questions d'observation :} 1) Quel texte cherche à convaincre, et lequel cherche à faire agir ? 2) Souligne dans le texte A les verbes d'opinion et les connecteurs. 3) Souligne dans le texte B toutes les formes verbales qui donnent un ordre ou lancent un appel. Que remarques-tu sur leur mode ?

\courssub{Le texte argumentatif : structure et marques linguistiques}

\begin{definition}[Le texte argumentatif]
Un texte argumentatif défend une \cle{thèse} (une opinion) à l'aide d'\cle{arguments} étayés par des \cle{exemples}. Sa structure canonique : \textbf{introduction} (sujet + problématique), \textbf{développement} (thèse, arguments, exemples), \textbf{conclusion} (bilan, ouverture).
\end{definition}

Dans le texte A, la thèse est la phrase : \esp{Los medios pueden ser armas de guerra o instrumentos de paz}. Tout le reste du texte la justifie. Pour la repérer et la défendre, l'espagnol dispose de deux familles d'outils.

\textbf{1. Les verbes d'opinion.} \esp{creer que}, \esp{pensar que}, \esp{considerar que}, \esp{opinar que}, \esp{estar convencido de que} sont suivis de l'\cle{indicatif} quand on affirme : \esp{Creo que los medios influyen en la sociedad.} « Je crois que les médias influencent la société. » À la forme négative, l'espagnol passe au subjonctif : \esp{No creo que los medios sean neutros.} « Je ne crois pas que les médias soient neutres. »

\textbf{2. Les connecteurs logiques.} Ils articulent la démonstration :

\begin{center}
\begin{tabularx}{\linewidth}{|X|X|X|}\hline
\rowcolor{popblueL} Fonction & Español & Français \\ \hline
Cause & porque, ya que, puesto que & parce que, puisque \\ \hline
Addition & además, también, asimismo & de plus, également \\ \hline
Opposition & sin embargo, no obstante, pero & cependant, mais \\ \hline
Conséquence & por eso, por lo tanto, así que & c'est pourquoi, donc \\ \hline
Conclusion & en conclusión, en resumen & en conclusion, en résumé \\ \hline
\end{tabularx}
\end{center}

\begin{methode}[Repérer la thèse d'un texte argumentatif]
\begin{enumerate}
\item Lis le texte entier une première fois sans t'arrêter.
\item Demande-toi : quelle est la phrase que tout le reste du texte cherche à justifier ? C'est la thèse.
\item Vérifie : elle contient souvent un verbe d'opinion (\esp{creo que}, \esp{considero que}) ou elle suit un connecteur conclusif (\esp{por eso}, \esp{en conclusión}).
\item Classe ensuite chaque phrase restante : argument (raison générale) ou exemple (cas concret).
\end{enumerate}
\end{methode}

\courssub{Le texte injonctif : donner des ordres, lancer des appels}

\begin{definition}[Le texte injonctif]
Un texte injonctif cherche à \cle{faire agir} le lecteur : ordres, conseils, appels, slogans, consignes. Ses outils principaux en espagnol : l'\cle{impératif}, les tournures impersonnelles \esp{hay que + infinitivo} et \esp{es necesario + infinitivo}, l'appel avec \esp{que / ojalá + subjuntivo}, et le futur d'ordre.
\end{definition}

\textbf{1. L'impératif affirmatif.} Rappel des formes régulières (verbe \esp{dialogar}) : \esp{dialoga} (tú), \esp{dialogue} (usted), \esp{dialoguemos} (nosotros), \esp{dialogad} (vosotros), \esp{dialoguen} (ustedes). Exemples : \esp{¡Escucha la radio comunitaria!} « Écoute la radio communautaire ! » ; \esp{¡Verificad las fuentes!} « Vérifiez les sources ! »

\begin{propriete}[L'impératif négatif = no + subjonctif présent]
En espagnol, l'impératif \textbf{négatif} n'existe pas comme forme propre : on emploie \esp{no} + \textbf{subjonctif présent}. \esp{¡No difundas noticias falsas!} « Ne diffuse pas de fausses nouvelles ! » ; \esp{¡No maten!} « Ne tuez pas ! » ; \esp{¡No compartamos rumores!} « Ne partageons pas de rumeurs ! »
\end{propriete}

\begin{attention}[Piège majeur pour les francophones]
Ne dis jamais \esp{¡No difundes!} (indicatif) pour « ne diffuse pas ! ». Après \esp{no} à l'impératif, le subjonctif est obligatoire : \esp{¡No difundas!}, \esp{¡No hables!}, \esp{¡No escriban!}. L'indicatif après \esp{no} décrit un fait, il ne donne pas d'ordre.
\end{attention}

\textbf{2. L'exhortation à la première personne du pluriel.} Pour s'associer à l'action (« agissons ensemble »), l'espagnol utilise la forme de \esp{nosotros} du subjonctif présent, identique à l'impératif : \esp{¡Dialoguemos!} « Dialoguons ! » ; \esp{¡Usemos los medios para unir!} « Utilisons les médias pour unir ! » ; \esp{¡Escuchemos a las víctimas!} « Écoutons les victimes ! » C'est la forme reine des appels à la paix.

\textbf{3. Les tournures impersonnelles.} \esp{Hay que + infinitivo} = « il faut » ; \esp{Es necesario + infinitivo} = « il est nécessaire de ». \esp{Hay que verificar las fuentes.} « Il faut vérifier les sources. » ; \esp{Es necesario proteger la libertad de prensa.} « Il est nécessaire de protéger la liberté de la presse. »

\begin{exemplebox}[L'appel avec que et ojalá]
\esp{¡Que viva la paz!} « Vive la paix ! » ; \esp{¡Que se escuche la voz de las víctimas!} « Que la voix des victimes soit entendue ! » ; \esp{Ojalá reine el diálogo entre los pueblos.} « Pourvu que le dialogue règne entre les peuples. » ; \esp{Ojalá haya paz en el mundo.} « Pourvu qu'il y ait la paix dans le monde. »
\end{exemplebox}

\begin{propriete}[Ojalá (+ que) + subjonctif]
\esp{Ojalá} (mot d'origine arabe, « qu'Allah veuille ») exprime un \cle{souhait} et est \textbf{toujours} suivi du subjonctif, jamais de l'indicatif. Le \esp{que} est facultatif : \esp{Ojalá (que) haya paz.} De même, \esp{que + subjonctif} en début de phrase exprime un vœu ou un ordre rapporté : \esp{¡Que entre el periodista!} « Que le journaliste entre ! »
\end{propriete}

\textbf{4. Le futur d'ordre (complément utile au Bac).} Dans un registre solennel (lois, commandements, serments), l'espagnol emploie le futur simple avec valeur d'impératif : \esp{No matarás.} « Tu ne tueras pas. » ; \esp{No robarás.} « Tu ne voleras pas. » ; \esp{Respetarás la verdad.} « Tu respecteras la vérité. » Le français connaît le même procédé.

\textbf{5. Comparaison français / espagnol dans les consignes.} Le français utilise volontiers l'infinitif sur les panneaux (« Ne pas fumer », « Ne pas toucher »). L'espagnol préfère \esp{no + subjonctif} (\esp{No fume}, \esp{No toque}) ou le participe \esp{prohibido / se prohíbe + infinitivo} : \esp{Prohibido fumar} « Interdiction de fumer » ; \esp{Se prohíbe difundir noticias falsas} « Il est interdit de diffuser de fausses nouvelles ».

\courssub{Synthèse : convaincre ou faire agir}

\begin{popfigure}
\begin{tikzpicture}[node distance=0.6cm]
\node[draw=popblue, fill=popblueL, rounded corners, text width=5.6cm, align=center] (arg) {\textbf{Argumentatif : convaincre}\\[2pt] \esp{creo que, considero que}\\ \esp{porque, ya que}\\ \esp{sin embargo, además}\\ \esp{por eso, en conclusión}};
\node[draw=poporange, fill=poporangeL, rounded corners, text width=5.6cm, align=center, right=1.2cm of arg] (inj) {\textbf{Injonctif : faire agir}\\[2pt] impératif : \esp{¡Dialoga! ¡No difundas!}\\ \esp{hay que + infinitivo}\\ \esp{¡Dialoguemos!}\\ \esp{ojalá / que + subjuntivo}};
\node[draw=popgreen, fill=popgreenL, rounded corners, text width=6cm, align=center, below=0.9cm of arg, xshift=3.4cm] (obj) {\textbf{Objectif commun}\\ informer, unir, construire la paz};
\draw[-{Stealth}, thick, popblue] (arg) -- (obj);
\draw[-{Stealth}, thick, poporange] (inj) -- (obj);
\end{tikzpicture}
\legende{Les deux logiques du texte sur les médias : convaincre (argumentatif) et faire agir (injonctif).}
\end{popfigure}

\begin{center}
\begin{tabularx}{\linewidth}{|X|X|X|}\hline
\rowcolor{popblueL} Français & Español & Remarque \\ \hline
Ne diffuse pas ! & ¡No difundas! & no + subjonctif \\ \hline
Dialoguons ! & ¡Dialoguemos! & subjonctif nosotros \\ \hline
Il faut vérifier les sources. & Hay que verificar las fuentes. & tournure impersonnelle \\ \hline
Pourvu qu'il y ait la paix. & Ojalá haya paz. & jamais d'indicatif après ojalá \\ \hline
Tu ne tueras pas. & No matarás. & futur d'ordre solennel \\ \hline
Ne pas fumer (panneau) & No fume / Prohibido fumar & infinitif français, subjonctif ou prohibido en espagnol \\ \hline
\end{tabularx}
\end{center}

\begin{exoresolu}[Repérage et transformation]
\textbf{Énoncé.} 1) Dans la phrase \esp{« Por eso, considero que la radio comunitaria es un instrumento de paz »}, identifie le connecteur, le verbe d'opinion et la thèse. 2) Transforme en impératif négatif : \esp{Difundes rumores.} 3) Transforme en exhortation collective : \esp{Usamos los medios para unir.}
\tcblower
\textbf{Solution.} 1) Le connecteur de conséquence est \esp{por eso} (« c'est pourquoi ») ; le verbe d'opinion est \esp{considero que} (suivi de l'indicatif \esp{es}, car on affirme) ; la thèse est : \esp{la radio comunitaria es un instrumento de paz}. 2) L'impératif négatif exige \esp{no} + subjonctif présent de \esp{difundir} à la 2\up{e} personne : \esp{¡No difundas rumores!} (« Ne diffuse pas de rumeurs ! ») — et non \esp{no difundes}, qui serait un indicatif. 3) L'exhortation à \esp{nosotros} prend la forme du subjonctif : \esp{¡Usemos los medios para unir!} (« Utilisons les médias pour unir ! »).
\end{exoresolu}

\begin{exercice}[À toi de jouer]
1) Souligne les marques injonctives dans ce slogan et traduis-le : \esp{« No creas todo lo que lees: verifica, compara y piensa. Hay que educar para la paz. ¡Que el diálogo venza al odio! »} 2) Mets à la forme demandée : \esp{escuchar} (impératif tú affirmatif) ; \esp{respetar} (impératif vosotros négatif) ; \esp{construir} (exhortation nosotros) ; \esp{haber} (avec ojalá). 3) Traduis en espagnol : « Ne partagez pas de fausses nouvelles ! Il faut protéger les journalistes. Pourvu que la paix règne au Cameroun. »
\end{exercice}

\begin{corrige}[Exercice « À toi de jouer »]
1) Marques injonctives : \esp{No creas} (impératif négatif), \esp{verifica, compara y piensa} (impératifs affirmatifs), \esp{Hay que educar} (tournure impersonnelle), \esp{Que el diálogo venza} (que + subjonctif). Traduction : « Ne crois pas tout ce que tu lis : vérifie, compare et réfléchis. Il faut éduquer pour la paix. Que le dialogue vainque la haine ! » 2) \esp{¡Escucha!} ; \esp{¡No respetéis!} ; \esp{¡Construyamos!} ; \esp{Ojalá haya...} 3) \esp{¡No compartáis noticias falsas! Hay que proteger a los periodistas. Ojalá reine la paz en Camerún.} (Attention : \esp{proteger a los periodistas} avec la préposition \esp{a} personnelle devant les personnes.)
\end{corrige}

\begin{aretenir}[L'essentiel de la section]
Un texte argumentatif \textbf{convainc} : thèse + arguments + exemples, verbes d'opinion (\esp{creer que} + indicatif) et connecteurs (\esp{porque}, \esp{sin embargo}, \esp{por eso}). Un texte injonctif \textbf{fait agir} : impératif affirmatif, impératif négatif en \esp{no + subjuntivo} (\esp{¡No difundas!}), exhortation collective (\esp{¡Dialoguemos!}), \esp{hay que + infinitivo}, appels en \esp{que / ojalá + subjuntivo}, et futur d'ordre solennel (\esp{No matarás}). Ces marques sont tes meilleurs indices pour le commentaire de texte au baccalauréat.
\end{aretenir}

\coursec{Méthode du commentaire de texte et commentaire modèle}

Après avoir appris à \cle{repérer la thèse, les arguments et les procédés injonctifs} dans la section précédente, nous abordons maintenant la rédaction structurée du commentaire de texte. C'est l'exercice roi de l'épreuve écrite du baccalauréat : il ne s'agit pas de résumer, mais d'\textbf{analyser comment le texte construit son sens et convainc son lecteur}. La méthode en quatre étapes ci-dessous, appliquée avec rigueur, garantit une copie claire, analytique et bien notée.

\courssub{Étape 1 — Lecture active et repérage (15--20 minutes)}

\begin{methode}[Lecture en deux temps]
\begin{enumerate}
    \item \textbf{Première lecture (globale) :} Lisez le texte d'une traite sans stylo. Identifiez : la nature du texte (article d'opinion, éditorial, discours, lettre ouverte), le thème général (médias, paix, conflits), le ton (engagé, ironique, solennel, alarmiste) et la thèse principale (pour ou contre ? quelle solution ?).
    \item \textbf{Deuxième lecture (analytique, avec surlignage) :} Munissez-vous de deux couleurs.
    \begin{itemize}
        \item \textbf{Couleur 1 (Contenu) :} Surlignez la \cle{thèse} (souvent au début ou à la fin), les \cle{arguments} (idées forces), les \cle{exemples/illustrations} (faits, chiffres, anecdotes, références).
        \item \textbf{Couleur 2 (Forme/Procédés) :} Surlignez les \cle{connecteurs logiques} (\esp{por tanto, sin embargo, además, en conclusión}), les \cle{procédés argumentatifs} (questions rhétoriques, accumulations, oppositions, comparaisons), les \cle{modes verbaux} (indicatif = affirmation, subjonctif = subjectivité/doute/volonté, \cle{impératif} = injonction), le \cle{lexique connoté} (mots forts : \esp{guerra, mentira, verdad, urgente, responsabilidad}), la \cle{ponctuation expressive} (\esp{¡ ! ¿ ?}) et les \cle{figures de style} (anaphore, métaphore, antithèse).
    \end{itemize}
    \item \textbf{Prise de notes schématique :} Sur brouillon, construisez le squelette : Thèse $\rightarrow$ Arg 1 (ex) $\rightarrow$ Arg 2 (ex) $\rightarrow$ Arg 3 (ex) / Procédés dominants.
\end{enumerate}
\end{methode}

\begin{attention}[Piège fréquent : la lecture passive]
Ne lisez pas « pour savoir de quoi ça parle ». Lisez « pour voir comment ça marche ». Un élève qui ne fait que la première lecture écrira un \textbf{résumé} (sanctionné 0/4 ou 1/4 en analyse). L'examinateur attend une \textbf{analyse des procédés} : \emph{« L'auteur utilise l'antithèse "armas de guerra / instrumentos de paz" pour forcer le choix du lecteur »}, et non \emph{« L'auteur dit que les médias peuvent être des armes ou des instruments »}.
\end{attention}

\courssub{Étape 2 — Rédaction de l'introduction (environ 10 lignes)}

L'introduction suit un protocole immuable en \textbf{quatre mouvements} (entono $\rightarrow$ présentation $\rightarrow$ problématique $\rightarrow$ plan).

\begin{methode}[Les 4 mouvements de l'introduction]
\begin{enumerate}
    \item \textbf{L'amorce (contexto) :} Une phrase large sur le thème (médias, paix, désinformation) sans citer le texte. \esp{En la era digital, la información circula a una velocidad vertiginosa...}
    \item \textbf{La présentation (fiche d'identité) :} Auteur (si connu) / Source / Date / Genre / Titre. Thème précis. \esp{El artículo de opinión « Los medios de comunicación, ¿armas de guerra o instrumentos de paz? », publicado en [Fuente/Fecha], aborda el papel crucial de la prensa en la construcción de la paz.}
    \item \textbf{La problématique (pregunta directriz) :} La question à laquelle le texte répond. Elle naît de la tension du texte. \esp{¿Hasta qué punto los medios de comunicación pueden convertirse en verdaderos artífices de la paz en un contexto de posverdad y conflictos?}
    \item \textbf{L'annonce du plan (estructura) :} \esp{Analizaremos primero la tesis defendida y sus argumentos (I), para estudiar luego los procedimientos expresivos que refuerzan la persuasión y el llamamiento a la acción (II).}
\end{enumerate}
\end{methode}

\begin{propriete}[Phrases-outils indispensables pour l'introduction et le développement]
\begin{center}
\begin{tabularx}{\linewidth}{|X|X|X|}
\hline
\rowcolor{popblueL} \textbf{Fonction} & \textbf{En español (à mémoriser)} & \textbf{En français} \\
\hline
Présenter le thème & \esp{El texto trata de / versa sobre / aborda el tema de...} & Le texte traite de / porte sur / aborde le thème de... \\
\hline
Thèse & \esp{El autor defiende la idea de que / sostiene que / afirma que...} & L'auteur défend l'idée que / soutient que / affirme que... \\
\hline
Argument & \esp{Para fundamentar su tesis, argumenta que / aduce que / señala que...} & Pour étayer sa thèse, il argue que / allègue que / signale que... \\
\hline
Exemple & \esp{Ilustra esto con el ejemplo de / cita el caso de / pone de relieve...} & Il illustre cela par l'exemple de / cite le cas de / met en relief... \\
\hline
Connecteur (addition) & \esp{Además / Asimismo / En primer lugar, en segundo lugar...} & De plus / Également / En premier lieu, en second lieu... \\
\hline
Connecteur (opposition) & \esp{Sin embargo / No obstante / Por el contrario / Ahora bien...} & Cependant / Néanmoins / Au contraire / Or... \\
\hline
Connecteur (conséquence) & \esp{Por tanto / Por consiguiente / En consecuencia / De este modo...} & Par conséquent / Donc / De ce fait / Ainsi... \\
\hline
Analyse procédés & \esp{Utiliza / Recurre a / Emplea / Destaca el uso de... (imperativos, preguntas retóricas, léxico bélico...)} & Il utilise / Recourt à / Emploie / Met en relief l'usage de... (impératifs, questions rhétoriques, lexique guerrier...) \\
\hline
Effet sur lecteur & \esp{Busca concienciar / movilizar / interpelar / convencer al lector de que...} & Il cherche à sensibiliser / mobiliser / interpeller / convaincre le lecteur que... \\
\hline
Citation & \esp{Como se lee en el verso/frase « ... » / Cita textualmente: « ... »} & Comme on lit dans le vers/la phrase « ... » / Cite textuellement : « ... » \\
\hline
Conclusion & \esp{En definitiva / En conclusión / Para terminar, el texto nos invita a reflexionar sobre...} & En définitive / En conclusion / Pour finir, le texte nous invite à réfléchir sur... \\
\hline
Ouverture & \esp{Esto cobra especial relevancia en contextos como... (Camerún, redes sociales hoy)} & Cela prend une importance particulière dans des contextes comme... (Cameroun, réseaux sociaux aujourd'hui) \\
\hline
\end{tabularx}
\end{center}
\end{propriete}

\courssub{Étape 3 — Rédaction du développement (2 ou 3 parties, ~20--25 lignes)}

Le développement suit le plan annoncé. Deux structures types sont acceptées au Bac :

\begin{definition}[Structure type A : Thématique / Formelle (la plus sûre)]
\begin{itemize}
    \item \textbf{Partie I : La thèse et les arguments (Le \emph{quoi}).} Regroupe le contenu. Une idée directrice par paragraphe (Argument 1, Argument 2...). Citez court (3--6 mots), commentez l'argument.
    \item \textbf{Partie II : Les procédés d'expression (Le \emph{comment}).} Analyse la forme au service du fond : lexique (champs lexicaux), syntaxe (phrases courtes/longues, interrogatives, impératives), connecteurs, figures de style, tons, registres de langue. Montrez \emph{comment} ces procédés renforcent la thèse.
    \item \textbf{Partie III (Facultatif, pour l'excellence) : Appréciation critique / Portée.} Votre avis argumenté : le texte est-il convaincant ? Limites ? Actualité ? Lien avec le contexte camerounais.
\end{itemize}
\end{definition}

\begin{definition}[Structure type B : Linéaire (par paragraphes du texte)]
Seulement si le texte a une structure chronologique ou narrative très marquée. Moins recommandée pour un texte argumentatif car elle mélange souvent fond et forme dans chaque paragraphe, risquant la répétition.
\end{definition}

\begin{methode}[Rédiger un paragraphe analytique (règle du « C.E.C. »)]
Pour chaque argument ou procédé analysé :
\begin{enumerate}
    \item \textbf{C}laim (Affirmation) : \esp{El autor denuncia el peligro de las « noticias falsas ».}
    \item \textbf{E}vidence (Citation) : \esp{Escribe: « la mentira se disfraza de verdad ».}
    \item \textbf{C}omment (Analyse) : \esp{La metáfora del « disfraz » sugiere la intención manipuladora y la dificultad de detectar la falsedad, lo que refuerza la tesis de la responsabilidad periodística.}
\end{enumerate}
\emph{Ne citez jamais sans commenter. Ne commentez jamais sans citer.}
\end{methode}

\courssub{Étape 4 — Rédaction de la conclusion (5--7 lignes)}

\begin{methode}[Conclusion en deux temps]
\begin{enumerate}
    \item \textbf{Bilan (síntesis) :} Reprenez la thèse et l'essentiel de la démonstration sans ajouter d'idée nouvelle. \esp{En definitiva, el texto demuestra que la paz no es ausencia de guerra, sino resultado de una información veraz y ética.}
    \item \textbf{Ouverture (apertura) :} Élargissez le débat. Liez au programme (Culture de la paix), à l'actualité (réseaux sociaux au Cameroun, \cle{fake news}, discours de haine), ou à une autre œuvre. \esp{Esta reflexión es vital para Camerún, donde el uso responsable de WhatsApp y Facebook puede prevenir conflictos comunitarios y fortalecer el « vivir juntos ».}
\end{enumerate}
\end{methode}

\courssub{Commentaire modèle sur le texte support}

Le texte support de la leçon (Section 1) s'intitule : \esp{« Los medios de comunicación, ¿armas de guerra o instrumentos de paz? »}. Voici un commentaire modèle respectant la structure, la citation intégrée et l'analyse.

\begin{textebox}[Commentaire modèle — « Los medios de comunicación, ¿armas de guerra o instrumentos de paz? »]
Este artículo de opinión trata del papel de los medios de comunicación en la construcción de la paz. El autor defiende la tesis de que los medios pueden ser instrumentos de paz si informan con verdad y ética. Para convencernos, presenta tres argumentos: el peligro de las noticias falsas que « siembran odio », la responsabilidad del periodista como « guardián de la verdad » y el uso positivo de las redes sociales para el diálogo. En el plano de la expresión, destaca la pregunta del título que interpela al lector, la antítesis « armas de guerra / instrumentos de paz » que estructura el debate, la oposición adversativa « sin embargo » que nuance el pesimismo inicial y el imperativo final « ¡usemos los medios para construir puentes! », que transforma el análisis en un verdadero llamamiento a la acción ciudadana. En conclusión, el texto nos invita a consumir y compartir la información con espíritu crítico. Esta exigencia cobra urgencia en Camerún, donde la viralización de rumores en WhatsApp amenaza la cohesión nacional.
\end{textebox}

\begin{aretenir}[Critères d'évaluation type Bac (sur 20 pts -- grille indicative)]
\begin{itemize}
    \item \textbf{Structure (4 pts) :} Introduction codifiée (4 mouvements), développement en 2 parties équilibrées, conclusion (bilan + ouverture), paragraphes distincts (saut de ligne).
    \item \textbf{Contenu / Analyse (8 pts) :} Thèse bien identifiée, arguments relevés et expliqués, procédés pertinents (lexique, syntaxe, énonciation) \textbf{liés à l'intention persuasive}, pas de paraphrase.
    \item \textbf{Citations (4 pts) :} Citations courtes (3--6 mots), intégrées dans la phrase, entre guillemets espagnols (\esp{« ... »}), commentées immédiatement.
    \item \textbf{Correction linguistique (4 pts) :} Orthographe (accents, \esp{ñ}, \esp{ü}), grammaire (ser/estar, subjonctif, accords), vocabulaire précis, connecteurs variés. \cle{Pénalité :} -0,5 par erreur grave (max -4 pts).
\end{itemize}
\end{aretenir}

\begin{popfigure}
\legende{Organigramme de la structure du commentaire de texte}
\begin{tikzpicture}[
    node distance=1.2cm and 2cm,
    box/.style={draw=popdark, thick, rounded corners, fill=popcream, text width=3.2cm, align=center, font=\small, minimum height=1.8cm},
    arrow/.style={-Stealth, thick, popblue},
    labelbox/.style={font=\tiny\bfseries, text=popblue, anchor=north}
]
% Nodes
\node[box, fill=popblueL, align=center] (intro){\textbf{INTRODUCTION}\\\small Amorce $\rightarrow$ Présentation $\rightarrow$ Problématique $\rightarrow$ Plan};
\node[box, fill=poporangeL, below=of intro, align=center] (part1){\textbf{I. THÈSE \& ARGUMENTS}\\\small (Le \emph{quoi})\\Contenu : Idées forces, exemples, illustrations};
\node[box, fill=popgreenL, right=of part1, align=center] (part2){\textbf{II. PROCÉDÉS D'EXPRESSION}\\\small (Le \emph{comment})\\Forme : Lexique, syntaxe, connecteurs, figures, ton, énonciation};
\node[box, fill=poppurpleL, below=of part1, align=center] (part3){\textbf{III. APPRÉCIATION (Optionnel)}\\\small Portée, limites, actualité, lien personnel/contexte};
\node[box, fill=poppinkL, below=of part3, align=center] (concl){\textbf{CONCLUSION}\\\small Bilan (Synthèse) $\rightarrow$ Ouverture (Élargissement)};

% Arrows
\draw[arrow] (intro.south) -- ++(0,-0.3) -| (part1.north) node[labelbox, pos=0.25, left] {Annonce};
\draw[arrow] (intro.south) -- ++(0,-0.3) -| (part2.north) node[labelbox, pos=0.25, right] {Annonce};
\draw[arrow] (part1.south) -- (part3.north);
\draw[arrow] (part2.south) -- (part3.north -| part2.south);
\draw[arrow] (part3.south) -- (concl.north);

% Alternative direct link if no Part III
\draw[arrow, dashed, popmuted] (part1.south) -- ++(0,-1.5) -| (concl.north) node[labelbox, pos=0.5, left] {Si pas III};
\draw[arrow, dashed, popmuted] (part2.south) -- ++(0,-1.5) -| (concl.north) node[labelbox, pos=0.5, right] {Si pas III};
\end{tikzpicture}
\end{popfigure}

\courssub{Boîte à outils : La citation intégrée (Règle d'or)}

\begin{methode}[Comment citer dans un commentaire ?]
\begin{enumerate}
    \item \textbf{Court :} 3 à 6 mots maximum. Pas de phrase entière.
    \item \textbf{Guillemets espagnols :} \esp{« ... »} (obligatoire).
    \item \textbf{Intégration syntaxique :} La citation doit pouvoir se lire comme une partie de votre phrase française/espagnole.
    \item \textbf{Commentaire immédiat :} Expliquez l'effet de la citation juste après.
\end{enumerate}
\textbf{Modèle :} \esp{El autor afirma que la paz « empieza con una información verdadera »: subraya así la responsabilidad del periodista frente a la desinformación.}
\textbf{Traduction :} L'auteur affirme que la paix « commence par une information vraie » : il souligne ainsi la responsabilité du journaliste face à la désinformation.
\end{methode}

\begin{attention}[Attention : Commentaire $\neq$ Résumé]
\begin{itemize}
    \item \textbf{Résumé (Hors sujet) :} \esp{« El autor dice que las noticias falsas son malas. Luego dice que los periodistas deben decir la verdad. Finalmente pide usar bien las redes. »} $\rightarrow$ Note : 2--4/20.
    \item \textbf{Commentaire (Attendu) :} \esp{« El autor denuncia la toxicidad de las « noticias falsas » mediante una metáfora agrícola (« siembran odio »), lo que dramatiza las consecuencias sociales de la mentira. »} $\rightarrow$ Note : 14--18/20.
\end{itemize}
\cle{Mots-clés :} \textbf{Analyser} = Dire \emph{comment} le texte produit du sens. \textbf{Résumer} = Dire \emph{quoi} le texte dit.
\end{attention}

\courssub{Exemples traduits : Phrases-outils du commentaire}

\begin{center}
\begin{tabularx}{\linewidth}{|X|X|}
\hline
\rowcolor{popblueL} \textbf{Français (votre pensée)} & \textbf{Espagnol (votre rédaction)} \\
\hline
Le texte traite des médias et de la paix. & \esp{El texto trata de los medios de comunicación y la paz.} \\
\hline
L'auteur défend l'idée que les médias construisent la paix. & \esp{El autor defiende la idea de que los medios construyen la paz.} \\
\hline
Il argue que la désinformation provoque des conflits. & \esp{Argumenta que la desinformación provoca conflictos.} \\
\hline
Il illustre cela par l'exemple des réseaux sociaux. & \esp{Ilustra esto con el ejemplo de las redes sociales.} \\
\hline
Il utilise l'antithèse « guerre / paix » pour choquer. & \esp{Utiliza la antítesis « guerra / paz » para impactar.} \\
\hline
L'impératif final appelle à l'action citoyenne. & \esp{El imperativo final llama a la acción ciudadana.} \\
\hline
En conclusion, le texte invite à la vigilance. & \esp{En conclusión, el texto invita a la vigilancia.} \\
\hline
Cela est pertinent au Cameroun aujourd'hui. & \esp{Esto es pertinente en Camerún hoy en día.} \\
\hline
\end{tabularx}
\end{center}

\courssub{Entraînement : Texte d'entraînement et exercice résolu}

Pour vous exercer, voici un texte original, court (env. 220 mots), sur le même thème, suivi d'un exercice résolu de rédaction d'introduction et d'analyse d'un argument.

\begin{textebox}[Texte d'entraînement : « WhatsApp en el aula: ¿rumor o herramienta de paz? »]
En los recreos del liceo bilingüe de Yaoundé, los teléfonos vibran sin cesar. Un mensaje reenviado cien veces: « ¡Mañana huelga violenta en el barrio Mvog-Ada! ». El pánico se apodera de los alumnos; algunos corren a casa, otros llaman a sus padres llorando. Al día siguiente, el barrio amaneció tranquilo. Era un bulo, una « fake news » nacida de un malentendido y amplificada por la velocidad del reenvío.

Sin embargo, ese mismo WhatsApp sirvió ayer para organizar una colecta de libros para la biblioteca de una escuela rural en el Norte. « La tecnología no es buena ni mala », dice la profesora Amina, coordinadora del club de paz. « Depende de quién la use y del corazón que la guíe ». El club acaba de lanzar la campaña « Verifica antes de compartir », con carteles en francés, inglés y fulfulde.

¿Y tú? ¿Eres eslabón de la cadena del miedo o constructor de puentes? La paz digital empieza en tu pulgar. ¡Piénsalo antes de reenviar!
\end{textebox}

\begin{definition}[Glossaire du texte d'entraînement]
\begin{itemize}
    \item \esp{recreo} : récréation, pause.
    \item \esp{reenviar} : transférer, réexpédier.
    \item \esp{bulo} : rumeur, infox, canular (anglicisme \esp{fake news} très usité).
    \item \esp{apoderarse de} : s'emparer de (ici : s'installer).
    \item \esp{amaneció} : il fit jour / la journée se leva (verbe \esp{amanecer}).
    \item \esp{colecta} : collecte, quête.
    \item \esp{eslabón} : maillon (de chaîne).
    \item \esp{fulfulde} : langue peule (langue nationale camerounaise).
\end{itemize}
\end{definition}

\begin{exoresolu}[Exercice résolu : Rédiger l'introduction et analyser l'argument 1]
\textbf{Consigne :} À partir du texte « WhatsApp en el aula... », rédigez l'introduction complète du commentaire (4 mouvements) puis analysez le premier argument (le danger de la rumeur) en appliquant la règle C.E.C.

\tcblower
\textbf{Solution détaillée :}

\textbf{1. Introduction rédigée :}
\esp{En la era de la mensajería instantánea, la información viaja más rápido que la verdad. El artículo « WhatsApp en el aula: ¿rumor o herramienta de paz? », de autor anónimo pero contextualizado en un liceo camerunés, aborda el doble filo de las redes sociales en el entorno escolar. La pregunta directriz es: ¿cómo transformar una herramienta de viralización del miedo en un instrumento de solidaridad y paz? Analizaremos primero el argumento del peligro de la desinformación viral (I), para estudiar luego los procedimientos expresivos que convierten el texto en un llamamiento a la responsabilidad ciudadana (II).}

\textbf{2. Analyse de l'argument 1 (C.E.C.) :}
\begin{itemize}
    \item \textbf{Claim :} El autor denuncia el poder destructivo de los bulos virales en un entorno juvenil vulnerable.
    \item \textbf{Evidence :} Cita la reacción en cadena: « El pánico se apodera de los alumnos; algunos corren a casa, otros llaman a sus padres llorando ».
    \item \textbf{Comment :} El campo léxico del miedo (\esp{pánico, corren, llorando}) y la enumeración de conductas irracionales (\esp{corren... llaman}) dramatizan la consecuencia inmediata de la « fake news »: la parálisis social. La contraposición brutal con la realidad del día siguiente (« El barrio amaneció tranquilo. Era un bulo ») revela el absurdo de la reacción, reforzando la tesis de la necesidad de verificación (\esp{Verifica antes de compartir}).
\end{itemize}
\end{exoresolu}

\begin{exercice}[À vous de jouer : Développement Partie II]
Rédigez la \textbf{Partie II (Procédés d'expression)} du commentaire sur ce même texte. Appuyez-vous sur : la question rhétorique finale, l'opposition \esp{Sin embargo}, la citation de la professeure (autorité), l'impératif final, le lexique de la construction (\esp{constructor, puentes, campaña}). \textbf{Objectif :} 10--12 lignes, 2 citations intégrées minimum.
\end{exercice}

\begin{corrige}[Exercice : Développement Partie II]
\textbf{Piste de correction :}
\esp{En el plano expresivo, el texto combina el análisis racional con la apelación emotiva. La pregunta retórica del título y la final (« ¿Eres eslabón... o constructor...? ») interpelan directamente la conciencia del lector, forzando una autoevaluación moral. La adversativa « Sin embargo » marca el giro argumentativo: del diagnóstico al pronóstico, del miedo a la esperanza. La voz de la profesora Amina (« La tecnología no es buena ni mala... ») aporta una autoridad pedagógica que legitima la tesis del uso responsable. El léxico de la construcción (\esp{constructor, puentes, campaña, eslabón}) reemplaza la metáfora de la cadena viral por una imagen de solidaridad activa. Finalmente, el imperativo « ¡Piénsalo antes de reenviar! » sintetiza la exigencia ética: la paz digital es una decisión individual, inmediata y concreta.}
\end{corrige}

\begin{savaistu}[Le saviez-vous ? Cameroun \& Guinée équatoriale]
La Guinée équatoriale (pays hispanophone d'Afrique centrale, comme le Cameroun est francophone/anglophone) fait face aux mêmes défis : la viralisation de fausses informations sur WhatsApp et Facebook lors des élections ou des crises sanitaires (Ebola, COVID). L'Union Africaine a lancé en 2023 la \esp{« Campaña África Verifica »} pour former des « vérificateurs de faits » (fact-checkers) dans les rédactions. Au Cameroun, des initiatives comme \esp{\#DefyHateNow} ou \esp{Local Youth Corner} forment des jeunes à la \esp{alfabetización mediática} (éducation aux médias) pour lutter contre les \esp{discursos de odio} (discours de haine). Le commentaire de texte que vous apprenez aujourd'hui est donc un entraînement direct à l'exercice de citoyenneté numérique !
\end{savaistu}

\coursec{Production guidée : écrire un appel à la paix}

Dans la section précédente, nous avons appris à \emph{commenter} un texte argumentatif : repérer la thèse, analyser les arguments, évaluer les procédés. Il est temps maintenant de passer de l'analyse à la \cle{production} : c'est toi qui vas écrire. Au baccalauréat, l'expression écrite porte souvent sur des sujets de société — médias, paix, jeunesse — et exige exactement les compétences que nous allons entraîner ici : argumenter, injonctiver, mobiliser le bon lexique. Nous procéderons par étapes, du plus guidé au plus libre, comme on monte un escalier.

\begin{popfigure}
\begin{center}
\begin{tikzpicture}[font=\small]
\draw[fill=popblueL, draw=popblue] (0,0) rectangle (3.2,1);
\node[align=center] at (1.6,0.5) {\textbf{Tâche 1}\\ Compléter};
\draw[fill=popgreenL, draw=popgreen] (1.2,1) rectangle (4.4,2);
\node[align=center] at (2.8,1.5) {\textbf{Tâche 2}\\ Transformer};
\draw[fill=poporangeL, draw=poporange] (2.4,2) rectangle (5.6,3);
\node[align=center] at (4.0,2.5) {\textbf{Tâche 3}\\ Slogans};
\draw[fill=poppurpleL, draw=poppurple] (3.6,3) rectangle (6.8,4);
\node[align=center] at (5.2,3.5) {\textbf{Tâche 4}\\ Tribune};
\draw[-{Stealth}, thick, popdark] (0.3,0.2) -- (6.6,3.8);
\node[align=center, popmuted] at (7.6,2) {du guidé\\ vers le libre};
\end{tikzpicture}
\end{center}
\legende{L'escalier de production : quatre tâches, du texte à trous à la tribune libre.}
\end{popfigure}

\courssub{Tâche 1 : compléter un paragraphe argumentatif à trous}

Avant d'écrire seul, on s'entraîne à \emph{articuler} : un texte argumentatif sans connecteurs est une liste, pas un raisonnement. Rappel des quatre connecteurs de la leçon :

\begin{center}
\begin{tabularx}{\linewidth}{|X|X|X|}\hline
\rowcolor{popblueL} \textbf{Connecteur} & \textbf{Valeur logique} & \textbf{Exemple} \\ \hline
\esp{porque} & cause, justification & \esp{Apago la televisión porque hay demasiada violencia.} \\ \hline
\esp{sin embargo} & opposition, nuance & \esp{Las redes informan; sin embargo, también engañan.} \\ \hline
\esp{además} & ajout d'un argument & \esp{La radio es barata; además, llega a las aldeas.} \\ \hline
\esp{por eso} & conséquence, conclusion & \esp{Hay conflictos; por eso debemos dialogar.} \\ \hline
\end{tabularx}
\end{center}

\begin{exoresolu}[Paragraphe à trous]
Énoncé : complète le paragraphe suivant avec \esp{porque}, \esp{sin embargo}, \esp{además}, \esp{por eso}.

\esp{Los medios de comunicación son muy poderosos \ldots\ldots\ldots{} llegan a millones de personas. A veces difunden noticias falsas; \ldots\ldots\ldots{}, también pueden promover el diálogo entre las comunidades. Informan sobre los conflictos; \ldots\ldots\ldots{}, educan a los jóvenes. \ldots\ldots\ldots{}, debemos usar los medios para construir la paz.}

\tcblower

Solution détaillée :
\begin{enumerate}
\item Premier trou : on \emph{justifie} la puissance des médias → \esp{porque}.
\item Deuxième trou : opposition entre « noticias falsas » et « promover el diálogo » → \esp{sin embargo}.
\item Troisième trou : on ajoute un second aspect positif → \esp{además}.
\item Quatrième trou : phrase conclusive, conséquence logique → \esp{por eso}.
\end{enumerate}
Texte complété : \esp{Los medios de comunicación son muy poderosos porque llegan a millones de personas. A veces difunden noticias falsas; sin embargo, también pueden promover el diálogo entre las comunidades. Informan sobre los conflictos; además, educan a los jóvenes. Por eso, debemos usar los medios para construir la paz.}

Traduction : « Les médias sont très puissants parce qu'ils atteignent des millions de personnes. Parfois ils diffusent de fausses nouvelles ; cependant, ils peuvent aussi promouvoir le dialogue entre les communautés. Ils informent sur les conflits ; de plus, ils éduquent les jeunes. C'est pourquoi nous devons utiliser les médias pour construire la paix. »
\end{exoresolu}

\courssub{Tâche 2 : transformer des infinitifs en injonctions}

Observe d'abord ces transformations :

\begin{exemplebox}[De l'infinitif à l'injonction]
\esp{verificar las fuentes} → \esp{¡Verifica las fuentes!} (à un ami) / \esp{¡Verifiquemos las fuentes!} (à tous) « Vérifie les sources ! / Vérifions les sources ! »

\esp{escuchar a los demás} → \esp{¡Escucha a los demás!} / \esp{¡Escuchemos a los demás!}

\esp{no compartir noticias falsas} → \esp{¡No compartas noticias falsas!} « Ne partage pas de fausses nouvelles ! »
\end{exemplebox}

\begin{propriete}[Former l'injonction : impératif et exhortation]
\begin{itemize}
\item \textbf{À une personne (tú), affirmatif} : 3\textsuperscript{e} personne du singulier du présent : \esp{verifica}, \esp{escucha}, \esp{dialoga}. Exceptions irrégulières : \esp{haz} (hacer), \esp{di} (decir), \esp{ve} (ir), \esp{pon} (poner), \esp{sal} (salir), \esp{ten} (tener), \esp{ven} (venir), \esp{sé} (ser).
\item \textbf{À une personne (tú), négatif} : \esp{no} + subjonctif présent : \esp{no compartas}, \esp{no insultes}, \esp{no odies}.
\item \textbf{À un groupe dont on fait partie (nosotros)} : subjonctif présent à la 1\textsuperscript{re} personne du pluriel : \esp{verifiquemos}, \esp{construyamos}, \esp{dialoguemos}. C'est la forme reine du slogan collectif.
\item \textbf{Souhait} : \esp{ojalá} + subjonctif, ou \esp{que} + subjonctif : \esp{Ojalá dialogue nuestra comunidad.} \esp{Que la paz reine en nuestro país.}
\end{itemize}
\end{propriete}

\begin{attention}[La place des pronoms : le piège majeur]
À l'\textbf{impératif affirmatif}, les pronoms se placent \emph{après} le verbe et s'y \emph{soudent}, avec un accent graphique pour conserver l'accent tonique : \esp{difundámoslo} « diffusons-le », \esp{verifícala} « vérifie-la », \esp{escúchame} « écoute-moi ». À la forme \textbf{négative}, le pronom revient \emph{avant} le verbe : \esp{no lo difundamos} « ne le diffusons pas », \esp{no la compartas}. Ne dis jamais \esp{¡difundamos lo!} ni \esp{¡no difundámoslo!}.
\end{attention}

\begin{exercice}[À ton tour]
Transforme en injonction (tú affirmatif, puis nosotros) : 1. \esp{respetar las opiniones} ; 2. \esp{informarse antes de hablar} ; 3. \esp{no insultar en las redes sociales} ; 4. \esp{sembrar la paz en el barrio}.
\end{exercice}

\begin{corrige}[Exercice : transformations]
1. \esp{¡Respeta las opiniones! / ¡Respetemos las opiniones!} 2. \esp{¡Infórmate antes de hablar! / ¡Informémonos antes de hablar!} (verbe pronominal : pronom soudé, accent). 3. \esp{¡No insultes en las redes sociales! / ¡No insultemos en las redes sociales!} (négatif : subjonctif). 4. \esp{¡Siembra la paz en el barrio! / ¡Sembremos la paz en el barrio!} (attention au changement de voyelle e→ie : \esp{siembra}).
\end{corrige}

\courssub{Tâche 3 : rédiger cinq slogans pour la paix}

Un slogan obéit à une grammaire de la brièveté : une phrase courte, un rythme, une injonction ou un souhait.

\begin{methode}[Écrire un appel à la paix en quatre étapes]
\begin{enumerate}
\item \textbf{Nommer le problème} : \esp{las noticias falsas}, \esp{el odio en las redes}, \esp{la violencia en el barrio}.
\item \textbf{Donner un ordre ou un souhait} : impératif (\esp{¡Verifica!}), exhortation collective (\esp{¡Construyamos!}) ou souhait (\esp{Ojalá reine la paz}).
\item \textbf{Justifier avec \esp{porque}} : \esp{porque una mentira puede destruir una comunidad}.
\item \textbf{Conclure par une formule forte et brève} : \esp{¡La paz empieza contigo!}
\end{enumerate}
\end{methode}

\begin{exemplebox}[Exemples traduits]
\esp{¡No compartas una noticia sin verificarla!} « Ne partage pas une nouvelle sans l'avoir vérifiée ! »

\esp{¡Escuchemos a los demás!} « Écoutons les autres ! »

\esp{Que la paz reine en nuestro país.} « Que la paix règne dans notre pays. »

\esp{¡Construyamos la paz juntos!} « Construisons la paix ensemble ! »

\esp{Ojalá dialogue nuestra comunidad.} « Puissions-nous dialoguer dans notre communauté. »
\end{exemplebox}

\begin{savaistu}[« La paz empieza conmigo »]
\esp{La paz empieza conmigo} (« La paix commence avec moi ») est un slogan utilisé dans de nombreuses campagnes scolaires hispanophones, de l'Espagne à l'Amérique latine. Un bon slogan est \textbf{court}, \textbf{rytmé} et contient un \textbf{impératif} ou une \textbf{exhortation}. À Yaoundé comme à Madrid, les murs des écoles se couvrent de ces phrases lors des journées de la paix : à toi d'inventer les tiennes !
\end{savaistu}

\begin{experience}[Atelier slogans]
Par groupes de quatre, rédigez cinq slogans pour la paix : deux à l'impératif \esp{tú}, deux à la forme \esp{nosotros}, un avec \esp{ojalá} ou \esp{que} + subjonctif. Écrivez-les au tableau ; la classe vote pour le slogan le plus rythmé, qui sera affiché dans la salle.
\end{experience}

\courssub{Tâche 4 : production guidée type Bac — la tribune}

Dernière marche de l'escalier : la rédaction complète. Sujet : \textbf{\esp{Los jóvenes, las redes sociales y la paz}} — une tribune de 10 à 12 lignes avec le plan imposé : \textbf{thèse → deux arguments → appel final}.

\begin{methode}[Rédiger la tribune]
\begin{enumerate}
\item \textbf{Introduction / thèse} (2 lignes) : pose le sujet et annonce ta position avec \esp{pienso que}, \esp{en mi opinión}, \esp{estoy convencido de que}.
\item \textbf{Argument 1} (3 lignes) : introduis-le par \esp{en primer lugar} ou \esp{primero}, justifie avec \esp{porque}, illustre avec un exemple (\esp{por ejemplo}).
\item \textbf{Argument 2} (3 lignes) : enchaîne avec \esp{además} ou \esp{en segundo lugar} ; tu peux nuancer avec \esp{sin embargo}.
\item \textbf{Appel final} (2-3 lignes) : conclus avec \esp{por eso} ou \esp{en conclusión}, puis lance au moins \textbf{trois marques injonctives} : impératifs, \esp{¡construyamos!}, \esp{ojalá}, \esp{que} + subjonctif.
\end{enumerate}
\end{methode}

\begin{textebox}[Modèle de tribune : « Los jóvenes, las redes sociales y la paz »]
\esp{En mi opinión, las redes sociales son una oportunidad magnífica para los jóvenes, pero también un peligro si las usamos sin pensar. Estoy convencido de que podemos convertirlas en instrumentos de paz.}

\esp{En primer lugar, las redes nos permiten informarnos y dialogar con jóvenes de todo el mundo, porque derriban las fronteras. Por ejemplo, un estudiante de Yaundé puede hablar cada día con un estudiante de Malabo o de Madrid. Además, podemos denunciar la injusticia y difundir mensajes de tolerancia. Sin embargo, las redes también propagan noticias falsas y mensajes de odio que dividen a nuestras comunidades.}

\esp{Por eso, debemos actuar. ¡Verifica las fuentes antes de compartir! ¡No difundas el odio, difunde el respeto! Construyamos juntos unas redes al servicio del diálogo, porque la paz empieza con cada uno de nosotros. Que la paz reine en nuestros corazones y en nuestras pantallas.}
\end{textebox}

\textbf{Glossaire} : \esp{derribar las fronteras} : abattre les frontières ; \esp{denunciar} : dénoncer ; \esp{propagar} : propager ; \esp{el odio} : la haine ; \esp{las fuentes} : les sources ; \esp{las pantallas} : les écrans.

\textbf{Questions d'observation} : 1. Où se trouve la thèse ? Souligne-la. 2. Quels connecteurs introduisent les deux arguments ? 3. Relève les marques injonctives de l'appel final et identifie leur forme (impératif tú, nosotros, \esp{que} + subjonctif).

\begin{exercice}[Ta tribune]
Rédige à ton tour une tribune de 10 à 12 lignes sur \esp{Los jóvenes, las redes sociales y la paz}, en respectant le plan imposé. Utilise au moins : cinq mots du lexique de la leçon (\esp{prensa, noticia, periodista, diálogo, conflicto, paz, redes sociales}\ldots), quatre connecteurs et trois marques injonctives.
\end{exercice}

\courssub{Grille d'auto-évaluation et correction collective}

Avant de rendre ta copie, évalue-toi honnêtement avec cette grille — c'est aussi celle du correcteur :

\begin{center}
\begin{tabularx}{\linewidth}{|X|c|}\hline
\rowcolor{popblueL} \textbf{Critère} & \textbf{Oui / Non} \\ \hline
Ai-je une \textbf{thèse claire}, annoncée dès l'introduction ? & \\ \hline
Ai-je développé \textbf{deux arguments} distincts et justifiés ? & \\ \hline
Ai-je utilisé au moins \textbf{trois marques injonctives} (impératif, \esp{ojalá}, \esp{que} + subjonctif) ? & \\ \hline
Ai-je employé le \textbf{lexique de la leçon} (médias et paix) ? & \\ \hline
Ai-je placé correctement les \textbf{pronoms} (\esp{difundámoslo} / \esp{no lo difundamos}) ? & \\ \hline
Ai-je vérifié les \textbf{accents} et la ponctuation \esp{¡\ldots!} ? & \\ \hline
\end{tabularx}
\end{center}

\begin{experience}[Correction collective au tableau]
Le professeur recopie au tableau une production d'élève (anonymisée). Démarche collective : 1. un élève lit le texte à voix haute ; 2. la classe repère la thèse et les arguments et les entoure ; 3. on corrige ensemble les erreurs de forme (impératifs, place des pronoms, accents) ; 4. on améliore l'appel final en ajoutant une marque injonctive ; 5. chacun reprend sa propre copie avec la grille. Exemple de correction type : \esp{¡No difundamos lo!} → \esp{¡No lo difundamos!} ; \esp{verificá las fuentes} → \esp{verifica las fuentes} (pas d'accent à l'impératif \esp{tú} des verbes en -ar).
\end{experience}

\begin{aretenir}[L'essentiel de la production guidée]
Pour écrire un appel à la paix : nomme le problème, injoncte (impératif ou \esp{ojalá}), justifie avec \esp{porque}, conclus par une formule brève. Pour une tribune type Bac : thèse + deux arguments + appel final, reliés par \esp{en primer lugar}, \esp{además}, \esp{sin embargo}, \esp{por eso}. À l'impératif affirmatif, les pronoms se soudent après le verbe avec accent (\esp{verifícala}) ; au négatif, ils précèdent (\esp{no la verifiques}). Un slogan est court, rythmé et injonctif : \esp{¡La paz empieza contigo!}
\end{aretenir}

\newpage

\coursec{Activité d'intégration}

\courssub{Objectifs de l'activité}

À l'issue de cette activité d'intégration, l'élève sera capable de :
\begin{itemize}
\item mobiliser dans une tâche réelle le \cle{lexique} des médias (\esp{prensa, periodista, noticia, informativo, redes sociales}) et de la culture de la paix (\esp{paz, conflicto, diálogo, tolerancia}) ;
\item produire un \cle{texte argumentatif} bref et cohérent : thèse, deux arguments, appel final ;
\item employer correctement les \cle{formes injonctives} : impératif affirmatif et négatif, \esp{ojalá} + subjonctif, \esp{hay que} + infinitif ;
\item présenter oralement un projet collectif et justifier un choix avec des connecteurs d'opinion (\esp{me gusta porque..., en mi opinión..., estoy de acuerdo con...}) ;
\item coopérer en groupe dans une perspective de communication authentique et citoyenne.
\end{itemize}

\courssub{Situation}

Le lycée bilingue de Yaoundé organise sa \emph{Semana de la Paz}. À la suite de plusieurs conflits entre élèves provoqués par des rumeurs et de fausses informations partagées sur WhatsApp et Facebook, le club d'espagnol a été chargé par le proviseur de lancer une campagne de sensibilisation intitulée \esp{« Medios responsables, paz duradera »}. Chaque groupe de quatre élèves de Terminale A devra concevoir une affiche, rédiger un texte argumentatif et le présenter oralement devant la classe. La meilleure campagne sera affichée à l'entrée du lycée et lue lors de la cérémonie de clôture de la semaine.

Voici le message officiel publié par le club d'espagnol pour lancer le concours :

\begin{textebox}[Convocatoria del club de español]
¡Atención, alumnos de Terminale!\\
El club de español organiza la campaña « Medios responsables, paz duradera ».\\
Las redes sociales informan, pero también pueden dividir: una noticia falsa compartida en un grupo de WhatsApp puede crear un conflicto en la escuela.\\
Por eso os invitamos a participar: cread un afiche con un eslogan, escribid un texto argumentativo y presentad vuestra campaña ante la clase.\\
¡Verificad antes de compartir! ¡Construyamos juntos la paz en nuestro liceo!\\
Fecha límite: viernes, 17 horas. La clase votará la mejor campaña.
\end{textebox}

\textbf{Glossaire :} \esp{afiche} : affiche (mot masculin, courant en Amérique latine ; en Espagne on dit plutôt \esp{cartel}) ; \esp{eslogan} : slogan ; \esp{compartir} : partager ; \esp{fecha límite} : date limite ; \esp{noticia falsa} : fausse information ; \esp{liceo} : lycée.

\begin{attention}[Accord et genre dans la convocatoria]
Remarquez l'accord : \esp{un afiche} (masculin) et non \esp{una afiche} ; \esp{una noticia falsa} (féminin). De même, on écrit \esp{el eslogan} (masculin). Le participe passé employé comme adjectif s'accorde : \esp{una noticia falsa \textbf{compartida} en un grupo}.
\end{attention}

\courssub{Consignes}

\begin{enumerate}
\item \textbf{Relecture et repérage.} Relisez la convocatoria du club. Soulignez en rouge les deux injonctions finales et dites à quel mode (impératif ou subjonctif) appartiennent les verbes \esp{verificad} et \esp{construyamos}. Relevez ensuite trois mots du lexique des médias et deux mots du lexique de la paix.

\item \textbf{Création de l'affiche.} En groupe, rédigez un \cle{slogan injonctif} d'une phrase maximum, en utilisant l'impératif (\esp{¡Verifica la noticia!}) ou la structure \esp{ojalá} + subjonctif (\esp{¡Ojalá reine la paz en nuestro liceo!}). Décrivez ensuite en espagnol, en trois phrases, le visuel de votre affiche (couleurs, symboles, personnages).

\item \textbf{Rédaction du texte argumentatif.} Écrivez un texte de 8 à 10 lignes répondant à la question : \esp{¿Por qué debemos verificar las informaciones antes de compartirlas?} Respectez le plan : une thèse claire, deux arguments introduits par \esp{en primer lugar} et \esp{además}, puis un appel final à l'impératif ou avec \esp{ojalá}.

\item \textbf{Présentation orale.} Présentez votre campagne en deux minutes : chaque membre du groupe prend la parole (slogan, description de l'affiche, lecture du texte, appel final). Parlez lentement, articulez et regardez l'auditoire.

\item \textbf{Vote argumenté.} La classe élit la meilleure campagne. Chaque électeur justifie son vote à voix haute avec une phrase du type : \esp{Voto por el grupo 2 porque su eslogan promueve el diálogo y su texto tiene argumentos convincentes.}
\end{enumerate}

\courssub{Ressources à mobiliser}

\begin{aretenir}[Boîte à outils de la campagne]
\textbf{Lexique des médias :} \esp{la prensa, la radio, la televisión, el/la periodista, la noticia, el informativo, las redes sociales, verificar, compartir, informar}.\\
\textbf{Lexique de la paix :} \esp{la paz, el conflicto, el diálogo, la tolerancia, el respeto, la convivencia, la verdad}.\\
\textbf{Injonctions :} impératif affirmatif (\esp{¡Verifica! ¡Lee! ¡Dialoga!}), impératif négatif avec subjonctif (\esp{¡No compartas noticias falsas!}), \esp{ojalá} + subjonctif (\esp{¡Ojalá haya paz!}), \esp{hay que} + infinitif (\esp{Hay que verificar las fuentes}).\\
\textbf{Connecteurs argumentatifs :} \esp{en primer lugar, además, por eso, en conclusión, en mi opinión}.\\
\textbf{Justifier un choix :} \esp{me gusta porque..., voto por... porque..., estoy de acuerdo con...}.
\end{aretenir}

\begin{propriete}[Rappel : les trois formes de l'appel]
\begin{itemize}
\item \textbf{Impératif affirmatif} (2\textsuperscript{e} personne) : \esp{¡Verifica la fuente!} (tú), \esp{¡Verificad las noticias!} (vosotros).
\item \textbf{Impératif négatif} = \esp{no} + subjonctif présent : \esp{¡No compartas rumores!}, \esp{¡No compartáis mentiras!}.
\item \textbf{Exhortation collective} = subjonctif à la 1\textsuperscript{re} personne du pluriel : \esp{¡Construyamos la paz!} ; ou \esp{ojalá} + subjonctif : \esp{¡Ojalá tengamos paz!} ; ou encore la tournure impersonnelle \esp{hay que} + infinitif : \esp{Hay que dialogar}.
\end{itemize}
\end{propriete}

\begin{attention}[Pièges à éviter]
Ne dites pas \esp{las médias} : le mot correct est \esp{los medios (de comunicación)}. Après \esp{ojalá}, le subjonctif est obligatoire : \esp{ojalá \textbf{tengamos} paz}, et non un indicatif. À l'impératif négatif, on utilise le subjonctif : \esp{no compartas}, jamais \esp{no compartes}. Attention au faux ami : \esp{una noticia} signifie « une nouvelle / une information », pas « une notice ». Enfin, \esp{informaciones} au pluriel désigne des « informations » précises ; pour « l'information » en général, on dit \esp{la información}.
\end{attention}

\courssub{Proposition de production et corrigé commenté}

\begin{exoresolu}[Campagne modèle du groupe 3]
Produisez l'affiche (slogan + description du visuel) et le texte argumentatif de votre campagne « Medios responsables, paz duradera ».
\tcblower
\textbf{Eslogan del afiche :} \esp{¡No compartas la mentira, comparte la paz!}

\textbf{Descripción del visual :} \esp{En el centro del afiche hay una paloma blanca que vuela sobre un teléfono móvil. A la izquierda, unos jóvenes discuten por una noticia falsa; a la derecha, los mismos jóvenes dialogan y se dan la mano. Los colores son el verde, el rojo y el amarillo, como la bandera de Camerún.}

\textbf{Texto argumentativo :}

\esp{En nuestro liceo, las redes sociales son muy útiles para informarnos, pero también pueden ser peligrosas. En mi opinión, debemos verificar las informaciones antes de compartirlas. En primer lugar, una noticia falsa puede crear conflictos entre los alumnos: el mes pasado, un rumor falso en un grupo de WhatsApp dividió a toda una clase. Además, un periodista responsable siempre verifica sus fuentes; nosotros debemos hacer lo mismo, porque la verdad es la base de la convivencia. Por eso, pensar antes de compartir es un acto de paz. En conclusión, ¡verificad las noticias, dialogad entre vosotros y construyamos juntos una paz duradera en nuestro liceo!}

\textbf{Commentaire des choix de langue :}
\begin{itemize}
\item Le slogan utilise l'\cle{impératif négatif} (\esp{no compartas}, subjonctif) suivi d'un impératif affirmatif (\esp{comparte}) : l'antithèse \esp{mentira/paz} rend le message percutant.
\item La thèse est annoncée par \esp{en mi opinión} ; les deux arguments sont introduits par \esp{en primer lugar} et \esp{además}, avec un exemple contextualisé (le groupe WhatsApp de la classe).
\item L'appel final combine l'impératif pluriel (\esp{verificad, dialogad}) et le subjonctif d'exhortation (\esp{construyamos}), ce qui illustre les deux formes injonctives de la leçon.
\item Le lexique imposé est bien présent : \esp{redes sociales, noticia falsa, periodista, fuentes, conflicto, diálogo, convivencia, paz duradera}.
\item Accords vérifiés : \esp{una paloma blanca}, \esp{una noticia falsa}, \esp{las redes sociales son útiles}, \esp{los mismos jóvenes}.
\end{itemize}
\end{exoresolu}

\begin{savaistu}[La colombe, symbole universel de la paix]
La \esp{paloma blanca} (colombe blanche) est le symbole de la paix dans le monde hispanophone comme ailleurs : le peintre espagnol Pablo Picasso l'a rendue célèbre avec ses dessins de colombes pour les congrès de la paix. En Guinée équatoriale, pays hispanophone voisin du Cameroun, les campagnes scolaires pour la \esp{convivencia} utilisent aussi ce symbole, souvent associé à la poignée de main (\esp{darse la mano}) et au dialogue (\esp{el diálogo}).
\end{savaistu}

\courssub{Bilan de l'activité}

Cette activité a permis de réinvestir l'ensemble des acquis de la leçon dans une tâche de communication authentique et citoyenne. Sur le plan de la \cle{compréhension}, les élèves ont relu un document injonctif réel (la convocatoria) et repéré ses formes verbales. Sur le plan du \cle{lexique}, ils ont mobilisé activement le vocabulaire des médias et de la paix, à l'écrit comme à l'oral. Sur le plan \cle{grammatical}, ils ont manipulé l'impératif affirmatif et négatif, \esp{ojalá} + subjonctif et \esp{hay que} + infinitif. Sur le plan de l'\cle{expression}, ils ont construit un texte argumentatif structuré (thèse, arguments, appel) et défendu oralement un projet collectif. Enfin, le vote argumenté a développé la compétence d'interaction et l'esprit critique : savoir justifier un choix, c'est déjà pratiquer le dialogue, première étape de la culture de la paix. Le slogan gagnant, affiché à l'entrée du lycée, prolongera la leçon dans la vie quotidienne de l'établissement.

\newpage

\coursec{Méthodes et exercices résolus}

\coursec{Leçon E22 : Lectura y comentario : « Los medios de comunicación » ; textes argumentatifs / injonctifs ; cultura de la paz}

\courssub{Texte support de la leçon}
\begin{textebox}[Texte original — Los medios de comunicación, ¿armas de guerra o instrumentos de paz?]
Los medios de comunicación ocupan un lugar central en la vida de nuestras sociedades. En Camerún, como en España o en Guinea Ecuatorial, la prensa, la radio, la televisión y, hoy sobre todo, las redes sociales moldean la opinión pública. Sin embargo, su poder es una espada de doble filo: pueden construir puentes de diálogo o levantar muros de odio.

Cuando un periodista respeta la veracidad y el rigor, la noticia se convierte en herramienta de paz. Informar con equilibrio sobre un conflicto étnico en el Norte o sobre una tensión política en Yaundé permite a los ciudadanos comprender las causas profundas y fomenta la empatía. Por el contrario, la manipulación, los titulares sensacionalistas y la difusión de bulos —las llamadas \emph{fake news}— siembran la desconfianza y pueden desencadenar violencia. Un informe de la UNESCO recuerda que \textbf{« la libertad de expresión y el acceso a la información son pilares de la cultura de la paz »}.

Ante este panorama, la responsabilidad es compartida. Los profesionales deben cumplir su código deontológico: contrastar fuentes, evitar el lenguaje estigmatizante y dar voz a las víctimas, no solo a los verdugos. Los gobiernos, por su parte, han de garantizar la libertad de prensa sin censuras, pero también legislar contra el discurso de odio. Finalmente, cada ciudadano tiene el deber de desarrollar su espíritu crítico: verificar antes de compartir, escuchar al que piensa distinto y rechazar la polarización.

En definitiva, los medios no son ni buenos ni malos en sí mismos; son espejos de quienes los hacen y de quienes los consumen. La pregunta no es solo qué dicen los medios, sino qué hacemos nosotros con lo que dicen. \textbf{¡Transformemos la información en encuentro, no en trinchera!}
\par\medskip
\textbf{Traduction / Glossaire :} \emph{moldean} (façonnent, modelent) — \emph{espada de doble filo} (arme à double tranchant) — \emph{veracidad} (véracité) — \emph{rigor} (rigueur) — \emph{titulares sensacionalistas} (titres racoleurs/sensationnalistes) — \emph{bulos} (rumeurs, infox) — \emph{siembran} (semencent) — \emph{desconfianza} (méfiance) — \emph{desencadenar} (déclencher) — \emph{código deontológico} (code de déontologie) — \emph{contrastar fuentes} (recouper les sources) — \emph{lenguaje estigmatizante} (langage stigmatisant) — \emph{dar voz a} (donner la parole à) — \emph{legislar contra} (légiférer contre) — \emph{discurso de odio} (discours de haine) — \emph{espíritu crítico} (esprit critique) — \emph{polarización} (polarisation) — \emph{trinchera} (tranchée) — \emph{Yaundé} (Yaoundé, capitale du Cameroun).
\end{textebox}

\courssub{Lexique thématique : Médias et culture de la paix}
\begin{propriete}[Vocabulaire essentiel — Módulo 5]
\textbf{1. Les médias (Los medios de comunicación)}
\begin{center}
\begin{tabularx}{\linewidth}{|X|X|X|}
\hline
\rowcolor{popblueL} \textbf{Español} & \textbf{Français} & \textbf{Exemple / Contexte} \\
\hline
\emph{la prensa} (escrita) & la presse (écrite) & \esp{Leer la prensa digital cada mañana.} \\
\hline
\emph{la radio} & la radio & \esp{La radio comunitaria en zonas rurales.} \\
\hline
\emph{la televisión} / \emph{la tele} & la télévision & \esp{El informativo de la noche.} \\
\hline
\emph{las redes sociales} & les réseaux sociaux & \esp{Twitter (X), Facebook, WhatsApp, TikTok.} \\
\hline
\emph{el periodista} / \emph{la periodista} & le/la journaliste & \esp{Un periodista de investigación.} \\
\hline
\emph{la noticia} & la nouvelle, l'info & \esp{Dar la noticia en primicia.} \\
\hline
\emph{el informativo} / \emph{el telediario} & le journal télévisé & \esp{Ver el informativo de las 21h.} \\
\hline
\emph{el titular} & le titre (une) & \emph{Titular sensacionalista} (titre racoleur). \\
\hline
\emph{el reportaje} & le reportage & \esp{Un reportaje de guerra.} \\
\hline
\emph{la opinión pública} & l'opinion publique & \esp{Moldear la opinión pública.} \\
\hline
\emph{la libertad de prensa} & la liberté de la presse & \esp{Garantizar la libertad de prensa.} \\
\hline
\emph{la censura} & la censure & \esp{Evitar la censura previa.} \\
\hline
\emph{el bulo} / \emph{la fake news} & l'infox, la rumeur & \esp{Combatir los bulos virales.} \\
\hline
\emph{verificar} / \emph{contrastar} & vérifier / recouper & \esp{Contrastar fuentes fiables.} \\
\hline
\end{tabularx}
\end{center}

\textbf{2. Paix, conflit et dialogue (Paz, conflicto y diálogo)}
\begin{center}
\begin{tabularx}{\linewidth}{|X|X|X|}
\hline
\rowcolor{popgreenL} \textbf{Español} & \textbf{Français} & \textbf{Exemple / Contexte} \\
\hline
\emph{la paz} & la paix & \esp{Cultura de paz, proceso de paz.} \\
\hline
\emph{el conflicto} & le conflit & \esp{Conflicto étnico, resolución de conflictos.} \\
\hline
\emph{el diálogo} & le dialogue & \esp{Iniciar un diálogo nacional.} \\
\hline
\emph{la violencia} & la violence & \esp{Desencadenar violencia.} \\
\hline
\emph{el odio} / \emph{el discurso de odio} & la haine / discours de haine & \esp{Legislar contra el discurso de odio.} \\
\hline
\emph{la reconciliación} & la réconciliation & \esp{Comisión de la verdad y reconciliación.} \\
\hline
\emph{la empatía} & l'empathie & \esp{Fomentar la empatía.} \\
\hline
\emph{la tolerancia} & la tolérance & \esp{Educación en la tolerancia.} \\
\hline
\emph{el respeto} & le respect & \esp{Respeto a los derechos humanos.} \\
\hline
\emph{la convivencia} & la coexistence, le vivre-ensemble & \esp{Convivencia pacífica.} \\
\hline
\emph{el extremismo} & l'extrémisme & \esp{Prevenir el extremismo violento.} \\
\hline
\emph{la polarización} & la polarisation & \esp{Rechazar la polarización social.} \\
\hline
\end{tabularx}
\end{center}

\textbf{3. Verbes d'action et de responsabilité}
\begin{itemize}
\item \esp{Informar} (informer) — \esp{Manipular} (manipuler) — \esp{Denunciar} (dénoncer) — \esp{Contrastar} (recouper/vérifier) — \esp{Siembrar} (semer : \emph{siembra dudas/miedo/paz}) — \esp{Desencadenar} (déclencher) — \esp{Fomentar} (favoriser) — \esp{Garantizar} (garantir) — \esp{Legislar} (légiférer) — \esp{Transformar} (transformer) — \esp{Construir} / \esp{Levantar} (construire / élever : \emph{puentes/muros}).
\end{itemize}
\end{propriete}

\courssub{Méthode 1 : Commenter un texte argumentatif (type Bac)}
\begin{methode}[Commentaire composé en 4 étapes]
\textbf{Objectif :} Produire une analyse structurée (introduction, développement en 2-3 parties, conclusion) qui montre comment l'auteur construit son argumentation sur les médias et la paix.

\textbf{Étape 1 — Lecture active et repérage (15 min)}
\begin{enumerate}
\item Identifier le \cle{genre} (article d'opinion, tribune, éditorial), la \cle{source} (presse, blog, discours), le \cle{thème} (médias, paix, responsabilité) et la \cle{thèse} (position de l'auteur : les médias sont neutres, l'usage qu'on en fait détermine leur impact).
\item Surligner les \cle{connecteurs logiques} (\esp{por tanto, sin embargo, por el contrario, en definitiva}) qui structurent le raisonnement.
\item Repérer les \cle{arguments} (exemples concrets : Camerún, UNESCO, code déontologique) et les \cle{procédés rhétoriques} (métaphore \esp{espada de doble filo}, antithèse \esp{armas de guerra / instrumentos de paz}, impératif final).
\end{enumerate}

\textbf{Étape 2 — Construction du plan détaillé (10 min)}
\begin{itemize}
\item \textbf{Introduction :} Accroche (actualité médias/paix) $\rightarrow$ Présentation du texte (titre, auteur, source, date, thème) $\rightarrow$ Thèse résumée $\rightarrow$ Annonce du plan (ex: I. Le pouvoir ambivalent des médias ; II. Les conditions d'une information au service de la paix ; III. L'appel à la responsabilité collective).
\item \textbf{Développement :} 2 ou 3 parties équilibrées. Chaque paragraphe = 1 idée force + citation courte + analyse (procédé + effet de sens). Utiliser \esp{En primer lugar, Además, No obstante, Por último}.
\item \textbf{Conclusion :} Bilan synthétique $\rightarrow$ Ouverture (autre texte, expérience personnelle, rôle des réseaux sociaux au Cameroun aujourd'hui).
\end{itemize}

\textbf{Étape 3 — Rédaction soignée (25 min)}
\begin{itemize}
\item Écrire en espagnol correct (accords, subjonctif après \esp{es importante que}, connecteurs).
\item Citer le texte avec guillemets espagnols \esp{« ... »} et référence de ligne.
\item Varier les verbes d'analyse : \esp{el autor sostiene, denuncia, ilustra, contrapone, insta a, concluye}.
\end{itemize}

\textbf{Étape 4 — Relecture croisée (10 min)}
\begin{itemize}
\item Vérifier : thèse bien restituée ? arguments cités ? connecteurs logiques ? orthographe (accents, \esp{ñ}, \esp{ü}) ? grammaire (ser/estar, indicatif/subjonctif, accords) ? présentation (paragraphes sautés) ?
\end{itemize}
\end{methode}

\begin{exoresolu}[Commentaire modèle : « Los medios de comunicación, ¿armas de guerra o instrumentos de paz? »]
\textbf{Consigne :} Rédigez un commentaire composé d'environ 200 mots en espagnol, structuré en introduction, développement (2 parties) et conclusion, sur le texte ci-dessus.
\tcblower
\textbf{Solution détaillée}

\textbf{1. Analyse préalable (brouillon) :}
\begin{itemize}
\item \textbf{Thèse :} Les médias ont un pouvoir ambivalent (guerre/paix) ; la responsabilité incombe aux journalistes, États et citoyens pour en faire des instruments de paix.
\item \textbf{Structure du texte :} 
\begin{itemize}
\item §1 : Constat du pouvoir central des médias (métaphore épée double tranchant).
\item §2 : Argument 1 — Information rigoureuse = paix (exemple Camerún/Yaundé). Argument 2 — Manipulation/fake news = violence.
\item §3 : Argument d'autorité (UNESCO). Responsabilité partagée : journalistes (déontologie), États (liberté/législation), citoyens (esprit critique).
\item §4 : Conclusion synthétique + injonction finale (impératif \esp{¡Transformemos...!}).
\end{itemize}
\end{itemize}

\textbf{2. Plan retenu :}
\begin{itemize}
\item \textbf{Introduction} : Présentation + thèse + plan (I. L'ambivalence du pouvoir médiatique ; II. La responsabilité partagée pour une culture de paix).
\item \textbf{I. L'ambivalence du pouvoir médiatique : entre construction et destruction}
\begin{itemize}
\item Métaphore \esp{« espada de doble filo »} (l. 3) : pouvoir neutre, usage déterminant.
\item Antithèse \esp{« armas de guerra o instrumentos de paz »} (titre) : double potentialité.
\item Exemple concret (Camerún, Yaundé) : information équilibrée $\rightarrow$ empathie / fake news $\rightarrow$ violence.
\end{itemize}
\item \textbf{II. La responsabilité partagée : journalistes, États, citoyens}
\begin{itemize}
\item Journalistes : déontologie (\esp{contrastar fuentes, evitar lenguaje estigmatizante}).
\item États : garantir liberté + légiférer contre discours de haine.
\item Citoyens : esprit critique (\esp{verificar antes de compartir}).
\item Appel final à l'action : impératif \esp{¡Transformemos...!} $\rightarrow$ engagement personnel.
\end{itemize}
\item \textbf{Conclusion} : Bilan + ouverture sur le défi numérique actuel (réseaux sociaux, IA).
\end{itemize}

\textbf{3. Rédaction du commentaire (espagnol) :}

\begin{textebox}[Commentaire rédigé]
\textbf{Introducción}
En un mundo saturado de información, el texto « Los medios de comunicación, ¿armas de guerra o instrumentos de paz? » aborda con lucidez el papel ambivalente de los medios en la construcción de la paz. El autor defiende la tesis de que los medios no son intrínsecamente buenos ni malos, sino que su impacto depende de la responsabilidad de quienes los producen y consumen. Analizaremos primero la ambivalencia de su poder y, en segundo lugar, la triple responsabilidad que implica una cultura de paz mediática.

\textbf{Desarrollo}
\textbf{I. La ambivalencia del poder mediático: construir puentes o levantar muros}
Desde el título, la antítesis « armas de guerra o instrumentos de paz » plantea la dicotomía central. La metáfora de la « espada de doble filo » (l. 3) ilustra esta neutralidad técnica: el mismo medio puede fomentar el diálogo o la violencia. El autor lo demuestra con un ejemplo contextualizado: en Camerún, una información equilibrada sobre conflictos étnicos o tensiones en Yaundé genera empatía, mientras que los « titulares sensacionalistas » y los « bulos » siembran desconfianza y pueden « desencadenar violencia » (l. 10-11). El conectivo « por el contrario » marca la oposición entre rigor periodístico y manipulación.

\textbf{II. Una responsabilidad compartida: periodistas, Estados y ciudadanía}
La solución no es técnica, sino ética y política. El autor invoca a la UNESCO para recordar que « la libertad de expresión y el acceso a la información son pilares de la cultura de la paz » (l. 13-14). Luego, enumera tres niveles de acción mediante estructuras impersonales de obligación: los periodistas « deben cumplir su código deontológico » (contrastar fuentes, dar voz a las víctimas); los gobiernos « han de garantizar la libertad de prensa » y « legislar contra el discurso de odio »; finalmente, cada ciudadano tiene el « deber de desarrollar su espíritu crítico ». El texto culmina con una injunción directa en « nosotros »: « ¡Transformemos la información en encuentro, no en trinchera! », transformando el lector en actor de la paz.

\textbf{Conclusión}
En definitiva, el texto demuestra que la paz mediática es una construcción colectiva. Su fuerza reside en pasar del diagnóstico a la interpelación concreta de cada agente. Como apertura, cabe preguntarse cómo la inteligencia artificial y los algoritmos de las redes sociales complican hoy esta responsabilidad compartida, especialmente en contextos electorales sensibles en África y América Latina.
\end{textebox}

\textbf{4. Justifications grammaticales clés :}
\begin{itemize}
\item \esp{aborda con lucidez} (indicatif présent : fait réel).
\item \esp{defiende la tesis de que los medios no \textbf{son}...} (Indicatif après \esp{defiende que} affirmatif ; le subjonctif n'apparaît qu'avec la négation du verbe d'opinion : \esp{no defiende que / niega que / duda que} + subjonctif).
\item \esp{genera empatía / siembran desconfianza} (indicatif : conséquences factuelles).
\item \esp{deben cumplir / han de garantizar / tiene el devoir de} : valeurs modales (obligation morale/légale).
\item \esp{¡Transformemos!} : \textbf{impératif 1ère personne pluriel (nosotros)} du subjonctif présent — forme exhortative inclusive (\esp{nosotros} = auteur + lecteurs).
\item \textbf{Accents vérifiés :} \esp{análisis, ética, también, acción, trinchera (sans accent), país, Yaundé}.
\end{itemize}
\end{exoresolu}

\courssub{Textes argumentatifs et injonctifs : Repérage et grammaire de l'appel}
\begin{propriete}[L'argumentation et l'injonction : formes et valeurs]
\textbf{1. Repérer la thèse et les arguments (Rappel Bac)}
\begin{itemize}
\item \textbf{Thèse} = Position de l'auteur (souvent dans le titre, l'introduction ou la conclusion). Question : \esp{¿Qué defiende el autor?}
\item \textbf{Arguments} = Preuves avancées. Types : \cle{Logique} (raisonnement), \cle{D'autorité} (expert, ONU, UNESCO), \cle{Exemplificatif} (cas concret : Camerún), \cle{D'analogie/métaphore} (\esp{espada de doble filo}).
\item \textbf{Connecteurs} = Squelette logique (voir tableau Méthode 2).
\end{itemize}

\textbf{2. Le texte injonctif (appel, exhortation, ordre, conseil)}
Le texte injonctif vise à \textbf{modifier le comportement} du destinataire. Il utilise des \cle{formes de l'obligation/nécessité} :

\begin{center}
\begin{tabularx}{\linewidth}{|X|X|X|}
\hline
\rowcolor{poporangeL} \textbf{Structure} & \textbf{Valeur / Niveau} & \textbf{Exemples du texte / Nouveaux} \\
\hline
\esp{Imperativo affirmative} (\esp{¡Transformemos!}) & Ordre direct, exhortation collective (nous), engagement fort. & \esp{¡Construyamos puentes! / ¡No compartas bulos!} \\
\hline
\esp{Deber + infinitif} (\esp{deben cumplir, tiene el deber de}) & Obligation morale, déontologique, interne. & \esp{Los periodistas deben contrastar fuentes.} \\
\hline
\esp{Haber de + infinitif} (\esp{han de garantizar}) & Obligation légale, institutionnelle, plus formelle. & \esp{Los gobiernos han de legislar.} \\
\hline
\esp{Tener que + infinitif} & Nécessité concrète, contrainte réelle. & \esp{Tenemos que verificar antes de compartir.} \\
\hline
\esp{Hay que + infinitif} & Impersonnel, obligation générale, conseil universel. & \esp{Hay que fomentar el diálogo.} \\
\hline
\esp{Es necesario / imprescindible / urgente + QUE + SUBJONCTIF} & Nécessité évaluée par le locuteur $\rightarrow$ \textbf{Subjonctif obligatoire}. & \esp{Es imprescindible que la prensa sea libre.} \\
\hline
\esp{Infinitif comme impératif} (affiches, pancartes, règlements) & Ordre neutre, dépourvu de personne, très courant dans l'espace public. & \esp{Prohibido pasar. / No tirar basura. / Verificar fuentes.} \\
\hline
\esp{Que + Subjonctif} (exhortatif) & Souhait, vœu, ordre indirect (quejussif). & \esp{¡Que cese la violencia! / ¡Que viva la paz!} \\
\hline
\end{tabularx}
\end{center}

\end{propriete}

\begin{attention}[Pièges fréquents — Francophones]
\begin{itemize}
\item \textbf{Imperatif \emph{nosotros} :} \esp{Transformemos} (subj. prés. 1pl), \textbf{PAS} \esp{*transformamos} (indicatif). Règle : Imperatif \emph{nosotros/nosotras} = Forme du Subjonctif Présent 1pl. Idem \emph{vosotros} : \esp{Transformad} (imperatif propre) vs \esp{Transforméis} (subj).
\item \textbf{Subjonctif après \emph{Es necesario que} :} \esp{Es necesario que \textbf{verifiquemos}} (subj), \textbf{PAS} \esp{*verificamos}.
\item \textbf{\emph{Deber} vs \emph{Haber de} :} \esp{Deber} = devoir moral/conseil (\esp{Debes estudiar}). \esp{Haber de} = obligation future, souvent légale/formelle (\esp{Has de presentar el pasaporte}).
\item \textbf{Faux ami :} \esp{Informativo} = Journal télévisé (nom masc.), \textbf{PAS} "informative" (adj). \esp{La noticia} = Une info / une nouvelle (singulier), \esp{Las noticias} = Les infos / le journal d'infos.
\item \textbf{Accent sur \emph{Yaundé} :} S'écrit avec accent aigu sur le \emph{é} en espagnol (règle des oxytones en voyelle + n/s).
\end{itemize}
\end{attention}

\courssub{Méthode 2 : Repérer la thèse, les arguments et les connecteurs dans un texte argumentatif}
\begin{methode}[Analyse structurelle d'un texte argumentatif]
\textbf{Objectif :} Identifier la thèse (position de l'auteur), les arguments (preuves, exemples) et les connecteurs logiques (articulation) pour répondre aux questions de compréhension « globale » et « détaillée » du Bac.

\textbf{Étape 1 — Trouver la thèse (souvent dans le titre, l'intro ou la conclusion)}
\begin{itemize}
\item Question clé : \esp{¿Qué opina el autor? ¿A favor o en contra de qué?}
\item La thèse = une phrase affirmative qui résume le point de vue. Elle n'est pas un thème (ex: « les médias ») mais une prise de position (ex: « les médias doivent être régulés pour servir la paix »).
\end{itemize}

\textbf{Étape 2 — Isoler les arguments (le « corps » de la démonstration)}
\begin{itemize}
\item Un argument = une idée force + un développement (exemple, chiffre, citation d'autorité, analogie).
\item Repérer les \cle{marqueurs d'argumentation} : \esp{En primer lugar, Además, Por otra parte, Por ejemplo, Según..., Es evidente que...}
\item Classer les arguments : \cle{logiques} (raisonnement), \cle{d'autorité} (expert, institution), \cle{d'analogie}, \cle{exemplificatifs} (cas concret : Camerún, UNESCO).
\end{itemize}

\textbf{Étape 3 — Cartographier les connecteurs logiques (le « squelette »)}

\begin{center}
\begin{tabularx}{\linewidth}{|X|X|X|}
\hline
\rowcolor{popblueL} \textbf{Fonction} & \textbf{Connecteurs espagnols fréquents} & \textbf{Exemple du texte} \\
\hline
Addition / Énumération & \esp{Además, Asimismo, También, En primer lugar, En segundo lugar} & \esp{Además, la manipulación...} \\
\hline
Opposition / Concession & \esp{Sin embargo, No obstante, Por el contrario, Aunque, A pesar de} & \esp{Por el contrario, la manipulación...} \\
\hline
Cause / Explication & \esp{Porque, Puesto que, Ya que, Dado que, Debido a} & \esp{Puesto que informan con equilibrio...} \\
\hline
Conséquence / Conclusion & \esp{Por tanto, Por consiguiente, En consecuencia, En definitiva} & \esp{En definitiva, los medios no son...} \\
\hline
Illustration / Exemple & \esp{Por ejemplo, Como por ejemplo, Un claro ejemplo es} & \esp{Un ejemplo es el conflicto en el Norte...} \\
\hline
Autorité / Renvoi & \esp{Según, Como señala, De acuerdo con} & \esp{Según la UNESCO...} \\
\hline
\end{tabularx}
\end{center}

\textbf{Étape 4 — Répondre aux questions type Bac}
\begin{itemize}
\item \esp{« ¿Cuál es la tesis del autor? »} $\rightarrow$ Réponse en une phrase complète (indicatif).
\item \esp{« Cite dos argumentos del texto. »} $\rightarrow$ Citer l'idée + la preuve (ex: \esp{Argumento de autoridad: cita a la UNESCO; Argumento exemplificativo: el caso de Camerún}).
\item \esp{« ¿Qué conectores marcan la oposición? »} $\rightarrow$ \esp{Por el contrario, Sin embargo}.
\end{itemize}
\end{methode}

\courssub{Production guidée : Écrire un appel à la paix}
\begin{exercice}[Rédaction d'un texte injonctif — Appel à la paix sur les réseaux sociaux]
\textbf{Consigne :} Vous êtes membre d'un club « Jeunesse pour la Paix » à Douala. Rédigez un \textbf{appel à la paix} (80-100 mots) destiné à être publié sur WhatsApp/Facebook, à l'attention des jeunes camerounais, pour les inciter à un usage responsable des réseaux sociaux en période de tensions (élections, conflits locaux).

\textbf{Contraintes obligatoires :}
\begin{enumerate}
\item Utiliser \textbf{au moins 3 structures injonctives différentes} (ex: \esp{¡No compartas...!}, \esp{Hay que verificar...}, \esp{Es urgente que...}, \esp{¡Transformemos...!}, \esp{Debemos...}).
\item Intégrer \textbf{5 mots du lexique} de la leçon (ex: \esp{bulos, discurso de odio, empatía, polarización, verificar, convivencia}).
\item Ton : solennel mais accessible, inclusif (\emph{nosotros}).
\item Orthographe soignée (accents, \esp{ñ}, \esp{¿ ¡}).
\end{enumerate}
\end{exercice}

\begin{corrige}[Exercice — Appel à la paix]
\textbf{Proposition de rédaction (92 mots) :}

\begin{textebox}[Appel : « Jóvenes de Camerún, ¡construyamos paz digital! »]
\textbf{¡Jóvenes de Camerún, no dejemos que las redes sean trincheras!} \esp{En tiempos de tensión, cada clic cuenta.} \esp{\textbf{Hay que verificar} toda información antes de compartirla: los \textbf{bulos} y el \textbf{discurso de odio} solo siembran \textbf{polarización} y destruyen nuestra \textbf{convivencia}.} \esp{\textbf{Es urgente que desarrollemos} un espíritu crítico activo; \textbf{debemos contrastar fuentes} y \textbf{dar voz a la empatía}, no al rencor.} \esp{\textbf{¡Transformemos nuestros muros en puentes!} \textbf{¡Que la paz empiece en tu pantalla!} \esp{Compartir es poder: úsalo para unir.}}
\end{textebox}

\textbf{Analyse des structures utilisées :}
\begin{itemize}
\item \esp{¡No dejemos...!} (Impératif \emph{nosotros} négatif = Subj. prés. 1pl).
\item \esp{Hay que verificar} (Impersonnel + infinitif).
\item \esp{Es urgente que desarrollemos} (Impersonnel + \textbf{que} + Subj.-prés. 1pl).
\item \esp{Debemos contrastar / dar} (\emph{Deber} + infinitif = obligation morale).
\item \esp{¡Transformemos...! / ¡Que la paz empiece...!} (Impératif \emph{nosotros} / Que-jussif 3ème pers).
\item Lexique intégré : \esp{bulos, discurso de odio, polarización, convivencia, empatía, verificar, contrastar fuentes}.
\end{itemize}
\end{corrige}

\courssub{Le saviez-vous ? — Cameroun, Guinée équatoriale et médias}
\begin{savaistu}[Écosystème médiatique et paix en zone francophone/hispanophone africaine]
\begin{itemize}
\item \textbf{Guinée équatoriale :} Seul pays africain hispanophone officiel. Médias : \emph{TVGE} (télévision d'État), \emph{Radio Nacional}, presse écrite (\emph{La Gaceta}). Défi : transition vers le pluralisme et régulation des réseaux sociaux (WhatsApp très dominant).
\item \textbf{Cameroun :} Pays bilingue (FR/EN) mais forte présence hispanophone (frontière avec G.E., diaspora, apprentissage LV2). Médias : \emph{CRTV} (bilingue), radios communautaires (essentielles en zones rurales anglophones et francophones), presse en ligne (\emph{Actu Cameroun}, \emph{Journal du Cameroun}). Rôle clé des \emph{bendskins} et artistes (Stanley Enow, Locko) comme relais d'opinion sur les réseaux.
\item \textbf{UNESCO \& Culture de paix :} Programme \emph{« Youth and Peacebuilding »} au Cameroun. Formation de jeunes blogueurs à la vérification de l'info (\emph{fact-checking}) via \emph{Africa Check} et \emph{DataCheck}.
\item \textbf{Éthique journalistique :} La \emph{Déclaration de principes sur la liberté d'expression en Afrique} (2002, Commission africaine des droits de l'homme) inspire les codes déontologiques locaux : exactitude, équilibre, droit de réponse, protection des sources.
\end{itemize}
\end{savaistu}

\newpage

\coursec{Exercices}

\courssub{Je vérifie mes connaissances}

\begin{exercice}[QCM : les médias et la paix]
Pour chaque question, entoure la bonne réponse.
\begin{enumerate}
\item Un \esp{informativo} est :
\begin{enumerate}
\item un journal télévisé ou radiophonique
\item un journaliste d'investigation
\item une publicité à la radio
\end{enumerate}
\item La thèse d'un texte argumentatif est :
\begin{enumerate}
\item un exemple qui illustre un argument
\item l'idée principale que l'auteur défend
\item la conclusion du texte
\end{enumerate}
\item \esp{Ojalá} + subjonctif exprime :
\begin{enumerate}
\item un ordre
\item un souhait
\item une certitude
\end{enumerate}
\item Le contraire de \esp{la paz} est :
\begin{enumerate}
\item el diálogo
\item la noticia
\item el conflicto
\end{enumerate}
\end{enumerate}
\end{exercice}

\begin{exercice}[Vrai ou faux ? Justifie]
Réponds par \esp{verdadero} ou \esp{falso} et justifie ta réponse en français.
\begin{enumerate}
\item \esp{La prensa} désigne uniquement les journaux gratuits distribués dans la rue.
\item Un texte injonctif cherche à faire agir le lecteur (ordres, conseils, appels).
\item \esp{No compartamos las noticias falsas} est une exhortation à la première personne du pluriel.
\item \esp{Periodista} est un faux ami : il signifie « périodique ».
\end{enumerate}
\end{exercice}

\begin{exercice}[Vocabulaire : trouve le mot]
Complète chaque phrase avec le mot du lexique qui convient : \esp{prensa, periodista, noticia, informativo, diálogo, conflicto, redes sociales, paz}.
\begin{enumerate}
\item Mi hermana es \_\_\_\_\_\_\_\_ : escribe artículos para un periódico de Douala.
\item Anoche vi el \_\_\_\_\_\_\_\_ de las nueve en la televisión.
\item Las \_\_\_\_\_\_\_\_ falsas circulan muy rápido por WhatsApp.
\item El \_\_\_\_\_\_\_\_ entre las comunidades es el camino hacia la \_\_\_\_\_\_\_\_.
\item Muchos jóvenes se informan a través de las \_\_\_\_\_\_\_\_.
\end{enumerate}
\end{exercice}

\begin{exercice}[Conjugaison : l'impératif et l'exhortation]
Mets les verbes entre parenthèses à la forme qui convient (impératif ou subjonctif d'exhortation).
\begin{enumerate}
\item ¡\_\_\_\_\_\_\_\_ (escuchar, tú) la radio para informarte bien !
\item ¡\_\_\_\_\_\_\_\_ (verificar, ustedes) las fuentes antes de publicar !
\item ¡No \_\_\_\_\_\_\_\_ (difundir, nosotros) el odio en las redes !
\item ¡\_\_\_\_\_\_\_\_ (dialogar, vosotros) en lugar de pelear !
\item ¡\_\_\_\_\_\_\_\_ (construir, nosotros) la paz juntos !
\end{enumerate}
\end{exercice}

\courssub{Je m'entraîne}

\begin{exercice}[Version : traduis en espagnol]
Traduis les phrases suivantes en espagnol correct (attention aux subjonctifs et aux exhortations).
\begin{enumerate}
\item Les journalistes doivent vérifier leurs sources.
\item Ne partageons pas les fausses nouvelles.
\item Pourvu que le dialogue règne entre les communautés !
\item Il faut que les médias parlent de paix et non de guerre.
\item Écoutez les informations avant de juger.
\end{enumerate}
\end{exercice}

\begin{exercice}[Texte à trous : un appel à la paix]
Complète ce message publié par un lycéen de Yaoundé avec les verbes suivants, conjugués à l'impératif ou au subjonctif d'exhortation : \esp{respetar, creer, unirse, callar, sembrar}.
\begin{textebox}[Un mensaje para la juventud]
¡Jóvenes de Camerún ! No \_\_\_\_\_\_\_\_ todo lo que leen en las redes sociales : \_\_\_\_\_\_\_\_ la verdad antes de compartir. No \_\_\_\_\_\_\_\_ ante la injusticia : \_\_\_\_\_\_\_\_ su voz a la voz de los que defienden la paz. ¡\_\_\_\_\_\_\_\_ el diálogo en sus barrios y en sus escuelas !
\end{textebox}
\end{exercice}

\begin{exercice}[Appariements : médias, paix, conflit]
Associe chaque mot espagnol (1-8) à sa traduction française (a-h).
\begin{center}
\begin{tabularx}{\linewidth}{|X|X|}
\hline
\rowcolor{popblueL} Español & Français \\
\hline
1. la prensa & a. le dialogue \\
\hline
2. el periodista & b. la guerre \\
\hline
3. el informativo & c. la presse \\
\hline
4. el diálogo & d. le journaliste \\
\hline
5. la guerra & e. le journal télévisé \\
\hline
6. la noticia falsa & f. la rumeur \\
\hline
7. el rumor & g. la fausse nouvelle \\
\hline
8. el locutor & h. le présentateur (radio/TV) \\
\hline
\end{tabularx}
\end{center}
Puis classe ces huit mots dans trois colonnes : \esp{medios / paz / conflicto}.
\end{exercice}

\begin{exercice}[Transformations : du discours direct à l'injonction]
Transforme chaque phrase affirmative en injonction à l'impératif (forme indiquée entre parenthèses), puis traduis le résultat en français.
\begin{enumerate}
\item Ustedes verifican las fuentes. (ustedes)
\item Tú escuchas las dos versiones del conflicto. (tú)
\item Nosotros no difundimos rumores. (nosotros, négatif)
\item Vosotros defendéis la libertad de prensa. (vosotros)
\item Usted no insulta a los demás en las redes. (usted, négatif)
\end{enumerate}
\end{exercice}

\courssub{Je me prépare au Bac}

\begin{exercice}[Compréhension écrite type Bac]
Lis le texte suivant, puis réponds aux questions en espagnol, avec des phrases complètes.
\begin{textebox}[Las redes sociales y la paz]
En los últimos años, las redes sociales se han convertido en el medio de comunicación favorito de los jóvenes africanos. En Camerún, como en muchos países, millones de personas se informan cada día a través de Facebook, WhatsApp o TikTok. Sin embargo, estas herramientas pueden ser un arma de doble filo.

Por un lado, las redes permiten denunciar las injusticias, organizar campañas de solidaridad y difundir mensajes de paz. Durante las crisis, muchos jóvenes han utilizado sus teléfonos para llamar al diálogo y calmar los ánimos. Por otro lado, las noticias falsas circulan a una velocidad impresionante y pueden provocar conflictos entre comunidades. Un simple rumor compartido mil veces se convierte, para muchos, en una verdad.

Por eso, los periodistas y los educadores repiten el mismo mensaje : antes de compartir, hay que verificar. La paz también se construye con los dedos, sobre el teclado del teléfono. Cada usuario es responsable de lo que publica. Si queremos un país unido, debemos usar las redes para unir y no para dividir.
\end{textebox}
\textbf{Questions :}
\begin{enumerate}
\item Di si es verdadero o falso y justifica con una cita del texto : « Los jóvenes africanos prefieren la televisión a las redes sociales. »
\item ¿Cuál es la tesis del autor ? Exprésala con tus propias palabras.
\item ¿Qué significa la expresión \esp{un arma de doble filo} en este contexto ? Explica con tus palabras.
\item Cita dos ventajas y un peligro de las redes sociales según el texto.
\item Reformula en español, sin copiar el texto : « antes de compartir, hay que verificar ».
\end{enumerate}
\end{exercice}

\begin{exercice}[Thème et version type Bac]
\textbf{A. Thème.} Traduis en espagnol (5 phrases) :
\begin{enumerate}
\item Les médias doivent informer sans mentir.
\item Ne croyons pas toutes les rumeurs des réseaux sociaux.
\item Pourvu que la paix revienne dans notre région !
\item Il est important que les jeunes écoutent les informations chaque jour.
\item Parlez aux autres communautés : le dialogue est la meilleure arme.
\end{enumerate}
\textbf{B. Version.} Traduis en français le passage suivant :
\begin{textebox}[Para traducir]
La libertad de prensa es un tesoro frágil. Cuando los periodistas pueden trabajar sin miedo, la sociedad entera respira mejor : los ciudadanos conocen la verdad y los conflictos se resuelven con palabras y no con violencia. Defender a los medios libres es defender la paz.
\end{textebox}
\end{exercice}

\begin{exercice}[Expression écrite type Bac]
Choisis \textbf{un} des deux sujets suivants et rédige un texte en espagnol de \textbf{12 à 15 lignes} (environ 120 à 150 mots). Soigne la présentation, la grammaire et la ponctuation espagnole (¿ ¡).

\textbf{Sujet 1 :} \esp{¿Son las redes sociales un peligro para la paz ? Da tu opinión con al menos dos argumentos y un ejemplo concreto.}

\textbf{Sujet 2 :} \esp{Escribe un mensaje para las redes sociales en el que llames a los jóvenes de tu ciudad a rechazar las noticias falsas y a construir la paz. Utiliza al menos tres imperativos y una exhortación con « ojalá ».}
\end{exercice}

\newpage

\coursec{Corrigés des exercices}

\begin{corrige}[Exercice 1 — QCM : les médias et la paix]
\begin{enumerate}
\item \textbf{Bonne réponse : a.} Un \esp{informativo} est un journal télévisé ou radiophonique. Exemple : \esp{Veo el informativo de las nueve.} « Je regarde le journal de neuf heures. »
\item \textbf{Bonne réponse : b.} La thèse est l'idée principale que l'auteur défend. Les arguments servent à la soutenir ; la conclusion la reformule.
\item \textbf{Bonne réponse : b.} \esp{Ojalá} + subjonctif exprime un souhait : \esp{¡Ojalá venga la paz!} « Pourvu que la paix vienne ! »
\item \textbf{Bonne réponse : c.} Le contraire de \esp{la paz} est \esp{el conflicto} (ou \esp{la guerra}). \esp{El diálogo}, au contraire, est un chemin vers la paix.
\end{enumerate}
\end{corrige}

\begin{corrige}[Exercice 2 — Vrai ou faux ? Justifie]
\begin{enumerate}
\item \textbf{Falso.} \esp{La prensa} désigne l'ensemble des journaux et magazines (la presse écrite), payants ou gratuits, papier ou numérique. Exemple : \esp{la prensa nacional} « la presse nationale ».
\item \textbf{Verdadero.} Le texte injonctif vise à faire agir : il utilise l'impératif (\esp{¡Verifica las fuentes!}), l'infinitif de consigne (\esp{No fumar}) ou le subjonctif d'exhortation (\esp{Construyamos la paz}).
\item \textbf{Verdadero.} \esp{No compartamos} est le subjonctif présent de \esp{compartir}, première personne du pluriel, employé comme exhortation négative : « Ne partageons pas ».
\item \textbf{Falso.} \esp{Periodista} signifie « journaliste » (homme ou femme : \esp{el/la periodista}). Le mot français « périodique » se dit \esp{publicación periódica} ou \esp{revista}. C'est bien un piège de faux ami, mais la phrase inverse la réalité.
\end{enumerate}
\end{corrige}

\begin{corrige}[Exercice 3 — Vocabulaire : trouve le mot]
\begin{enumerate}
\item Mi hermana es \textbf{periodista} : escribe artículos para un periódico de Douala. (« Ma sœur est journaliste. »)
\item Anoche vi el \textbf{informativo} de las nueve en la televisión. (« Hier soir, j'ai regardé le journal de neuf heures. »)
\item Las \textbf{noticias} falsas circulan muy rápido por WhatsApp. (Accord pluriel avec \esp{falsas} : « les fausses nouvelles ».)
\item El \textbf{diálogo} entre las comunidades es el camino hacia la \textbf{paz}. (« Le dialogue entre les communautés est le chemin vers la paix. »)
\item Muchos jóvenes se informan a través de las \textbf{redes sociales}. (« Beaucoup de jeunes s'informent à travers les réseaux sociaux. »)
\end{enumerate}
\textbf{Point de vigilance :} à la phrase 3, deux mots de la liste pouvaient sembler possibles ; seul \esp{noticias} s'accorde avec l'adjectif féminin pluriel \esp{falsas}.
\end{corrige}

\begin{corrige}[Exercice 4 — Conjugaison : l'impératif et l'exhortation]
\begin{enumerate}
\item ¡\textbf{Escucha} la radio para informarte bien ! (Impératif affirmatif \esp{tú} = 3\up{e} personne du singulier du présent : \esp{escucha}.) « Écoute la radio pour bien t'informer ! »
\item ¡\textbf{Verifiquen} las fuentes antes de publicar ! (Impératif \esp{ustedes} = subjonctif présent : \esp{verifiquen}, avec changement orthographique c $\to$ qu.) « Vérifiez les sources avant de publier ! »
\item ¡No \textbf{difundamos} el odio en las redes ! (Exhortation négative \esp{nosotros} : subjonctif présent.) « Ne répandons pas la haine sur les réseaux ! »
\item ¡\textbf{Dialogad} en lugar de pelear ! (Impératif \esp{vosotros} : infinitif sans -r + -d : \esp{dialogar} $\to$ \esp{dialogad}.) « Dialoguez au lieu de vous battre ! »
\item ¡\textbf{Construyamos} la paz juntos ! (Exhortation \esp{nosotros} : subjonctif ; \esp{construir} ajoute un y : \esp{construyamos}.) « Construisons la paix ensemble ! »
\end{enumerate}
\textbf{Point de vigilance :} l'impératif négatif emprunte toujours le subjonctif : \esp{¡No escuches!}, \esp{¡No verifiquen!}. Ne jamais écrire \esp{no escucha}.
\end{corrige}

\begin{corrige}[Exercice 5 — Version : traduis en espagnol]
\begin{enumerate}
\item \esp{Los periodistas deben verificar sus fuentes.} (\esp{deber} + infinitif pour le devoir.)
\item \esp{No compartamos las noticias falsas.} (Exhortation négative à la 1\iere{} personne du pluriel : subjonctif.)
\item \esp{¡Ojalá reine el diálogo entre las comunidades!} (\esp{ojalá} + subjonctif : \esp{reine}, de \esp{reinar}.)
\item \esp{Es necesario que los medios hablen de paz y no de guerra.} (\esp{Es necesario que} + subjonctif : \esp{hablen} ; on accepte aussi \esp{Hay que} + infinitif si la phrase est impersonnelle, mais ici le sujet est précis, donc subjonctif.)
\item \esp{Escuchen las noticias antes de juzgar.} (Impératif \esp{ustedes} : \esp{escuchen}, changement c $\to$ qu non nécessaire ici car \esp{escuchar} $\to$ \esp{escuchen} garde le ch ; attention : \esp{antes de} + infinitif.)
\end{enumerate}
\textbf{Points de vigilance :} « il faut que » se traduit par \esp{es necesario que} ou \esp{hay que} ; avec un sujet exprimé, le subjonctif est obligatoire. « Pourvu que » de souhait = \esp{ojalá} + subjonctif, jamais d'indicatif.
\end{corrige}

\begin{corrige}[Exercice 6 — Texte à trous : un appel à la paix]
Texte complété :
\begin{quote}
\esp{¡Jóvenes de Camerún! No \textbf{crean} todo lo que leen en las redes sociales : \textbf{verifiquen} la verdad antes de compartir. No \textbf{callen} ante la injusticia : \textbf{unan} su voz a la voz de los que defienden la paz. ¡\textbf{Siembren} el diálogo en sus barrios y en sus escuelas!}
\end{quote}
Justifications :
\begin{itemize}
\item \esp{No crean} : impératif négatif \esp{ustedes} $\to$ subjonctif (\esp{creer} : \esp{crean}).
\item \esp{verifiquen} : le verbe attendu est \esp{respetar} ? Non : la liste impose \esp{respetar, creer, unirse, callar, sembrar} ; ici le sens « vérifier » n'existe pas dans la liste, donc la phrase se lit avec \esp{respeten} ? Correction : le deuxième trou prend \esp{respeten} serait incohérent. La bonne répartition est : no \textbf{crean} ; \textbf{respeten} la verdad ; no \textbf{callen} ; \textbf{únanse} ; ¡\textbf{siembran}? Non — voici la version correcte et définitive :
\end{itemize}
Réponses exactes, dans l'ordre des trous :
\begin{enumerate}
\item \textbf{crean} (de \esp{creer}) : « Ne croyez pas tout ce que vous lisez ».
\item \textbf{respeten} (de \esp{respetar}) : « respectez la vérité avant de partager ».
\item \textbf{callen} (de \esp{callar}) : « Ne vous taisez pas devant l'injustice ».
\item \textbf{unan} (de \esp{unir}, base d'\esp{unirse} ; ici construction transitive : \esp{unan su voz a la voz de...}) : « unissez votre voix ». On accepte aussi \esp{únanse} si l'on reformule : \esp{únanse a la voz de los que defienden la paz}.
\item \textbf{siembren} (de \esp{sembrar}, verbe à changement radical e $\to$ ie : \esp{siembren}) : « semez le dialogue ».
\end{enumerate}
\textbf{Traduction du message complet :} « Jeunes du Cameroun ! Ne croyez pas tout ce que vous lisez sur les réseaux sociaux : respectez la vérité avant de partager. Ne vous taisez pas devant l'injustice : unissez votre voix à la voix de ceux qui défendent la paix. Semez le dialogue dans vos quartiers et dans vos écoles ! »
\end{corrige}

\begin{corrige}[Exercice 7 — Appariements : médias, paix, conflit]
\textbf{Appariements :} 1-c ; 2-d ; 3-e ; 4-a ; 5-b ; 6-g ; 7-f ; 8-h.
\begin{center}
\begin{tabularx}{\linewidth}{|X|X|X|}
\hline
\rowcolor{popblueL} Medios & Paz & Conflicto \\
\hline
la prensa & el diálogo & la guerra \\
\hline
el periodista & & la noticia falsa \\
\hline
el informativo & & el rumor \\
\hline
el locutor & & \\
\hline
\end{tabularx}
\end{center}
\textbf{Remarque :} \esp{el diálogo} est le seul mot directement lié à la paix ; \esp{la noticia falsa} et \esp{el rumor} sont classés dans « conflicto » car ils alimentent les tensions. On peut discuter en classe : une \esp{noticia} vraie, elle, sert la paix.
\end{corrige}

\begin{corrige}[Exercice 8 — Transformations : du discours direct à l'injonction]
\begin{enumerate}
\item \esp{¡Verifiquen las fuentes!} « Vérifiez les sources ! » (Impératif \esp{ustedes} = subjonctif : \esp{verifiquen}.)
\item \esp{¡Escucha las dos versiones del conflicto!} « Écoute les deux versions du conflit ! » (Impératif \esp{tú} : \esp{escucha}.)
\item \esp{¡No difundamos rumores!} « Ne répandons pas de rumeurs ! » (Exhortation négative \esp{nosotros} : subjonctif.)
\item \esp{¡Defended la libertad de prensa!} « Défendez la liberté de la presse ! » (Impératif \esp{vosotros} : \esp{defender} $\to$ \esp{defended}, sans changement de radical à l'impératif.)
\item \esp{¡No insulte a los demás en las redes!} « N'insultez pas les autres sur les réseaux ! » (Impératif négatif \esp{usted} = subjonctif : \esp{no insulte}.)
\end{enumerate}
\textbf{Point de vigilance :} à la forme \esp{vosotros}, le changement de radical disparaît : \esp{defended} (et non \esp{defiended}). À la forme négative, toujours le subjonctif : \esp{no insulte}, jamais \esp{no insulta}.
\end{corrige}

\begin{corrige}[Exercice 9 — Compréhension écrite type Bac]
\textbf{Réponses modèles (en espagnol, phrases complètes) :}
\begin{enumerate}
\item \esp{Es falso. El texto dice : « las redes sociales se han convertido en el medio de comunicación favorito de los jóvenes africanos ».} (C'est faux : les jeunes Africains préfèrent les réseaux sociaux, pas la télévision.)
\item \esp{La tesis del autor es que las redes sociales son útiles pero peligrosas : depende de nosotros usarlas para unir y construir la paz, y no para dividir.} (La thèse : les réseaux sont un outil ambivalent ; leur usage responsable construit la paix.)
\item \esp{« Un arma de doble filo » significa que las redes tienen dos caras : pueden hacer el bien (denunciar injusticias, difundir paz) pero también el mal (difundir noticias falsas y provocar conflictos).} (Une arme à double tranchant : avantages et dangers à la fois.)
\item \esp{Dos ventajas : permiten denunciar las injusticias y difundir mensajes de paz. Un peligro : las noticias falsas circulan muy rápido y pueden provocar conflictos entre comunidades.}
\item \esp{Antes de difundir una noticia, debemos comprobar si es verdadera.} (Autre formulation acceptée : \esp{Hay que verificar la información antes de compartirla.})
\end{enumerate}
\textbf{Barème indicatif (10 points) :} 2 points par question (1 point pour la réponse correcte, 1 point pour la justification ou la qualité de l'espagnol).
\textbf{Points de vigilance :} répondre par des phrases complètes avec verbe conjugué ; pour la question 1, la citation doit être exacte et entre guillemets ; à la question 5, « reformuler » interdit de recopier le texte mot à mot.
\end{corrige}

\begin{corrige}[Exercice 10 — Thème et version type Bac]
\textbf{A. Thème — traductions modèles :}
\begin{enumerate}
\item \esp{Los medios deben informar sin mentir.} (\esp{mentir} : infinitif après la préposition \esp{sin}.)
\item \esp{No creamos todos los rumores de las redes sociales.} (Exhortation négative : subjonctif \esp{creamos}, de \esp{creer}.)
\item \esp{¡Ojalá vuelva la paz a nuestra región!} (\esp{ojalá} + subjonctif : \esp{vuelva}, de \esp{volver}, o $\to$ ue.)
\item \esp{Es importante que los jóvenes escuchen las noticias cada día.} (\esp{Es importante que} + subjonctif : \esp{escuchen}.)
\item \esp{Hablen con las otras comunidades : el diálogo es la mejor arma.} (Impératif \esp{ustedes} : \esp{hablen} ; on accepte \esp{Hablad} si le contexte est l'Espagne.)
\end{enumerate}
\textbf{B. Version — traduction modèle :}
\begin{quote}
« La liberté de la presse est un trésor fragile. Quand les journalistes peuvent travailler sans peur, la société entière respire mieux : les citoyens connaissent la vérité et les conflits se résolvent avec des mots et non avec la violence. Défendre les médias libres, c'est défendre la paix. »
\end{quote}
\textbf{Barème indicatif (10 points) :} thème 5 points (1 par phrase), version 5 points (fidélité 3, qualité du français 2).
\textbf{Points de vigilance :} \esp{tesoro} = « trésor » (faux ami avec « trésorerie ») ; \esp{sin miedo} = « sans peur » ; \esp{se resuelven} = « se résolvent » (présent, 3\up{e} personne du pluriel de \esp{resolver}) ; ne pas traduire \esp{defender} par « défendre de » + infinitif.
\end{corrige}

\begin{corrige}[Exercice 11 — Expression écrite type Bac]
\textbf{Proposition de rédaction modèle — Sujet 2 (message pour les réseaux sociaux) :}
\begin{quote}
\esp{¡Jóvenes de Yaoundé ! Nuestra ciudad necesita vuestra ayuda. Cada día, noticias falsas circulan por WhatsApp y Facebook, y siembran el odio entre nosotros. ¡No compartan esos rumores ! ¡Verifiquen siempre las fuentes antes de publicar ! ¡Escuchen a los periodistas serios y respeten la verdad ! Construyamos juntos una juventud unida : hablemos de paz en nuestras escuelas, en nuestros barrios y en nuestras redes. ¡Ojalá seamos todos mensajeros de paz y no de división ! El futuro de Camerún está en nuestras manos : usemos nuestros teléfonos para unir, nunca para dividir. ¡La paz empieza con un clic responsable !}
\end{quote}
\textbf{Traduction :} « Jeunes de Yaoundé ! Notre ville a besoin de votre aide. Chaque jour, de fausses nouvelles circulent sur WhatsApp et Facebook et sèment la haine entre nous. Ne partagez pas ces rumeurs ! Vérifiez toujours les sources avant de publier ! Écoutez les journalistes sérieux et respectez la vérité ! Construisons ensemble une jeunesse unie : parlons de paix dans nos écoles, nos quartiers et nos réseaux. Puissions-nous tous être des messagers de paix et non de division ! L'avenir du Cameroun est entre nos mains : utilisons nos téléphones pour unir, jamais pour diviser. La paix commence par un clic responsable ! »
\textbf{Barème indicatif (10 points) :} respect de la consigne et du format injonctif (2) ; au moins trois impératifs et un \esp{ojalá} + subjonctif corrects (3) ; richesse du lexique médias/paix (2) ; correction grammaticale et orthographique, accents et ¡ ¿ (2) ; présentation et longueur (1).
\textbf{Points de vigilance :} penser aux signes d'ouverture ¡ et ¿ ; vérifier chaque impératif (\esp{verifiquen}, \esp{escuchen}, \esp{respeten}) ; \esp{ojalá} est toujours suivi du subjonctif (\esp{seamos}) ; éviter le calque français « las redes sociales son un peligro \emph{pour} la paz » $\to$ \esp{para la paz}.
\end{corrige}

\newpage

\coursec{Fiche bilan}

\begin{popfigure}
\begin{tikzpicture}[mindmap, grow cyclic, every node/.style={concept, font=\small},
  root concept/.append style={concept color=popblue, font=\small\bfseries, text=white},
  level 1/.append style={level distance=3.1cm, sibling angle=51.5},
  level 1 concept/.append style={font=\footnotesize}]
\node[root concept]{Los medios de comunicación y la paz}
  child[concept color=poporange]{ node {Lexique des médias} }
  child[concept color=popgreen]{ node {Lexique de la paix} }
  child[concept color=poppurple]{ node {Texte argumentatif} }
  child[concept color=poppink]{ node {Texte injonctif} }
  child[concept color=popgold]{ node {Impératif et subjonctif} }
  child[concept color=popblue!60]{ node {Méthode du commentaire} }
  child[concept color=popgreen!70!black]{ node {Appel à la paix} };
\end{tikzpicture}
\legende{Carte mentale de la leçon : médias, argumentation et culture de la paix.}
\end{popfigure}

\begin{bilan}
\textbf{1. Lexique clé des médias et de la paix}

\begin{center}
\begin{tabularx}{\linewidth}{|X|X|X|X|}\hline
\rowcolor{popblueL} Español & Français & Español & Français \\ \hline
la prensa & la presse & la paz & la paix \\ \hline
la televisión & la télévision & el conflicto & le conflit \\ \hline
la radio & la radio & el diálogo & le dialogue \\ \hline
el/la periodista & le/la journaliste & la violencia & la violence \\ \hline
la noticia & la nouvelle, l'info & la tolerancia & la tolérance \\ \hline
el informativo & le journal télévisé & la reconciliación & la réconciliation \\ \hline
las redes sociales & les réseaux sociaux & los derechos humanos & les droits de l'homme \\ \hline
la noticia falsa & la fausse information & la campaña de paz & la campagne de paix \\ \hline
el titular & le titre (de presse) & el mensaje & le message \\ \hline
informar / opinar & informer / donner son avis & convivir & vivre ensemble \\ \hline
\end{tabularx}
\end{center}

\textbf{2. Grammaire à connaître par cœur}

\begin{center}
\begin{tabularx}{\linewidth}{|l|X|X|}\hline
\rowcolor{popblueL} Point & Règle & Exemple \\ \hline
Impératif affirmatif (tú) & 3\up{e} personne du singulier du présent & \esp{¡Dialoga!} « Dialogue ! » \\ \hline
Impératif (usted) & subjonctif présent & \esp{¡Escuche!} « Écoutez ! » \\ \hline
Impératif négatif & no + subjonctif présent & \esp{¡No odies!} « Ne hais pas ! » \\ \hline
Nosotros (exhortation) & subjonctif présent, 1\up{re} pers. plur. & \esp{¡Construyamos la paz!} « Construisons la paix ! » \\ \hline
Verbes à changement & e$\rightarrow$ie, o$\rightarrow$ue, e$\rightarrow$i & \esp{¡Piensa! ¡Duerme! ¡Pide!} \\ \hline
Pronoms & se placent après l'impératif affirmatif, avant le négatif & \esp{¡Díselo! / ¡No se lo digas!} \\ \hline
\end{tabularx}
\end{center}

\textbf{3. Texte argumentatif vs texte injonctif}
\begin{itemize}
\item \cle{Texte argumentatif} : une \cle{thèse} (opinion défendue) + des \cle{arguments} (raisons) + des \cle{exemples} ; connecteurs : \esp{porque, ya que, sin embargo, por lo tanto, en conclusión}.
\item \cle{Texte injonctif} : vise à faire agir ; marques : impératif, subjonctif d'exhortation, \esp{¡}, vocabulaire de l'appel (\esp{unámonos, digamos no a la violencia}).
\end{itemize}

\textbf{4. Méthode express du commentaire de texte}
\begin{enumerate}
\item Lire deux fois ; situer : source, thème, type de texte.
\item Dégager la thèse et numéroter les arguments.
\item Repérer les procédés (impératifs, questions rhétoriques, connecteurs).
\item Rédiger : introduction (présentation + problématique), développement (idées + procédés), conclusion (bilan + avis personnel justifié).
\end{enumerate}

\textbf{5. Pièges à éviter}
\begin{itemize}
\item \esp{la noticia} = la nouvelle, l'information (pas « la notice ») ; \esp{los medios} = les médias (pas « les moyens » au sens d'outils matériels).
\item Accent obligatoire quand on colle les pronoms : \esp{díselo}, \esp{escúchame}.
\item Ne pas oublier les signes ¡ et ¿ en tête de phrase.
\end{itemize}
\end{bilan}

\begin{aretenir}[Checklist avant le Bac]
\begin{itemize}
\item[$\square$] Je connais le lexique des médias : \esp{prensa, periodista, noticia, informativo, redes sociales}.
\item[$\square$] Je connais le lexique de la paix : \esp{paz, conflicto, diálogo, tolerancia, reconciliación}.
\item[$\square$] Je sais distinguer texte argumentatif (convaincre) et texte injonctif (faire agir).
\item[$\square$] Je sais repérer la thèse et les arguments d'un texte.
\item[$\square$] Je maîtrise l'impératif affirmatif de tú : \esp{¡habla! ¡escribe! ¡pide!}.
\item[$\square$] Je maîtrise l'impératif négatif avec le subjonctif : \esp{¡no hables! ¡no escribas!}.
\item[$\square$] Je sais former l'exhortation collectée : \esp{¡construyamos la paz!}.
\item[$\square$] Je place correctement les pronoms : \esp{¡díselo!} mais \esp{¡no se lo digas!}.
\item[$\square$] Je connais les connecteurs argumentatifs : \esp{sin embargo, por lo tanto, en conclusión}.
\item[$\square$] Je sais rédiger l'introduction d'un commentaire (source, thème, problématique).
\item[$\square$] Je n'oublie jamais les signes ¡ et ¿ ni les accents espagnols.
\item[$\square$] Je sais écrire un court appel à la paix avec au moins trois impératifs.
\end{itemize}
\end{aretenir}

\findecours

\end{document}
